\chapter{Local valuative pluripotential theory}
\index{pluripotential theory!local valuative}
\index{valuation space!normalized local}
\index{formal psh function}
\label{chap:local-valuative-psh}

The global theories in earlier parts are organized by a projective variety,
a cohomology or numerical class, and models of the whole space.  A local
singularity has none of this data.  It does, however, have a tower of
modifications over the marked point.  Boucksom--Favre--Jonsson showed that
this tower carries a genuine local non-Archimedean pluripotential theory
\cite{BFJ08}.  Its potentials are nef Weil divisors, its model functions come
from primary ideals, and its Monge--Amp\`ere measures are local intersection
measures on a normalized valuation link.

This chapter has three purposes.  First, it gives an account of that theory
with explicit source inputs, which is the natural predecessor of any proposed local
analogue of the Boucksom--Jonsson theory.  Second, it records the interfaces
proved in the sources cited here or in earlier chapters and keeps every
missing comparison arrow explicit.  Third, it develops the finite-energy
branch on polarized affine cones established by Wu.  That branch is a genuine
local analogue of part of the Boucksom--Jonsson theory, but its Reeb
polarization and its analytic link are additional data; it is not a theorem
for an arbitrary analytic germ.  We do not turn the remaining missing arrows
into definitions.  In particular, the arbitrary-germ finite-energy and
singular-product problems listed at the end remain research questions unless
a later Rethlas receipt proves them.

\section{The normalized valuation link}
\label{sec:local-normalized-link}

Let $(Y,0)$ be a smooth complex germ of dimension $n\geq1$.  After
choosing local coordinates, its algebraic local model is
\[
 R=\mathbf C[z_1,\ldots,z_n]_{(z_1,\ldots,z_n)},
 \qquad
 \mathfrak m=(z_1,\ldots,z_n),
 \qquad
 \widehat R=\mathbf C[[z_1,\ldots,z_n]].
\tag{70.1}\label{eq:local-regular-ring}
\]
The construction depends only on the smooth germ, not on the coordinates.
Working with $R$ makes the birational models algebraic; analytic psh germs
will re-enter in \Cref{sec:local-valuative-transform}.
Here and below, a modification above $0$ is projective, is an isomorphism
away from $0$, and has smooth source.  It is called \emph{simple} when
the scheme-theoretic inverse image of $0$ has simple-normal-crossings
support.

A centered semivaluation is a map
\[
 \nu:R\longrightarrow[0,+\infty]
\tag{70.2}\label{eq:local-semivaluation-map}
\]
that is trivial on $\mathbf C^*$, is multiplicative, satisfies the
ultrametric inequality, and is positive on $\mathfrak m$.  Put
\[
 \nu(\mathfrak a):=\min\{\nu(f):f\in\mathfrak a\},
 \qquad
 \nu(\mathfrak m):=\min_i\nu(z_i).
\tag{70.3}\label{eq:local-ideal-value}
\]
The minimum in the first formula can be taken over any finite generating
set.

\begin{definition}
\label{def:local-normalized-space}
The normalized valuation space of the germ is
\[
 \mathcal V_0
 :=\{\nu:\nu\text{ is centered at }0,\ \nu(\mathfrak m)=1\},
\tag{70.4}\label{eq:local-normalized-space}
\]
with the topology of pointwise convergence on $R$.
\end{definition}

Every centered semivaluation with
$0<\nu(\mathfrak m)<+\infty$ has a unique positive rescaling in
$\mathcal V_0$.  The semivaluation that is $+\infty$ on
$\mathfrak m$, whose home is the closed point, does not meet this slice.
Thus $\mathcal V_0$ is the link obtained by quotienting the finite-radial
part of the centered cone by the scaling action.  It is
also the fiberwise slice of the trivially valued Berkovich cone selected by
the equation $\nu(\mathfrak m)=1$.  This normalization is local boundary
data: without the ideal $\mathfrak m$, a homogeneous local potential has no
preferred scale.

\begin{theorem}\label{thm:local-valuation-compactness}
The space $\mathcal V_0$ is compact Hausdorff.  Divisorial valuations are
dense in it, and the union of its quasimonomial dual complexes is dense.
\end{theorem}

\begin{proof}
When $n=1$, every normalized semivaluation on the smooth curve ring is
$\ord_0$; the assertion follows directly from this one-point description.
For $n\geq2$,
BFJ identify $\mathcal V_0$ with the projective limit of the compact finite
dual complexes $\Delta(\pi)$, with transition maps given by retractions.
This gives compactness and Hausdorffness.  They identify the quasimonomial
locus with the corresponding inductive union, whose image is dense in the
projective limit.  Finally, the divisorial points of each simplex are exactly
its rational-weight points, hence are dense there.  This proves both density
statements.  See \cite[Sections~1.3--1.7, especially
Theorem~1.13]{BFJ08}.
\end{proof}

Let $\pi:Y_\pi\to Y$ be a simple modification above $0$, with reduced
snc exceptional divisor $E=\sum_{i\in I}E_i$.  Set
\[
                         b_i:=\ord_{E_i}(\mathfrak m),
 \qquad
                         \nu_{E_i}:=b_i^{-1}\ord_{E_i}.
\tag{70.5}\label{eq:local-divisorial-normalization}
\]
If a stratum is contained precisely in the components indexed by
$J\subset I$, its simplex is
\[
 \sigma_J
 =\left\{w\in\mathbf R_{\geq0}^J\ \middle|\
                    \sum_{j\in J}b_jw_j=1\right\}.
\tag{70.6}\label{eq:local-weight-simplex}
\]
An adapted coordinate system turns $w$ into a monomial valuation.  The
simplices glue to $\Delta(\pi)\subset\mathcal V_0$, and there is a
retraction
\[
                         r_\pi:\mathcal V_0\longrightarrow\Delta(\pi).
\tag{70.7}\label{eq:local-retraction}
\]
The union $\mathcal V_{0,\mathrm{qm}}:=\bigcup_\pi\Delta(\pi)$ is the
quasimonomial locus.  Besides the weak topology inherited from
$\mathcal V_0$, it carries the strong topology obtained by requiring the
Euclidean topology on every simplex.  The distinction is essential in the
definition of a formal psh function.

\begin{example}[A smooth curve germ]
\label{ex:local-smooth-curve-link}
If $n=1$, then every normalized centered valuation has
$\nu(z)=1$.  Hence $\mathcal V_0$ is the single point
$\ord_0$.  The elaborate local theory below reduces to one-variable order
of vanishing.  This is the correct unit test for every normalization: a
local Monge--Amp\`ere measure must be a multiple of
$\delta_{\ord_0}$.
\end{example}

% Adapted from the author-supplied normalized-valuation-links.pdf,
% 23 September 2026. Source bytes and review scope are recorded separately.
\subsection{Excellent germs and boundary ideals}
\label{sec:local-excellent-boundary-links}

The normalized boundary itself has a construction in every dimension.
In this subsection let $(R,\mathfrak m)$ be a positive-dimensional
excellent reduced Noetherian local ring. A nonnegative real semivaluation
has $v(0)=+\infty$, $v(1)=0$, $v(fg)=v(f)+v(g)$ and
$v(f+g)\geq\min\{v(f),v(g)\}$. We use the centered cone
\[
 \mathcal C_R
   :=\{v:R\longrightarrow[0,+\infty]\mid
          v\text{ is a semivaluation},\ 0<v(\mathfrak m)<+\infty\}.
\]
It has the topology of pointwise convergence, with the compact order
topology on $[0,+\infty]$, and the positive scaling action $v\mapsto tv$.
For an ideal $I=(f_1,\ldots,f_N)$, set $v(I)=\min_jv(f_j)$.
Expressing either generating set in terms of the other proves independence
of this choice. Units have value zero. The support
$\{f:v(f)=+\infty\}$ is prime; it need not be a minimal prime.
For an $\mathfrak m$-primary ideal $\mathfrak a$, write
$\mathcal V_{\mathfrak a}(R)=\{v\in\mathcal C_R:v(\mathfrak a)=1\}$,
consistently with the curve notation below.

\begin{theorem}\label{thm:p121-excellent-primary-links}
Every $\mathcal V_{\mathfrak a}(R)$ is nonempty compact Hausdorff,
and $(t,v)\mapsto tv$ is a homeomorphism
$(0,+\infty)\times\mathcal V_{\mathfrak a}(R)\simeq\mathcal C_R$.
For primary ideals $\mathfrak a,\mathfrak b$, the unique map of
sections preserving every ray is the homeomorphism
\[
 H_{\mathfrak a,\mathfrak b}(v)=\frac{v}{v(\mathfrak b)}.
\tag{70.7a}\label{eq:p121-general-boundary-change}
\]
It has inverse $H_{\mathfrak b,\mathfrak a}$ and satisfies
$H_{\mathfrak b,\mathfrak d}H_{\mathfrak a,\mathfrak b}
=H_{\mathfrak a,\mathfrak d}$ for every third primary ideal.
Thus the quotient $\mathcal C_R/\mathbf R_{>0}$ and these sections
are independent of generators and birational models.
\end{theorem}

\begin{proof}
The semivaluation axioms define a closed subset of
$[0,+\infty]^R$: extended addition and minimum are continuous,
and the order relation is closed. The additional condition
$v(\mathfrak m)=1$ is closed because $\mathfrak m$ is finitely
generated. Hence $\mathcal V_{\mathfrak m}(R)$ is compact Hausdorff.
Choose $r\geq s\geq1$ with
$\mathfrak m^r\subseteq\mathfrak a\subseteq\mathfrak m^s$.
For every centered $v$,
\[
                  s\,v(\mathfrak m)\leq v(\mathfrak a)
                                     \leq r\,v(\mathfrak m).
\tag{70.7b}\label{eq:p121-primary-gauge-bounds}
\]
Thus $v(\mathfrak a)$ is positive and finite. The continuous map
$v\mapsto v/v(\mathfrak a)$ sends
$\mathcal V_{\mathfrak m}(R)$ onto $\mathcal V_{\mathfrak a}(R)$,
proving compactness. The inverse of the asserted cone homeomorphism
is $u\mapsto(u(\mathfrak a),u/u(\mathfrak a))$; positive scaling
is jointly continuous also at infinite coordinate values.

If $\mathfrak m^u\subseteq\mathfrak b\subseteq\mathfrak m^t$,
then $t/r\leq v(\mathfrak b)\leq u/s$ on
$\mathcal V_{\mathfrak a}(R)$. This proves continuity of the change
of section with a denominator bounded away from zero and infinity.
Homogeneity proves its inverse and cocycle formulas. On the ray of $v$,
the only point having $\mathfrak b$-value one is $v/v(\mathfrak b)$.
The cone product decomposition identifies its ray quotient with any section.

For nonemptiness, choose a saturated prime chain ending at $\mathfrak m$
and let $\mathfrak p$ be its penultimate member. Then $R/\mathfrak p$
is an excellent one-dimensional local domain. Excellence implies the
Nagata property, and normalization is finite
\cite[Tags 07QV and 035S]{StacksProject}. Localizing that normalization
at a maximal ideal above $\mathfrak m/\mathfrak p$ gives a
one-dimensional normal Noetherian local domain, hence a DVR
\cite[Tag 00PD]{StacksProject}. Its order pulls back to an element of
$\mathcal C_R$, which can be rescaled into every primary section.
\end{proof}

Let $S=\prod_iS_i$ be the normalization of $R$ in its total ring of
fractions. It is finite, and $S_i$ is the normal domain associated to a
minimal component; see \cite[Tags 035P and 035S]{StacksProject}.
The conductor $\mathfrak c=\operatorname{Ann}_R(S/R)$ is also an
ideal of $S$, with $\mathfrak cS=\mathfrak c$. In particular
\[
                    R=S\mathbin{\times}_{S/\mathfrak cS}R/\mathfrak c.
\tag{70.7c}\label{eq:p121-conductor-pullback-square}
\]
Indeed, an element of $S$ whose residue comes from $R/\mathfrak c$
differs from an element of $R$ by an element of $\mathfrak cS\subset R$.

Let $\mathcal B_{\mathfrak a,i}$ be the nonnegative real semivaluations
on $S_i$ with $w(\mathfrak aS_i)=1$, and set
$\mathcal B_{\mathfrak a}=\coprod_i\mathcal B_{\mathfrak a,i}$.
Each summand is closed in $[0,+\infty]^{S_i}$, so this finite disjoint
union is compact Hausdorff. Restriction defines a continuous map
$\pi_{\mathfrak a}:\mathcal B_{\mathfrak a}\to
\mathcal V_{\mathfrak a}(R)$; the primary-ideal bounds prove the
centering condition downstairs.

This presentation retains all local normalization branches. The center
$\{s:w(s)>0\}$ contains $\mathfrak mS_i$ and is a maximal ideal
$\mathfrak n$ above $\mathfrak m$, since $S_i/\mathfrak mS_i$ is
finite over $R/\mathfrak m$. Every denominator outside $\mathfrak n$
has value zero, so $w$ extends uniquely to $(S_i)_{\mathfrak n}$.
There are finitely many such centers. Their pieces are clopen:
the Chinese remainder theorem supplies an element equal to zero modulo
one of these maximal ideals and to one modulo all the others; its
positive-value locus selects that center. Restriction from the local
branch link is a homeomorphism onto this piece, since localized
evaluations are the corresponding numerator evaluations.

\begin{lemma}\label{lem:p121-normalization-semivaluation-lift}
The restriction map $\pi_{\mathfrak a}$ is surjective.
\end{lemma}

\begin{proof}
Let $v$ have support $\mathfrak p$. Choose a minimal prime
$\mathfrak q_i\subseteq\mathfrak p$ and, by lying over, a prime
$P\subset S_i$ above $\mathfrak p/\mathfrak q_i$.
The field $L=\operatorname{Frac}(S_i/P)$ is finite over
$K=\operatorname{Frac}(R/\mathfrak p)$. The finite-valued part of
$v$ defines a valuation ring $A\subset K$ with value group
$\Gamma_A\subset\mathbf R$.

Let $A'$ be the integral closure of $A$ in $L$, choose a prime
$\mathfrak n\subset A'$ over the maximal ideal of $A$, and take
a valuation ring $W\subset L$ dominating $A'_{\mathfrak n}$
\cite[Tag 00IA]{StacksProject}. Then $W\cap K$ dominates $A$
and equals $A$: if $x\in(W\cap K)\setminus A$, its inverse
lies in the maximal ideal of $A$, hence in that of $W$, contradicting
$x\in W$.

The inclusion of value groups has finite index
\cite[Tag 0ASH]{StacksProject}. For $\gamma\in\Gamma_W$, choose
$N>0$ with $N\gamma\in\Gamma_A$ and assign to $\gamma$ the real
number $(N\gamma)/N$, using the given embedding of $\Gamma_A$
in $\mathbf R$. A common multiple of two choices proves independence;
the same construction respects addition and order. It therefore extends
the original real valuation. Every element of $S_i/P$ is integral
over $R/\mathfrak p\subset A$ and belongs to the integrally closed
ring $W$. Restrict this real valuation to $S_i/P$ and give elements
of $P$ value $+\infty$. The resulting semivaluation restricts to
$v$ and has $w(\mathfrak aS_i)=1$.
\end{proof}

\begin{theorem}\label{thm:p121-excellent-conductor-descent}
The map $\pi_{\mathfrak a}$ presents $\mathcal V_{\mathfrak a}(R)$
as the compact Hausdorff quotient of $\mathcal B_{\mathfrak a}$ by
the closed relation of equal restrictions to $R$.
Every fiber over $v$ with $v(\mathfrak c)<+\infty$ is a singleton.
Put
\[
 Z_R=\{v\in\mathcal V_{\mathfrak a}(R):v(\mathfrak c)=+\infty\},
 \qquad Z_B=\pi_{\mathfrak a}^{-1}(Z_R).
\]
These are closed. On $Z_B$, all lifts factor through $S/\mathfrak cS$,
and equality of their restrictions to $R/\mathfrak c$ is exactly the
remaining equivalence relation. There is a topological pushout
\[
              \mathcal V_{\mathfrak a}(R)
                  \simeq\mathcal B_{\mathfrak a}\sqcup_{Z_B}Z_R.
\]
Restriction, the conductor strata, and this quotient presentation commute
with every primary change \eqref{eq:p121-general-boundary-change}.
\end{theorem}

\begin{proof}
The continuous surjection in
\cref{lem:p121-normalization-semivaluation-lift} is a quotient map
because its domain is compact and its target is Hausdorff.
Its equivalence relation is the inverse image of the diagonal, hence closed.
Regard a point of $\mathcal B_{\mathfrak a,i}$ as a semivaluation on
$S$ by projection to $S_i$. This loses no information: multiplicativity
and $1=\sum_ie_i$ force a semivaluation on a finite product to take
value zero on exactly one primitive idempotent $e_i$ and infinity on
all the others.

If $v(\mathfrak c)<+\infty$, choose $c\in\mathfrak c$ with
$v(c)<+\infty$. Since $cs\in R$ for every $s\in S$, any lift satisfies
\[
                         w(s)=v(cs)-v(c).
\]
There is no subtraction of infinite quantities here. This determines
both the component and the lift uniquely. If $v(\mathfrak c)=+\infty$,
both $v$ and its lifts annihilate the conductor. Their equality on $R$
is precisely equality after restriction to $R/\mathfrak c$, consistently
with \eqref{eq:p121-conductor-pullback-square}. Finite generation of
$\mathfrak c$ makes the infinite-value condition closed. The indicated
pushout is a quotient of the compact space
$\mathcal B_{\mathfrak a}\coprod Z_R$, and its induced map to
$\mathcal V_{\mathfrak a}(R)$ is a continuous bijection. Compactness
and Hausdorffness make it a homeomorphism.
Finally $w(\mathfrak bS_i)=v(\mathfrak b)$, so all changes of section
commute with restriction and preserve the conductor strata.
\end{proof}

\begin{example}\label{ex:p121-surface-conductor-fiber}
For $R=k[[x,y,z]]/(xy)$, normalization is
$k[[x,z]]\times k[[y,z]]$. The image consists of pairs agreeing
on the $z$-axis, and the conductor is $(x,y)$: multiplying a pair
by the two primitive idempotents tests that its common $z$-axis
restriction vanishes. The semivaluation obtained from
$R/(x,y)=k[[z]]$ by $\ord_z$ has $v(\mathfrak m)=1$ and
$v(\mathfrak c)=+\infty$. It has two distinct lifts, with
$w(x)=+\infty$ on the first branch and $w(y)=+\infty$ on the
second. Thus conductor identifications really occur in dimension two.
The theorem confines identifications to the conductor stratum; it does
not assert that every point of that stratum has more than one lift.
\end{example}

\begin{proposition}\label{prop:p121-blowup-boundary-gauge}
Let $I=(f_1,\ldots,f_N)$ be an ideal of $R$. If $v\in\mathcal C_R$
and $v(I)<+\infty$, then $v$ has a canonical center on
$\operatorname{Bl}_I(\operatorname{Spec}R)$, and its value on the
invertible ideal $I\mathcal O_{\operatorname{Bl}_I R}$ is $v(I)$.
\end{proposition}

\begin{proof}
Choose $j$ with $v(f_j)=v(I)<+\infty$. The relevant blowup chart is
\[
                  R[f_1/f_j,\ldots,f_N/f_j]\subset R[1/f_j].
\]
The valuation ring of $\operatorname{Frac}(R/\operatorname{supp}v)$
receives this chart, since $v(f_i/f_j)=v(f_i)-v(f_j)\geq0$.
It determines a center there. Two minimizing generators have a ratio
of value zero, so the centers agree on their chart overlap. A generator
with larger finite value gives a negative-valued ratio and cannot give
a center; a generator of infinite value vanishes at the generic
support point and its chart cannot receive that point. On a minimizing
chart the pulled-back ideal is generated by $f_j$, of value $v(I)$.
\end{proof}

For a prime divisor $E$ over the closed point on a normal birational
model of a component of $\operatorname{Spec}R$, the change of primary boundary sends
\[
             \frac{\ord_E}{\ord_E(\mathfrak a)}
                    \longmapsto
             \frac{\ord_E}{\ord_E(\mathfrak b)}.
\]
Both denominators are positive and finite. On a chart where the two
pulled-back ideals are principal monomials
$u\prod_i z_i^{a_i}$ and $u'\prod_i z_i^{b_i}$, with $u,u'$ units,
a monomial valuation of weights $w_i$ satisfies
\[
                  v(\mathfrak a)=\sum_i a_iw_i,
                  \qquad v(\mathfrak b)=\sum_i b_iw_i.
\tag{70.7d}\label{eq:p121-monomial-gauge-slices}
\]
The two slices are the corresponding level-one hyperplanes; their change
rescales weights by $(\sum_i b_iw_i)^{-1}$. These coordinate formulas
require the named ideals to be monomial on the chosen chart. They do not
claim a resolution in arbitrary characteristic. Without that hypothesis,
the intrinsic ideal-value formula still applies, but coefficients along
the chart boundary alone need not compute a nonmonomial ideal's value.
For instance, in $k[[x,y]]$ the primary ideal $(x+y,x^2)$ has order
zero along both coordinate divisors, but value one at the monomial
valuation of weights $(1,1)$.

For a proper nonzero ideal $\mathfrak d$, the exact domain of its
normalization is the open scaling-invariant subcone
\[
 \mathcal C_R(\mathfrak d)
       =\{v\in\mathcal C_R:v(\mathfrak d)<+\infty\}.
\tag{70.7e}\label{eq:p121-nonprimary-gauge-domain}
\]
Indeed, $v(\mathfrak d)\geq v(\mathfrak m)>0$, and its finite-value
locus is open. Rescaling gives a homeomorphism from the corresponding
open subset of any primary section onto
$\mathcal V_{\mathfrak d}(R)=\{v\in\mathcal C_R:v(\mathfrak d)=1\}$.
Its inverse rescales by the primary ideal value. For arbitrary ideals
$\mathfrak d,\mathfrak e$, the same formula defines a partial change
on $\{v\in\mathcal V_{\mathfrak d}:0<v(\mathfrak e)<+\infty\}$,
with image $\{u\in\mathcal V_{\mathfrak e}:0<u(\mathfrak d)<+\infty\}$.
The inverse and cocycle identities hold on these domains and commute
with normalization-branch restriction. The unit ideal has zero gauge,
the zero ideal has infinite gauge, and neither has a level-one section.

\begin{example}\label{ex:p121-nonprimary-noncompact-link}
For $R=k[[x,y]]$ and $\mathfrak d=(x)$, the semivaluation from
$R/(x)=k[[y]]$ and $\ord_y$ lies in $\mathcal V_{\mathfrak m}(R)$
but has infinite $\mathfrak d$-value. Its ray is omitted by
\eqref{eq:p121-nonprimary-gauge-domain}. The monomial valuations
$v_j(x)=1$, $v_j(y)=1/j$ belong to $\mathcal V_{\mathfrak d}(R)$.
If this link were compact, its image under continuous $y$-evaluation
would be a closed compact subset of $[0,+\infty]$ containing every
$1/j$, hence zero. But every centered semivaluation is strictly positive
on $y\in\mathfrak m$. Thus this link is not compact.
\end{example}

Finally let $R$ be nonreduced, with nilradical $\mathcal N$.
If $n^k=0$, then $kv(n)=+\infty$, so $v(n)=+\infty$.
Applying the ultrametric inequality to $f+n$ and to $(f+n)-n$
gives $v(f+n)=v(f)$. Consequently reduction induces a homeomorphism
of centered cones and preserves every ideal gauge and change of section.
If $\dim R_{\mathrm{red}}=0$, that reduced local ring is a field,
so its centered cone is empty. Nilpotent multiplicities cannot be
recovered from this datum and must be retained separately for any
intersection or measure construction.

These results settle the normalizing-boundary question for excellent
local germs: the ray space, primary sections, partial nonprimary
sections, and conductor descent are explicit. They do not construct
a singular psh class, envelopes, products, energy, or general finite-map
pluripotential functoriality.


\subsection{Reduced curve links and ideal profiles}
\label{subsec:local-reduced-curve-profiles}

For a reduced singular curve germ the normalized boundary can already have
several points, even when the germ is irreducible.  In this dimension the
boundary datum and its scaling law admit a complete answer: the points are
the branches of the finite normalization, and a primary ideal fixes the
scale independently on each branch.  This is the first singular range of
the general boundary construction in
\cref{thm:p121-excellent-primary-links,thm:p121-excellent-conductor-descent}.
The
conductor-controlled result below also computes one explicit principal-ideal
profile family.  Its uniform closure and decreasing completion give a
book-local ideal-psh candidate and complete envelope calculus on the finite
link, with an intrinsic homogeneous realization on the branch-valuation
cone.  The construction is not identified with a separately defined
nef-b-divisor or analytic psh class, and it supplies no product or energy.

% Book-local P12.1 curve-link classification, 11 September 2026.
% The result identifies the boundary set and its normalization scale only;
% it does not construct a singular local psh or Monge--Ampere theory.
\begin{theorem}[Normalized link of a reduced Nagata curve germ]
\label{thm:p121-reduced-nagata-curve-link}
Let \((R,\mathfrak m)\) be a one-dimensional reduced Nagata noetherian
local ring, and let \(\mathfrak a\subset R\) be an
\(\mathfrak m\)-primary ideal.  Write \(R^\nu\) for the normalization of
\(R\) in its total ring of fractions and put
\[
 \operatorname{Br}(R,\mathfrak m)
 :=\{\mathfrak q\in\operatorname{Max}(R^\nu):
                         \mathfrak q\cap R=\mathfrak m\}.
\]
For \(\mathfrak q\in\operatorname{Br}(R,\mathfrak m)\), the local ring
\(D_{\mathfrak q}:=(R^\nu)_{\mathfrak q}\) is a discrete valuation
ring.  Pull its normalized order back to a semivaluation
\[
 w_{\mathfrak q}:R\longrightarrow
                    \mathbf Z_{\geq0}\cup\{+\infty\},
\]
where the value is \(+\infty\) on the kernel of
\(R\to D_{\mathfrak q}\), and set
\[
 d_{\mathfrak q}(\mathfrak a)
 :=w_{\mathfrak q}(\mathfrak a)
 =\min_{f\in\mathfrak a}w_{\mathfrak q}(f).
\]
Then \(d_{\mathfrak q}(\mathfrak a)\) is a positive integer, and
\[
 \operatorname{Br}(R,\mathfrak m)\xrightarrow{\ \sim\ }
 \mathcal V_{\mathfrak a}(R),
 \qquad
 \mathfrak q\longmapsto
       \frac{w_{\mathfrak q}}{d_{\mathfrak q}(\mathfrak a)},
 \tag{\(\star\)}\label{eq:p121-curve-branch-link}
\]
is a bijection.  Here \(\mathcal V_{\mathfrak a}(R)\) is the set of
nonnegative real semivaluations centred at \(\mathfrak m\) and normalized
by \(v(\mathfrak a)=1\).  With the topology of pointwise convergence on
\(R\), it is a finite discrete space.

For another \(\mathfrak m\)-primary ideal \(\mathfrak b\), the change of
normalization is the canonical homeomorphism
\[
 H_{\mathfrak a,\mathfrak b}:\mathcal V_{\mathfrak a}(R)
       \longrightarrow\mathcal V_{\mathfrak b}(R),
 \qquad
 H_{\mathfrak a,\mathfrak b}(v)=\frac{v}{v(\mathfrak b)}.
 \tag{\(\dagger\)}\label{eq:p121-change-normalizing-ideal}
\]
It fixes the branch in \eqref{eq:p121-curve-branch-link}, has inverse
\(H_{\mathfrak b,\mathfrak a}\), and satisfies
\[
 H_{\mathfrak b,\mathfrak c}\circ H_{\mathfrak a,\mathfrak b}
       =H_{\mathfrak a,\mathfrak c}
\]
for every third \(\mathfrak m\)-primary ideal \(\mathfrak c\).
In particular, the theorem applies to every one-dimensional reduced
excellent local complex germ.
\end{theorem}

\begin{proof}
The normalization is finite over \(R\) and is a semilocal finite product of
Dedekind domains by \cite[Tag~0C1R]{StacksProject}; hence the displayed
branch set is finite and every \(D_{\mathfrak q}\) is a one-dimensional
normal noetherian local domain.  It is therefore a discrete valuation ring
by \cite[Tag~00PD]{StacksProject}.
Moreover, \(\mathfrak a\) is contained in no minimal prime of \(R\), while
its image on every branch lies in the maximal ideal of
\(D_{\mathfrak q}\).  Thus
\(0<d_{\mathfrak q}(\mathfrak a)<+\infty\), as asserted.

Let \(v\in\mathcal V_{\mathfrak a}(R)\), and let
\(\mathfrak p=v^{-1}(+\infty)\) be its support.  Since
\(v(\mathfrak a)=1\), the ideal \(\mathfrak a\) is not contained in
\(\mathfrak p\).  The equality
\(\sqrt{\mathfrak a}=\mathfrak m\) therefore rules out
\(\mathfrak p=\mathfrak m\).  Every nonmaximal prime of a
one-dimensional local ring is minimal, so \(\mathfrak p\) is a minimal
prime of \(R\).

The semivaluation factors through \(A=R/\mathfrak p\) and extends uniquely
to a real valuation of \(K=\operatorname{Frac}(A)\).  Let \(V\subset K\)
be its valuation ring.  A valuation ring is normal
\cite[Tag~00IC]{StacksProject}, so \(V\) contains the integral closure
\(\overline A\) of \(A\) in \(K\).  The centre
\[
                 \overline{\mathfrak q}
                  :=\mathfrak m_V\cap\overline A
\]
lies over \(\mathfrak m/\mathfrak p\).  It determines a unique
\(\mathfrak q\in\operatorname{Br}(R,\mathfrak m)\), and
\(D_{\mathfrak q}=\overline A_{\overline{\mathfrak q}}\subset V\), with
\(V\) dominating \(D_{\mathfrak q}\).

Choose a uniformizer \(\pi\) of \(D_{\mathfrak q}\).  Every nonzero
element of \(K\) has a unique expression \(u\pi^n\), where
\(u\in D_{\mathfrak q}^{\times}\) and \(n\in\mathbf Z\).  Domination gives
\(v(u)=0\) and \(c:=v(\pi)>0\), whence
\[
                              v=cw_{\mathfrak q}.
\]
Thus every real valuation occurring here has a discrete cyclic value group.
The normalization \(v(\mathfrak a)=1\) forces
\(c=d_{\mathfrak q}(\mathfrak a)^{-1}\).  Conversely, every normalized
branch order in \eqref{eq:p121-curve-branch-link} is centred at
\(\mathfrak m\) and has \(\mathfrak a\)-value one.  This proves
surjectivity and the asserted formula.

Branches with different supports are separated by their infinite loci.  If
two branches have the same support and their normalized orders agree on
\(R\), then their extensions to \(K\) are positive multiples of each other
and have the same valuation ring; their centres on \(\overline A\) are
therefore equal.  Hence the map is injective.  Pointwise convergence embeds
the finite set into the Hausdorff product
\([0,+\infty]^R\), so the resulting topology is discrete.

If \(\mathfrak b\) is \(\mathfrak m\)-primary, then every branch value
\(v(\mathfrak b)\) is finite and strictly positive.  On the branch
\(\mathfrak q\), formula \eqref{eq:p121-change-normalizing-ideal} reads
\[
 \frac{w_{\mathfrak q}}{d_{\mathfrak q}(\mathfrak a)}
 \longmapsto
 \frac{w_{\mathfrak q}}{d_{\mathfrak q}(\mathfrak b)}.
\]
The inverse and cocycle identities follow immediately.  Finally, every
excellent ring is Nagata by \cite[Tag~07QV]{StacksProject}, proving the last
assertion.
\end{proof}

% Book-local P12.6 normalized-link functoriality for finite curve maps.
\begin{proposition}[Finite curve maps on normalized branch links]
\label{prop:p126-finite-reduced-curve-link}
Let
\[
 \varphi:(R,\mathfrak m)\longrightarrow(S,\mathfrak n)
\]
be a finite local homomorphism between one-dimensional reduced Nagata
noetherian local rings, and let \(\mathfrak a\subset R\) be
\(\mathfrak m\)-primary.  Then \(\mathfrak aS\) is
\(\mathfrak n\)-primary and restriction defines a continuous map
\[
 \varphi^*: \mathcal V_{\mathfrak aS}(S)\longrightarrow
             \mathcal V_{\mathfrak a}(R),
 \qquad w\longmapsto w\circ\varphi.
 \tag{\(\ddagger\)}\label{eq:p126-finite-curve-link-map}
\]

More precisely, let \(\mathfrak q'\in\operatorname{Br}(S,\mathfrak n)\).
It lies over a unique branch
\(\mathfrak q\in\operatorname{Br}(R,\mathfrak m)\), and the induced local
map of discrete valuation rings has a ramification index
\(e_{\mathfrak q'/\mathfrak q}\geq1\).  With the notation of
\Cref{thm:p121-reduced-nagata-curve-link},
\[
 \begin{aligned}
 w_{\mathfrak q'}(\varphi(f))
   &=e_{\mathfrak q'/\mathfrak q}\,w_{\mathfrak q}(f)
       &&(f\in R),\\
 d_{\mathfrak q'}(\mathfrak aS)
   &=e_{\mathfrak q'/\mathfrak q}\,
       d_{\mathfrak q}(\mathfrak a),\\
 \varphi^*\!\left(
   \frac{w_{\mathfrak q'}}{d_{\mathfrak q'}(\mathfrak aS)}
                 \right)
   &=\frac{w_{\mathfrak q}}{d_{\mathfrak q}(\mathfrak a)}.
 \end{aligned}
 \tag{\(\heartsuit\)}\label{eq:p126-ramification-cancellation}
\]
Here the first equality uses the convention
\(e\cdot(+\infty)=+\infty\).  These maps commute with composition of
finite local homomorphisms and with the change-of-ideal maps
\eqref{eq:p121-change-normalizing-ideal}.
\end{proposition}

\begin{proof}
The quotient \(S/\mathfrak aS\) is finite over the artinian ring
\(R/\mathfrak a\), hence is artinian.  It is a nonzero local ring, since
\(\mathfrak aS\subseteq\mathfrak mS\subseteq\mathfrak n\).  Its unique
prime is therefore \(\mathfrak n/\mathfrak aS\), proving that
\(\mathfrak aS\) is \(\mathfrak n\)-primary.  For every nonnegative
semivaluation \(w\) on \(S\),
\[
 w(\mathfrak aS)=\min_{f\in\mathfrak a}w(\varphi(f)),
 \tag{\(\clubsuit\)}\label{eq:p126-extended-ideal-value}
\]
because coefficients from \(S\) have nonnegative value and the coefficient
one occurs.  Hence restriction preserves the displayed normalization and
is continuous for pointwise convergence.  Since \(\varphi\) is local,
positivity on \(\mathfrak n\) restricts to positivity on \(\mathfrak m\), so
the restricted semivaluation remains centred at the closed point.

Let \(\mathfrak p'\) be the minimal prime of \(S\) corresponding to
\(\mathfrak q'\), and put
\(\mathfrak p=\varphi^{-1}(\mathfrak p')\).  The induced map
\[
             R/\mathfrak p\lhook\joinrel\longrightarrow
             S/\mathfrak p'
\]
is finite and integral.  Since the target is a one-dimensional domain,
dimension invariance for integral injective extensions
\cite[Tag~00OK]{StacksProject} shows that
\(\dim(R/\mathfrak p)=1\).  Thus \(\mathfrak p\) is a minimal prime of
\(R\).

The inclusion of the component fraction fields carries the normalization
of \(R/\mathfrak p\) into that of \(S/\mathfrak p'\).  Contracting the
centre \(\mathfrak q'\) and localizing therefore gives an injective local
map
\[
                       D_{\mathfrak q}\longrightarrow D_{\mathfrak q'}
\]
of discrete valuation rings.  By the definition of ramification index
\cite[Tag~09E4]{StacksProject}, a uniformizer of the first ring is a unit
times the \(e_{\mathfrak q'/\mathfrak q}\)-th power of a uniformizer of
the second.  This proves the first line of
\eqref{eq:p126-ramification-cancellation}.  Taking the minimum over
\(\mathfrak a\) and using \eqref{eq:p126-extended-ideal-value} proves the
second; dividing cancels the ramification index and proves the third.

Restriction of functions is associative.  Equivalently, the branchwise
ramification indices multiply in a tower by
\cite[Tag~0BRL]{StacksProject}.  This proves compatibility with composition.
The branchwise description of
\eqref{eq:p121-change-normalizing-ideal} proves compatibility with changing
the normalizing ideal.
\end{proof}

\begin{example}[Branches, minimal primes, and normalization scale]
\label{ex:p121-node-cusp-normalized-links}
For the reducible node
\[
 R=\mathbf C[[x,y]]/(xy),
 \qquad \mathfrak m=(x,y),
\]
the normalization is \(\mathbf C[[x]]\times\mathbf C[[y]]\).  Hence
\(\mathcal V_{\mathfrak m}(R)\) consists of two points: the orders on the
two axes, with respective supports \((y)\) and \((x)\).

Minimal primes alone do not record all branches.  The irreducible nodal
local ring
\[
 R'=\left(\mathbf C[x,y]/(y^2-x^2(1+x))\right)_{(x,y)}
\]
is a domain, but its normalization has two points above the closed point.
Indeed, in its fraction field put \(t=y/x\); then
\(x=t^2-1\) and \(y=t(t^2-1)\), and the two centres are \(t=1\) and
\(t=-1\).  Thus \(\mathcal V_{\mathfrak m}(R')\) again has two points,
although \(R'\) has only one minimal prime.

For the cusp
\[
 R''=\mathbf C[[t^2,t^3]]\subset S=\mathbf C[[t]],
 \qquad \mathfrak m''=(t^2,t^3),
\]
there is one branch and
\[
                 \mathcal V_{\mathfrak m''}(R'')
                    =\left\{\tfrac12\operatorname{ord}_t\right\}.
\]
The maximal-ideal normalization on \(S\) instead gives
\(\mathcal V_{(t)}(S)=\{\operatorname{ord}_t\}\).  With the compatible
ideal \(\mathfrak m''S=(t^2)\), however, the upstairs normalized order is
\(\tfrac12\operatorname{ord}_t\), and
\Cref{prop:p126-finite-reduced-curve-link} restricts it exactly to the
point of \(\mathcal V_{\mathfrak m''}(R'')\).  This is the simplest reason
that the change-of-ideal factor in
\eqref{eq:p121-change-normalizing-ideal} cannot be omitted.
\end{example}

% Layout-only edition, 22 September 2026.
% Controlling mathematical source: editorial/solver_dossiers_20260911/p121_reduced_excellent_curve_links_20260911/CONDUCTOR_IDEALS.tex
% Source SHA-256: 73eb8ce8754222fda5623965f92a6eeefb536a33d5949109c8387aaafd7e732f
% Only display grouping and line breaks differ; the source is retained.
% Book-local P12.2/P12.3/P12.6 conductor-controlled curve-link result.
% This is an ideal-function theorem on a finite normalized link, not a
% construction of the full singular local psh or Monge--Ampere theory.
\begin{proposition}[Conductor realization of branch orders]
\label{prop:p126-curve-conductor-principal-realization}
Let \((R,\mathfrak m)\) be a one-dimensional reduced Nagata noetherian
local ring, put \(D=R^\nu\), and enumerate
\[
             \operatorname{Br}(R,\mathfrak m)
                     =\{\mathfrak q_1,\ldots,\mathfrak q_s\}.
\]
Let
\[
 \mathfrak c=(R:D):=\{b\in D:bD\subseteq R\}
\]
be the conductor.  For every \(i\), define \(c_i\geq0\) by
\[
 \mathfrak cD_{\mathfrak q_i}
     =\mathfrak q_i^{c_i}D_{\mathfrak q_i},
 \qquad C_i:=\max\{c_i,1\}.
 \tag{\(\clubsuit\)}\label{eq:p126-curve-conductor-orders}
\]
Then, for every tuple of integers \(k_i\geq C_i\), there is an
\(R\)-nonzerodivisor \(f\in\mathfrak m\) such that
\[
                              w_{\mathfrak q_i}(f)=k_i
                       \qquad(1\leq i\leq s).
 \tag{\(\spadesuit\)}\label{eq:p126-prescribed-branch-orders}
\]
In particular, \((f)\) is an \(\mathfrak m\)-primary principal ideal.
Thus every sufficiently deep integral effective divisor on the finite
normalization fibre is cut out by one nonzerodivisor downstairs.
\end{proposition}

\begin{proof}
By \Cref{thm:p121-reduced-nagata-curve-link}, the finite \(R\)-algebra
\(D\) is a semilocal finite product of Dedekind domains and the displayed
\(\mathfrak q_i\) are all of its maximal ideals.  The conductor contains a
nonzerodivisor.  Indeed, choose finite \(R\)-module generators
\(b_1,\ldots,b_t\) of \(D\).  Since each \(b_j\) lies in the total ring of
fractions of \(R\), there is an \(R\)-nonzerodivisor \(a_j\) with
\(a_jb_j\in R\).  Their product \(a=\prod_ja_j\) is again a
nonzerodivisor and satisfies \(aD\subseteq R\), so \(a\in\mathfrak c\).

Consequently \(\mathfrak c\) is a nonzero ideal on every Dedekind factor.
Unique ideal factorization gives
\[
                         \mathfrak c
                    =\prod_{i=1}^s\mathfrak q_i^{c_i}
\]
with the exponents in \eqref{eq:p126-curve-conductor-orders}; see
\cite[Tags~034W and~034X]{StacksProject}.  For the prescribed tuple put
\[
                              J=\prod_{i=1}^s\mathfrak q_i^{k_i}.
\]
Then \(J\subseteq\mathfrak c\subseteq R\).  The ideal \(J\) is invertible,
hence each \(J/\mathfrak q_iJ\) is a nonzero one-dimensional vector space
over \(D/\mathfrak q_i\).  The Chinese remainder theorem for the pairwise
comaximal maximal ideals \(\mathfrak q_i\) lets us choose \(f\in J\) whose
class is nonzero in every \(J/\mathfrak q_iJ\).  Nakayama's lemma then shows
that \(f\) generates \(J_{\mathfrak q_i}\) over \(D_{\mathfrak q_i}\) for
every \(i\).  These are all the maximal ideals of \(D\), so \(J/fD=0\) and
\(J=fD\).  The inclusion \(J\subseteq R\) gives \(f\in R\), while
localization at \(\mathfrak q_i\) gives
\[
                         fD_{\mathfrak q_i}
                    =\mathfrak q_i^{k_i}D_{\mathfrak q_i},
\]
which is \eqref{eq:p126-prescribed-branch-orders}.

Multiplication by this basis element identifies \(D\) with \(J\), so \(f\)
is a nonzerodivisor of \(D\), and hence of \(R\subseteq D\).  Since every
\(k_i\geq1\), the element \(f\) lies in every \(\mathfrak q_i\), and thus
in \(\mathfrak m\).  A prime of the one-dimensional local ring \(R\)
which contains \(f\) cannot be minimal, because \(f\) is a
nonzerodivisor.  It must therefore be \(\mathfrak m\).  Hence
\(\sqrt{(f)}=\mathfrak m\); a noetherian ideal with maximal radical is
primary.
\end{proof}

\begin{theorem}[Principal-ideal profiles and envelopes on a curve link]
\label{thm:p122-curve-principal-ideal-envelope}
In the setting of
\Cref{prop:p126-curve-conductor-principal-realization}, let
\(\mathfrak a\subset R\) be \(\mathfrak m\)-primary and write
\[
 d_i=w_{\mathfrak q_i}(\mathfrak a),
 \qquad v_i=\frac{w_{\mathfrak q_i}}{d_i}
       \in\mathcal V_{\mathfrak a}(R).
\]
For an \(R\)-nonzerodivisor \(f\) such that \((f)\) is
\(\mathfrak m\)-primary and an integer \(N\geq1\), put
\[
 u_{f,N}(v):=\frac1N\log|(f)|(v)=-\frac{v(f)}N,
 \qquad
 \mathscr P_{\mathfrak a}^{\mathrm{pr}}
 :=\{u_{f,N}:f,N\text{ as above}\}.
 \tag{\(\diamondsuit\)}\label{eq:p122-principal-ideal-profile-family}
\]
Then the following assertions hold.
\begin{enumerate}
\item The uniform closure of this explicitly defined family is the complete
      nonpositive orthant.  More precisely,
      \[
       \begin{gathered}
       \mathscr P_{\mathfrak a}^{\mathrm{pr}}
       =\{u:u(v_i)\in\mathbf Q_{<0}\text{ for every }i\},\\
       \qquad
       \overline{\mathscr P_{\mathfrak a}^{\mathrm{pr}}}^{\,\|\cdot\|_\infty}
       =\{u:\mathcal V_{\mathfrak a}(R)\to\mathbf R:u\leq0\}.
       \end{gathered}
       \tag{\(\heartsuit\)}\label{eq:p122-principal-ideal-orthant}
      \]
\item Every nonpositive real-valued function is the increasing uniform limit
      of a sequence in \(\mathscr P_{\mathfrak a}^{\mathrm{pr}}\).  Every
      strictly negative real-valued function is also the decreasing uniform
      limit of such a sequence.  If \(u(v_i)=0\) for some branch, no
      decreasing sequence in \(\mathscr P_{\mathfrak a}^{\mathrm{pr}}\)
      can converge pointwise to \(u\).
\item For every finite real-valued obstacle \(\phi\) on the link, define the
      principal-ideal envelope by
      \[
       P_{\mathfrak a}^{\mathrm{pr}}\phi(v)
       :=\sup\{u(v):u\in\mathscr P_{\mathfrak a}^{\mathrm{pr}},\ u\leq\phi\}.
       \tag{\(\star\)}\label{eq:p123-principal-ideal-envelope}
      \]
      Then
      \[
                         P_{\mathfrak a}^{\mathrm{pr}}\phi
                              =\min\{\phi,0\}.
       \tag{\(\dagger\)}\label{eq:p123-principal-ideal-envelope-formula}
      \]
      The supremum in \eqref{eq:p123-principal-ideal-envelope} admits an
      increasing uniformly convergent sequence of principal-ideal
      competitors.  The same formula results if the competitor family is
      replaced by its uniform closure in
      \eqref{eq:p122-principal-ideal-orthant}.  The original supremum need
      not be attained by a member of
      \(\mathscr P_{\mathfrak a}^{\mathrm{pr}}\).
\end{enumerate}
These are statements about the displayed ideal-function family on a finite
link.  They do not identify that family with an independently constructed
singular nef or psh cone and do not construct a Monge--Ampere product or an
energy class.
\end{theorem}

\begin{proof}
Write \(u(v_i)=-\alpha_i\), where \(\alpha_i\geq0\).  With \(C_i\) as in
\eqref{eq:p126-curve-conductor-orders}, set
\[
 k_{n,i}=\left\lceil2^nd_i\alpha_i\right\rceil+C_i.
 \]
The conductor realization proposition supplies \(f_n\) with these branch
orders.  For \(N=2^n\),
\[
 u_{f_n,2^n}(v_i)=-\frac{k_{n,i}}{2^nd_i}\leq-\alpha_i=u(v_i),
 \qquad
 0\leq u(v_i)-u_{f_n,2^n}(v_i)
       <\frac{C_i+1}{2^nd_i}.
 \tag{\(\ddagger\)}\label{eq:p122-increasing-principal-approximation}
\]
Since \(\lceil2x\rceil\leq2\lceil x\rceil\), one has
\(k_{n+1,i}\leq2k_{n,i}\).  Thus the functions in
\eqref{eq:p122-increasing-principal-approximation} increase pointwise, and
the finite branch set makes the convergence uniform.

If every \(\alpha_i>0\), then for all sufficiently large \(n\) the integers
\[
                    \ell_{n,i}=\left\lfloor2^nd_i\alpha_i\right\rfloor
\]
satisfy \(\ell_{n,i}\geq C_i\).  Choose \(g_n\) with those branch orders.
Then
\[
 u_{g_n,2^n}(v_i)=-\frac{\ell_{n,i}}{2^nd_i}\geq u(v_i),
 \qquad
 0\leq u_{g_n,2^n}(v_i)-u(v_i)<\frac1{2^nd_i}.
\]
The inequality \(\lfloor2x\rfloor\geq2\lfloor x\rfloor\) shows that this
sequence decreases.  Conversely, every member of
\(\mathscr P_{\mathfrak a}^{\mathrm{pr}}\) is strictly negative at every
\(v_i\).  A decreasing sequence starting with such a function therefore
cannot have limit zero at any branch.

If all \(\alpha_i\) are positive rationals, choose one sufficiently
divisible \(N\) such that every \(k_i=Nd_i\alpha_i\) is an integer at least
\(C_i\).  A single application of the conductor proposition then gives
\(u=u_{f,N}\).  Conversely, every member of the displayed family has
strictly negative rational coordinates.  Together with the increasing
approximation, this proves
\eqref{eq:p122-principal-ideal-orthant} and the first two assertions.

Finally put \(w=\min\{\phi,0\}\).  The increasing construction applied to
\(w\) gives principal-ideal profiles \(u_n\leq w\leq\phi\) converging
uniformly to \(w\).  Hence the envelope is at least \(w\).  Every admissible
profile is both nonpositive and bounded above by \(\phi\), so it is bounded
above by \(w\); this proves
\eqref{eq:p123-principal-ideal-envelope-formula}.  The closure version is
identical by \eqref{eq:p122-principal-ideal-orthant}.
\end{proof}

% Book-local P12.2/P12.3 curve-level ideal-psh construction.
% The superscript ``id'' records that this is the closure of the displayed
% ideal-model family; no comparison with a separate nef or CLD psh cone is
% asserted.
\begin{theorem}[Ideal-psh completion and envelopes on reduced curve links]
\label{thm:p122-curve-ideal-psh-functoriality}
Let \((R,\mathfrak m)\) be a one-dimensional reduced Nagata noetherian
local ring, let \(\mathfrak a\subset R\) be \(\mathfrak m\)-primary, and
put \(V_{\mathfrak a}:=\mathcal V_{\mathfrak a}(R)\).  With the
principal-ideal model family
\(\mathscr P_{\mathfrak a}^{\mathrm{pr}}\) of
\eqref{eq:p122-principal-ideal-profile-family}, write
\(u^{\mathfrak a}_{f,N}\) when its normalizing ideal must be visible, and
define
\[
 \operatorname{CPSH}^{\mathrm{id}}_{\mathfrak a}(R)
 :=\overline{\mathscr P_{\mathfrak a}^{\mathrm{pr}}}
                         ^{\,\|\cdot\|_\infty},
 \qquad
 \operatorname{PSH}^{\mathrm{id},\bot}_{\mathfrak a}(R)
 :=\left\{\lim_{j\to\infty}u_j:
       u_j\in\operatorname{CPSH}^{\mathrm{id}}_{\mathfrak a}(R),\
       u_j\searrow\right\}.
 \tag{\(\triangle\)}\label{eq:p122-curve-ideal-psh-definition}
\]
The second class retains the identically \(-\infty\) bottom element;
removing it gives
\(\operatorname{PSH}^{\mathrm{id}}_{\mathfrak a}(R)\).  Then:
\begin{enumerate}
\item The two classes have the intrinsic descriptions
      \[
       \operatorname{CPSH}^{\mathrm{id}}_{\mathfrak a}(R)
          =\{u:V_{\mathfrak a}\to\mathbf R:u\leq0\},
       \qquad
       \operatorname{PSH}^{\mathrm{id},\bot}_{\mathfrak a}(R)
          =\{u:V_{\mathfrak a}\to[-\infty,0]\}.
       \tag{\(\square\)}\label{eq:p122-curve-ideal-psh-orthants}
      \]
      In particular, the continuous class is a closed convex cone stable
      under finite maxima, and the completed class is a complete lattice.
      Replacing the principal model family by all positive rational
      multiples of \(\log|\mathfrak b|\), with \(\mathfrak b\) ranging
      over the \(\mathfrak m\)-primary ideals, gives the same closure.
\item For an arbitrary extended-real obstacle
      \(\phi:V_{\mathfrak a}\to[-\infty,+\infty]\), its completed
      ideal-psh envelope
      \[
       P^{\mathrm{id}}_{\mathfrak a}\phi
       :=\sup\{u\in
          \operatorname{PSH}^{\mathrm{id},\bot}_{\mathfrak a}(R):
          u\leq\phi\}
      \]
      is attained and is given coordinatewise by
      \[
                       P^{\mathrm{id}}_{\mathfrak a}\phi
                              =\min\{\phi,0\}.
       \tag{\(\lozenge\)}\label{eq:p123-curve-ideal-psh-envelope}
      \]
      For finite \(\phi\), it belongs to
      \(\operatorname{CPSH}^{\mathrm{id}}_{\mathfrak a}(R)\), and the
      same function is the supremum of principal-ideal competitors in
      \eqref{eq:p123-principal-ideal-envelope}.
\item These constructions are independent of the chosen primary-ideal
      cross-section with an explicit homogeneous realization.  Put
      \[
       \mathcal C_R^{\mathrm{br}}
          :=\{t w_{\mathfrak q}:t>0,\
                  \mathfrak q\in\operatorname{Br}(R,\mathfrak m)\}.
      \]
      Restriction to \(V_{\mathfrak a}\) identifies the completed class
      in \eqref{eq:p122-curve-ideal-psh-orthants} with the nonpositive
      degree-one homogeneous functions
      \[
       G:\mathcal C_R^{\mathrm{br}}\longrightarrow[-\infty,0],
       \qquad G(t\nu)=tG(\nu).
       \tag{\(\circ\)}\label{eq:p122-curve-homogeneous-branch-psh}
      \]
      Finite-valued homogeneous functions correspond to the continuous
      class.  In particular, this homogeneous object is independent of
      \(\mathfrak a\).  For another
      \(\mathfrak m\)-primary ideal \(\mathfrak b\), let
      \(H_{\mathfrak a,\mathfrak b}\) be
      \eqref{eq:p121-change-normalizing-ideal} and define
      \[
       (T_{\mathfrak a,\mathfrak b}u)
          \bigl(H_{\mathfrak a,\mathfrak b}(v)\bigr)
          :=\frac{u(v)}{v(\mathfrak b)}.
       \tag{\(\diamondsuit\)}\label{eq:p122-curve-psh-normalization-transport}
      \]
      Division by a positive number fixes both infinities.  The map
      \(T_{\mathfrak a,\mathfrak b}\) is an order and lattice
      isomorphism on both classes in
      \eqref{eq:p122-curve-ideal-psh-definition}, carries each
      principal-ideal profile to the profile of the same pair \((f,N)\),
      extends by the same formula to extended-real obstacles, and satisfies
      \[
       T_{\mathfrak b,\mathfrak c}T_{\mathfrak a,\mathfrak b}
          =T_{\mathfrak a,\mathfrak c},
       \qquad
       T_{\mathfrak a,\mathfrak b}
          (P^{\mathrm{id}}_{\mathfrak a}\phi)
          =P^{\mathrm{id}}_{\mathfrak b}
             (T_{\mathfrak a,\mathfrak b}\phi).
       \tag{\(\clubsuit\)}\label{eq:p122-curve-psh-transport-cocycle}
      \]
\item Let
      \(\varphi:(R,\mathfrak m)\to(S,\mathfrak n)\) be a finite local
      homomorphism between one-dimensional reduced Nagata noetherian local
      rings, and let
      \(\rho_{\mathfrak a}:V_{\mathfrak aS}\to V_{\mathfrak a}\) be the
      restriction map of \Cref{prop:p126-finite-reduced-curve-link}.  Then
      \(\rho_{\mathfrak a}^*u=u\circ\rho_{\mathfrak a}\) preserves the
      continuous and completed ideal-psh classes (including the bottom
      element), sends
      \[
                u^{\mathfrak a}_{f,N}
                    \longmapsto
                u^{\mathfrak aS}_{\varphi(f),N},
      \]
      and commutes with completed envelopes:
      \[
       \rho_{\mathfrak a}^*
          (P^{\mathrm{id}}_{\mathfrak a}\phi)
        =P^{\mathrm{id}}_{\mathfrak aS}
          (\rho_{\mathfrak a}^*\phi).
       \tag{\(\spadesuit\)}\label{eq:p123-curve-psh-finite-map-envelope}
      \]
      Pullback is compatible with composition and with the transports
      \eqref{eq:p122-curve-psh-normalization-transport}; explicitly, for
      every second \(\mathfrak m\)-primary ideal \(\mathfrak b\),
      \[
       \rho_{\mathfrak b}^*T^R_{\mathfrak a,\mathfrak b}
          =T^S_{\mathfrak aS,\mathfrak bS}
             \rho_{\mathfrak a}^*.
      \tag{\(\bigstar\)}\label{eq:p122-curve-psh-map-normalization-square}
      \]
      On either class, the image of \(\rho_{\mathfrak a}^*\) consists
      exactly of the functions on \(V_{\mathfrak aS}\) that are constant on
      every fiber of \(\rho_{\mathfrak a}\).  Consequently,
      \(\rho_{\mathfrak a}^*\) is injective if and only if
      \(\rho_{\mathfrak a}\) is surjective, and it is surjective if and only
      if \(\rho_{\mathfrak a}\) is injective.  When
      \(\rho_{\mathfrak a}\) is surjective, pullback on the continuous class
      is an isometry:
      \[
       \|\rho_{\mathfrak a}^*u-\rho_{\mathfrak a}^*v\|_\infty
          =\|u-v\|_\infty.
      \]
      On the completed lattices, pullback has the explicit left and right
      adjoints
      \[
       (\rho_{\mathfrak a,!}w)(x)
          :=\sup\bigl(\{w(y):\rho_{\mathfrak a}(y)=x\}
                        \cup\{-\infty\}\bigr),
      \]
      \[
       (\rho_{\mathfrak a,*}w)(x)
          :=\inf\bigl(\{w(y):\rho_{\mathfrak a}(y)=x\}
                        \cup\{0\}\bigr),
      \]
      so that
      \[
       \rho_{\mathfrak a,!}w\leq u
          \Longleftrightarrow
       w\leq\rho_{\mathfrak a}^*u,
       \qquad
       \rho_{\mathfrak a}^*u\leq w
          \Longleftrightarrow
       u\leq\rho_{\mathfrak a,*}w.
       \tag{\(\heartsuit\)}\label{eq:p122-curve-psh-finite-map-adjoints}
      \]
      If \(\rho_{\mathfrak a}\) is surjective, both composites with
      pullback on the right are the identity on the base lattice; if it is
      injective, both composites with pullback on the left are the identity
      on the source lattice.  For arbitrary \(\rho_{\mathfrak a}\),
      \(\rho_{\mathfrak a}^*\rho_{\mathfrak a,*}w\) is the greatest
      fiber-constant minorant of \(w\), while
      \(\rho_{\mathfrak a}^*\rho_{\mathfrak a,!}w\) is its least
      fiber-constant majorant.  Thus their common fixed-point set is exactly
      the pullback image.
\end{enumerate}
Thus the ideal-closure candidate has a complete psh-style order, envelope,
normalization, and pullback calculus on reduced Nagata curve links.  The
superscript \({\rm id}\) is a scope marker: naming this closure does not by
itself identify it with a separately defined nef-b-divisor, CLD, or
complex-analytic psh class, and hence does not by itself retire P12.2 or
P12.3.  The theorem constructs no local intersection product or energy.
\end{theorem}

\begin{proof}
The link is finite and discrete by
\Cref{thm:p121-reduced-nagata-curve-link}.  The first equality in
\eqref{eq:p122-curve-ideal-psh-orthants} is exactly the uniform-closure
calculation in
\eqref{eq:p122-principal-ideal-orthant}.  If
\(u:V_{\mathfrak a}\to[-\infty,0]\), then
\[
                            u_j:=\max\{u,-j\}
\]
is finite, nonpositive, and decreases pointwise to \(u\).  This proves the
second equality.  The asserted cone and lattice properties are now
coordinatewise.  Every \(\mathfrak m\)-primary ideal function has strictly
negative rational coordinates, so
\eqref{eq:p122-principal-ideal-orthant} realizes the same function by a
principal model; this proves the last assertion of the first item.

The function \(w=\min\{\phi,0\}\) belongs to the completed class, lies
below \(\phi\), and dominates every nonpositive competitor.  This proves
\eqref{eq:p123-curve-ideal-psh-envelope}.  If \(\phi\) is finite, then
\(w\) is finite and hence continuous on the finite link.  Its equality with
the raw principal-ideal supremum is
\eqref{eq:p123-principal-ideal-envelope-formula}.

Write
\(v_{\mathfrak q}^{\mathfrak a}
  =w_{\mathfrak q}/d_{\mathfrak q}(\mathfrak a)\).
Every nonpositive homogeneous \(G\) restricts to a member of the completed
class by \eqref{eq:p122-curve-ideal-psh-orthants}.  Conversely, a function
\(u\) in that class has the unique homogeneous extension
\[
 G_u(tw_{\mathfrak q})
    :=t d_{\mathfrak q}(\mathfrak a)
          u(v_{\mathfrak q}^{\mathfrak a}).
\]
The same formula preserves finite values.  This proves
\eqref{eq:p122-curve-homogeneous-branch-psh}.  Moreover,
\[
 v_{\mathfrak q}^{\mathfrak a}(\mathfrak b)
   =\frac{d_{\mathfrak q}(\mathfrak b)}
          {d_{\mathfrak q}(\mathfrak a)},
 \qquad
 H_{\mathfrak a,\mathfrak b}
       (v_{\mathfrak q}^{\mathfrak a})
   =v_{\mathfrak q}^{\mathfrak b}.
\]
Hence \eqref{eq:p122-curve-psh-normalization-transport} is multiplication
by one positive scalar on each coordinate.  It is therefore an order and
lattice isomorphism, preserves decreasing limits, and has inverse
\(T_{\mathfrak b,\mathfrak a}\).  Moreover,
\[
 \frac{u^{\mathfrak a}_{f,N}(v)}{v(\mathfrak b)}
   =-\frac{v(f)}{Nv(\mathfrak b)}
   =u^{\mathfrak b}_{f,N}
       \bigl(H_{\mathfrak a,\mathfrak b}(v)\bigr).
\]
For a third ideal \(\mathfrak c\), putting
\(w=H_{\mathfrak a,\mathfrak b}(v)\) gives
\(w(\mathfrak c)=v(\mathfrak c)/v(\mathfrak b)\); the two denominators
cancel and prove the first identity in
\eqref{eq:p122-curve-psh-transport-cocycle}.  The second follows from
\(T(\min\{\phi,0\})=\min\{T\phi,0\}\).

The restriction map \(\rho_{\mathfrak a}\) is a map of finite sets, so
composition with it preserves finite nonpositive functions and pointwise
decreasing limits.  If \((f)\) is \(\mathfrak m\)-primary and principal,
then \(f\) avoids every minimal prime of \(R\).  The proof of
\Cref{prop:p126-finite-reduced-curve-link} shows that every minimal prime
of \(S\) contracts to a minimal prime of \(R\).  Thus \(\varphi(f)\)
avoids every minimal prime of the reduced ring \(S\), so it is a
nonzerodivisor.  It lies in \(\mathfrak n\), and in the one-dimensional
local ring \(S\) its principal ideal is therefore \(\mathfrak n\)-primary.
For \(w\in V_{\mathfrak aS}\),
\[
 (\rho_{\mathfrak a}^*u^{\mathfrak a}_{f,N})(w)
       =-\frac{(w\circ\varphi)(f)}N
       =-\frac{w(\varphi(f))}N
       =u^{\mathfrak aS}_{\varphi(f),N}(w),
\]
which proves the model formula.  Equation
\eqref{eq:p123-curve-psh-finite-map-envelope} follows by pulling back the
coordinatewise minimum formula.  Associativity of function composition
gives compatibility in a tower.  Finally,
\(w(\mathfrak bS)=(w\circ\varphi)(\mathfrak b)\) by
\eqref{eq:p126-extended-ideal-value}; inserting this equality into
\eqref{eq:p122-curve-psh-normalization-transport} proves
\eqref{eq:p122-curve-psh-map-normalization-square}.

The descriptions in \eqref{eq:p122-curve-ideal-psh-orthants} reduce the
remaining claims to finite sets.  A target function is a pullback precisely
when it is constant on every fiber: necessity is immediate, and for
sufficiency define the source function on the image of
\(\rho_{\mathfrak a}\) by the common fiber value and give it the value zero
at every missed point.  If \(\rho_{\mathfrak a}\) is surjective, equality of
two pullbacks implies equality of the source functions; if it is not
surjective, changing a source function only at a missed point disproves
injectivity.  If \(\rho_{\mathfrak a}\) is injective, the fiber condition is
vacant and every target function lifts; if it is not injective, a target
function taking the values zero and \(-1\) at two points in the same fiber
does not lift.  These witnesses are finite-valued, so the criteria hold for
both classes.  Finally, when \(\rho_{\mathfrak a}\) is surjective, taking
the maximum over its domain samples every point of its target, which proves
the displayed sup-norm identity.  For the completed classes, the two
fiberwise formulas take values in \([-\infty,0]\); the adjoined values
\(-\infty\) and zero are exactly the supremum and infimum over an empty
fiber.  The first equivalence in
\eqref{eq:p122-curve-psh-finite-map-adjoints} says coordinatewise that
every value of \(w\) over \(x\) is at most \(u(x)\), and the second says
that \(u(x)\) is at most every value of \(w\) over \(x\).  This proves both
adjunctions.  On a nonempty fiber, applying either fiberwise operation to a
pullback recovers the original value; this proves the two identity claims
when the map is surjective.  On a singleton fiber, pulling either fiberwise
operation back recovers \(w\); this proves the identity claims when the map
is injective.  Pointwise on the fiber through \(y\), the two source-side
composites take respectively the infimum and the supremum of the values of
\(w\) on that fiber.  They are therefore the claimed minorant and majorant;
equality for either composite is equivalent to constancy on every fiber,
which is already the pullback-image criterion.
\end{proof}

\begin{theorem}[Generic-norm trace descent for finite curve maps]
\label{thm:p126-curve-generic-norm-trace-descent}
Let
\(\varphi:(R,\mathfrak m)\to(S,\mathfrak n)\) be a finite local
homomorphism between one-dimensional reduced Nagata noetherian local rings,
let \(\mathfrak a\subset R\) be \(\mathfrak m\)-primary, and assume that the
branch map
\[
 \rho_{\mathfrak a}:V_{\mathfrak aS}\longrightarrow V_{\mathfrak a}
\]
is surjective.  Put \(D=R^\nu\).  For a base branch \(\mathfrak q\), write
\[
 K_{\mathfrak q}:=\operatorname{Frac}(D_{\mathfrak q}),\qquad
 T_{\mathfrak q}:=S\otimes_RK_{\mathfrak q},\qquad
 \delta_{\mathfrak q}:=\dim_{K_{\mathfrak q}}T_{\mathfrak q}.
\]
For every source branch \(\mathfrak q'\) over \(\mathfrak q\), let
\(e_{\mathfrak q'/\mathfrak q}\) and
\(f_{\mathfrak q'/\mathfrak q}\) denote its ramification index and residue
degree.  Then:
\begin{enumerate}
\item The algebra \(T_{\mathfrak q}\) is a nonzero product of finite field
      extensions of \(K_{\mathfrak q}\), and
      \[
       \delta_{\mathfrak q}
          =\sum_{\mathfrak q'\mid\mathfrak q}
             e_{\mathfrak q'/\mathfrak q}
             f_{\mathfrak q'/\mathfrak q}.
       \tag{\(\triangleleft\triangleleft\)}
       \label{eq:p126-curve-generic-degree}
      \]
\item On the continuous or completed ideal-psh classes define
      \[
       (\operatorname{Tr}_{\varphi,\mathfrak a}u)
          (v_{\mathfrak q}^{\mathfrak a})
       :=\sum_{\mathfrak q'\mid\mathfrak q}
          e_{\mathfrak q'/\mathfrak q}
          f_{\mathfrak q'/\mathfrak q}\,
          u(v_{\mathfrak q'}^{\mathfrak aS}),
      \]
      and put
      \[
       (\operatorname{Av}_{\varphi,\mathfrak a}^{\mathrm{gen}}u)
          (v_{\mathfrak q}^{\mathfrak a})
       :=\delta_{\mathfrak q}^{-1}
          (\operatorname{Tr}_{\varphi,\mathfrak a}u)
          (v_{\mathfrak q}^{\mathfrak a}).
       \tag{\(\triangleright\triangleright\)}
       \label{eq:p126-curve-generic-average}
      \]
      A positive sum containing \(-\infty\) is interpreted as
      \(-\infty\).  Both maps preserve order and decreasing limits, and
      \[
       \rho_{\mathfrak a,*}u
          \leq\operatorname{Av}_{\varphi,\mathfrak a}^{\mathrm{gen}}u
          \leq\rho_{\mathfrak a,!}u,
       \qquad
       \operatorname{Av}_{\varphi,\mathfrak a}^{\mathrm{gen}}
          \rho_{\mathfrak a}^*=\operatorname{id}.
       \tag{\(\blacktriangleright\blacktriangleright\)}
       \label{eq:p126-curve-generic-average-retraction}
      \]
      Conversely, if the branch map is not surjective, pullback admits no
      left inverse even as a map of sets.  Thus branch surjectivity is the
      exact condition for a finite curve-link pullback to admit a
      pushforward retraction.  Extending the normalized average by the same
      formula to finite real-valued obstacles \(\psi\), one has
      \[
       \operatorname{Av}^{\mathrm{gen}}_{\varphi,\mathfrak a}
          (P^{\mathrm{id}}_{\mathfrak aS}\psi)
       \leq
       P^{\mathrm{id}}_{\mathfrak a}
          (\operatorname{Av}^{\mathrm{gen}}_{\varphi,\mathfrak a}\psi).
       \tag{\(\blacktriangle\blacktriangle\)}
       \label{eq:p126-curve-generic-average-envelope}
      \]
      Equality holds exactly when no fiber contains two points \(y,z\) with
      \(\psi(y)<0<\psi(z)\).
\item On finite real-valued functions the normalized average is a positive
      linear contraction for the sup norm.  Hence
      \[
       \Pi_{\varphi,\mathfrak a}
          :=\rho_{\mathfrak a}^*
             \operatorname{Av}^{\mathrm{gen}}_{\varphi,\mathfrak a}
       \tag{\(\circledast\)}
       \label{eq:p126-curve-generic-barycentric-projection}
      \]
      is a positive norm-one idempotent projection onto the fiber-constant
      functions, which are exactly the pullback image.  Write
      \(\mathsf{Meas}_{\mathbf R}(V)\) for the finite signed measures on a
      finite set \(V\).  Dually, for
      \(\mu\in\mathsf{Meas}_{\mathbf R}(V_{\mathfrak a})\), define
      \[
       (\operatorname{Lift}_{\varphi,\mathfrak a}\mu)
          (\{v_{\mathfrak q'}^{\mathfrak aS}\})
       :=\frac{e_{\mathfrak q'/\mathfrak q}
                    f_{\mathfrak q'/\mathfrak q}}
               {\delta_{\mathfrak q}}\,
          \mu(\{v_{\mathfrak q}^{\mathfrak a}\}).
       \tag{\(\circledcirc\)}
       \label{eq:p126-curve-generic-measure-lift}
      \]
      The lift is positive and linear, is an isometry for the total-variation
      norm, and satisfies
      \((\rho_{\mathfrak a})_\#
          \operatorname{Lift}_{\varphi,\mathfrak a}\mu=\mu\).  It obeys
      the exact pairing identity
      \[
       \sum_{v\in V_{\mathfrak a}}
          (\operatorname{Av}^{\mathrm{gen}}_{\varphi,\mathfrak a}u)(v)
          \mu(\{v\})
       =\sum_{w\in V_{\mathfrak aS}}
          u(w)(\operatorname{Lift}_{\varphi,\mathfrak a}\mu)(\{w\}).
       \tag{\(\bullet\)}
       \label{eq:p126-curve-generic-function-measure-pairing}
      \]
      for every finite real-valued \(u\).  In particular, it sends
      probability measures to probability measures.  Moreover,
      \[
       Q_{\varphi,\mathfrak a}
         :=\operatorname{Lift}_{\varphi,\mathfrak a}
              (\rho_{\mathfrak a})_\#
       \tag{\(\circledcirc\circledast\)}
       \label{eq:p126-curve-generic-measure-projection}
      \]
      is a positive norm-one idempotent projection on
      \(\mathsf{Meas}_{\mathbf R}(V_{\mathfrak aS})\).  Its image is exactly
      the signed
      measures \(\nu\) whose restriction to every fiber has the canonical
      conditional weights:
      \[
       \nu(\{v_{\mathfrak q'}^{\mathfrak aS}\})
        =\frac{e_{\mathfrak q'/\mathfrak q}
                    f_{\mathfrak q'/\mathfrak q}}
                {\delta_{\mathfrak q}}
          ((\rho_{\mathfrak a})_\#\nu)
             (\{v_{\mathfrak q}^{\mathfrak a}\}).
       \tag{\(\bullet\circledast\)}
       \label{eq:p126-curve-generic-measure-projection-image}
      \]
      Under the pairing above, \(Q_{\varphi,\mathfrak a}\) is the Banach
      adjoint of \(\Pi_{\varphi,\mathfrak a}\).  More precisely, let
      \(\mu\) be a full-support probability measure on
      \(V_{\mathfrak a}\), and set
      \(\nu=\operatorname{Lift}_{\varphi,\mathfrak a}\mu\).  For every
      real-valued base function \(h\) and source function \(u\),
      \[
       \operatorname{Av}^{\mathrm{gen}}_{\varphi,\mathfrak a}
          ((\rho_{\mathfrak a}^*h)u)
        =h\operatorname{Av}^{\mathrm{gen}}_{\varphi,\mathfrak a}u,
       \qquad
       \Pi_{\varphi,\mathfrak a}((\rho_{\mathfrak a}^*h)u)
        =(\rho_{\mathfrak a}^*h)\Pi_{\varphi,\mathfrak a}u.
       \tag{\(\mathrm{CE}\)}
       \label{eq:p126-curve-generic-module-identity}
      \]
      Thus \(\Pi_{\varphi,\mathfrak a}\) is the conditional expectation
      onto the fiber-constant functions for the finite probability space
      \((V_{\mathfrak aS},\nu)\).  For every \(1\leq p\leq\infty\),
      \[
       \|\Pi_{\varphi,\mathfrak a}u\|_{L^p(\nu)}
          \leq \|u\|_{L^p(\nu)}.
       \tag{\(L^p\)}
       \label{eq:p126-curve-generic-lp-contraction}
      \]
      If \(1<p<\infty\), equality holds exactly when \(u\) is constant
      on every fiber.  For \(p=2\), the projection is the orthogonal
      projection onto \(\rho_{\mathfrak a}^*L^2(\mu)\), and one has the
      exact variance decomposition
      \[
       \begin{split}
       \|u\|_{L^2(\nu)}^2
        &=\|\Pi_{\varphi,\mathfrak a}u\|_{L^2(\nu)}^2
          +\|u-\Pi_{\varphi,\mathfrak a}u\|_{L^2(\nu)}^2,\\
       \|u-\Pi_{\varphi,\mathfrak a}u\|_{L^2(\nu)}^2
        &=\sum_{\mathfrak q}\mu(\{v_{\mathfrak q}^{\mathfrak a}\})
          \sum_{\mathfrak q'\mid\mathfrak q}
          \frac{e_{\mathfrak q'/\mathfrak q}
                    f_{\mathfrak q'/\mathfrak q}}
               {\delta_{\mathfrak q}}
          \left|u(v_{\mathfrak q'}^{\mathfrak aS})
             -(\operatorname{Av}^{\mathrm{gen}}_{\varphi,\mathfrak a}u)
                (v_{\mathfrak q}^{\mathfrak a})\right|^2.
       \end{split}
       \tag{\(\mathrm{Var}\)}
       \label{eq:p126-curve-generic-variance}
      \]
\item The trace and average commute with the change-of-normalizing-ideal
      transports of
      \eqref{eq:p122-curve-psh-normalization-transport}.  In a tower of
      branch-surjective finite local maps, the unnormalized traces compose
      exactly.  As operators on all finite real-valued functions, the
      normalized averages compose if and only if the upper generic degree is
      constant on the intermediate branches above each fixed base branch.
      Dually, the signed-measure lifts compose under exactly the same
      condition.  In particular, both compose when the upper map is finite
      flat.
\item Let \(g\in\mathfrak n\) be an \(S\)-nonzerodivisor, and set
      \[
       \beta_{\mathfrak q}(g)
          :=\sum_{\mathfrak q'\mid\mathfrak q}
             f_{\mathfrak q'/\mathfrak q}\,
             w_{\mathfrak q'}(g).
       \tag{\(\blacktriangleleft\blacktriangleleft\)}
       \label{eq:p126-curve-generic-norm-order}
      \]
      This is the order of the generic determinant norm of \(g\) over the
      normalized base DVR.  For every sufficiently large integer \(M\), and
      every sufficiently large common multiple \(L\) of the positive
      integers \(\delta_{\mathfrak q}\), there are
      \(R\)-nonzerodivisors \(h_M,\widetilde h_L\in\mathfrak m\) with
      \[
       w_{\mathfrak q}(h_M)=M\beta_{\mathfrak q}(g),
       \qquad
       w_{\mathfrak q}(\widetilde h_L)
          =\frac{L}{\delta_{\mathfrak q}}\beta_{\mathfrak q}(g),
      \]
      and, for every \(N\geq1\),
      \[
       \operatorname{Tr}_{\varphi,\mathfrak a}
          u^{\mathfrak aS}_{g,N}
          =u^{\mathfrak a}_{h_M,MN},
       \qquad
       \operatorname{Av}^{\mathrm{gen}}_{\varphi,\mathfrak a}
          u^{\mathfrak aS}_{g,N}
          =u^{\mathfrak a}_{\widetilde h_L,LN}.
       \tag{\(\blacklozenge\blacklozenge\)}
       \label{eq:p126-curve-generic-principal-descent}
      \]
\end{enumerate}
Thus every branch-surjective finite curve map has a canonical generic-degree
weighted link pushforward, and its principal source models descend exactly
after a conductor-controlled common scaling.  The next theorem identifies
the finite-flat range in which the generic norm is already one canonical
downstairs determinant norm, with no conductor scaling.  The signed-measure
lift and conditional expectation above are only finite-link operations;
they are not a construction or functoriality assertion for a singular
Monge--Amp\`ere measure.
\end{theorem}

\begin{proof}
Fix \(\mathfrak q\), let \(\mathfrak p\) be its corresponding minimal prime
of \(R\), and put \(A=D_{\mathfrak q}\).  Since
\(K_{\mathfrak q}\simeq R_{\mathfrak p}\), the algebra
\(T_{\mathfrak q}=S_{R\setminus\mathfrak p}\) is reduced and finite over
\(K_{\mathfrak q}\).  Branch surjectivity makes it nonzero, so it is a
product of finite field extensions.  The DVR \(A\) is Nagata.  Hence its
integral closure \(B\) in \(T_{\mathfrak q}\) is finite and torsion-free,
and therefore free of rank \(\delta_{\mathfrak q}\).  Its maximal ideals
are precisely the source branches above \(\mathfrak q\).  If \(\pi\) is a
uniformizer of \(A\), filtering the local factors of \(B/\pi B\) gives
\[
 \delta_{\mathfrak q}=\operatorname{length}_A(B/\pi B)
   =\sum_{\mathfrak q'\mid\mathfrak q}
      e_{\mathfrak q'/\mathfrak q}
      f_{\mathfrak q'/\mathfrak q},
\]
which proves \eqref{eq:p126-curve-generic-degree}.

The coefficients of the normalized average are positive and sum to one.
The two inequalities in
\eqref{eq:p126-curve-generic-average-retraction} are therefore the elementary
minimum--average--maximum inequalities on each fiber.  By
\eqref{eq:p126-ramification-cancellation}, a pullback is constant on every
fiber, so the degree formula gives the retraction identity.  The remaining
order and limit assertions are coordinatewise on a finite set.  If the
branch map is not surjective, \(\rho_{\mathfrak a}^*\) is not injective by
\Cref{thm:p122-curve-ideal-psh-functoriality}; a map with a left inverse is
injective, proving necessity.  The function \(s\mapsto\min\{s,0\}\) is
concave.  Its finite weighted Jensen inequality on each branch fiber gives
\eqref{eq:p126-curve-generic-average-envelope}.  Since every weight is
strictly positive, equality in this piecewise-linear Jensen inequality fails
precisely when the fiber contains both a strictly negative and a strictly
positive value, which proves the stated criterion.

For finite real-valued functions, a positive weighted average has sup norm
at most the source sup norm.  Together with
\(\operatorname{Av}^{\mathrm{gen}}\rho^*=\operatorname{id}\), this makes
\(\Pi_{\varphi,\mathfrak a}\) idempotent with the fiber-constant image; its
norm is one because it fixes constants.  Formula
\eqref{eq:p126-curve-generic-degree} says that the masses assigned in
\eqref{eq:p126-curve-generic-measure-lift} sum over each fiber to the original
base mass.  This proves both the probability and pushforward identities.
Expanding the two finite sums proves
\eqref{eq:p126-curve-generic-function-measure-pairing}.
All lifted atoms over one base point have the sign of that base mass, so
summing absolute values fiber by fiber gives
\[
 \|\operatorname{Lift}_{\varphi,\mathfrak a}\mu\|_{\mathrm{TV}}
   =\|\mu\|_{\mathrm{TV}}.
\]
Together with the evident positivity and linearity, this proves the stated
total-variation isometry on signed measures.  Since
\((\rho_{\mathfrak a})_\#\operatorname{Lift}_{\varphi,\mathfrak a}
=\operatorname{id}\), the operator \(Q_{\varphi,\mathfrak a}\) is
idempotent and has precisely the image described in
\eqref{eq:p126-curve-generic-measure-projection-image}.  Measure pushforward
is a total-variation contraction, while the lift is an isometry; hence
\(\|Q_{\varphi,\mathfrak a}\|\leq1\), and equality follows because \(Q\)
fixes its nonzero image.  The pairing identity applied to
\((\rho_{\mathfrak a})_\#\nu\) says exactly that \(Q\) is the adjoint of
\(\Pi\).

The first identity in
\eqref{eq:p126-curve-generic-module-identity} follows by taking the base
factor \(h\) outside each weighted fiber sum, and the second follows after
pullback.  With \(\nu=\operatorname{Lift}\mu\), Jensen's inequality for
\(t\mapsto |t|^p\), applied on every fiber and then integrated against
\(\mu\), proves \eqref{eq:p126-curve-generic-lp-contraction} for finite
\(p\); the case \(p=\infty\) was proved above.  Since all fiber weights and
all base masses are positive, strict convexity gives the stated equality
criterion when \(1<p<\infty\).  If \(z=\rho_{\mathfrak a}^*h\) is
fiber-constant, the pairing identity and the retraction formula give
\[
 \langle u-\Pi_{\varphi,\mathfrak a}u,z\rangle_{L^2(\nu)}
 =\left\langle
    \operatorname{Av}^{\mathrm{gen}}_{\varphi,\mathfrak a}
       (u-\Pi_{\varphi,\mathfrak a}u),h
  \right\rangle_{L^2(\mu)}=0.
\]
Thus \(\Pi_{\varphi,\mathfrak a}\) is the asserted orthogonal projection.
Pythagoras gives the first line of
\eqref{eq:p126-curve-generic-variance}, and expanding the lifted measure
fiber by fiber gives its second line.

Changing the normalizing ideal multiplies all source coordinates over one
base branch by the same positive factor, which proves transport
compatibility.  In a composable tower, ramification indices multiply by
\cite[Tag~0BRL]{StacksProject}, and residue degrees multiply by the field
tower law.  Expanding the two finite sums proves exact composition of the
unnormalized traces.  If the upper degree has the common value \(c_{\mathfrak
q}\) over the intermediate branches above \(\mathfrak q\), the composite
degree is \(c_{\mathfrak q}\delta_{\mathfrak q}\); division proves the
stated composition law for normalized averages.
Conversely, let \(D_{\mathfrak q}\) be the composite generic degree and let
\(c_{\mathfrak q'}\) be the upper generic degree at an intermediate branch.
Compare the coefficient of a top branch over \(\mathfrak q'\) in the direct
and iterated normalized averages.  After cancelling the positive product of
ramification and residue degrees, equality of the two operators forces
\[
                   \delta_{\mathfrak q}c_{\mathfrak q'}=D_{\mathfrak q}
\]
for every intermediate branch over \(\mathfrak q\).  The right side is
fixed, so the upper generic degree is constant over that intermediate fiber.
The lift statement follows either from the same coefficient comparison or from
the exact function--measure duality.

Finally, \(g\) is nonzero in every field factor of
\(T_{\mathfrak q}\).  Multiplication by \(g\) on \(B\) is injective, and
the DVR length--determinant formula
\cite[\emph{Lengths and determinants}, Tag~02QA]{stacks-project} gives
\[
 \operatorname{ord}_A\det(m_g\mid B)
   =\operatorname{length}_A(B/gB)
   =\sum_{\mathfrak q'\mid\mathfrak q}
      f_{\mathfrak q'/\mathfrak q}w_{\mathfrak q'}(g)
   =\beta_{\mathfrak q}(g).
\]
Every \(\beta_{\mathfrak q}(g)\) is a positive integer.  Let
\(C_{\mathfrak q}\) be the conductor thresholds of
\eqref{eq:p126-curve-conductor-orders}.  For all sufficiently large \(M\),
the integers \(M\beta_{\mathfrak q}(g)\) exceed those thresholds.  For all
sufficiently large common multiples \(L\) of the
\(\delta_{\mathfrak q}\), the same holds for
\(L\beta_{\mathfrak q}(g)/\delta_{\mathfrak q}\).  Applying
\Cref{prop:p126-curve-conductor-principal-realization} to these two vectors
produces \(h_M\) and \(\widetilde h_L\).  Since
\(d_{\mathfrak q'}(\mathfrak aS)
 =e_{\mathfrak q'/\mathfrak q}d_{\mathfrak q}(\mathfrak a)\), direct
substitution gives
\[
 (\operatorname{Tr}_{\varphi,\mathfrak a}
       u^{\mathfrak aS}_{g,N})(v_{\mathfrak q}^{\mathfrak a})
   =-\frac{\beta_{\mathfrak q}(g)}
           {N d_{\mathfrak q}(\mathfrak a)},
\]
and division by \(\delta_{\mathfrak q}\) gives the average.  The prescribed
orders of \(h_M\) and \(\widetilde h_L\) now yield exactly
\eqref{eq:p126-curve-generic-principal-descent}.
\end{proof}

\begin{example}[Variable generic degree obstructs normalized composition]
\label{ex:p126-curve-variable-degree-average}
Let \(K\) be a field and consider the finite local maps
\[
 \begin{aligned}
 R&=K[[t]]\longrightarrow
 S=K[[x,y]]/(xy),& t&\longmapsto x+y,\\
 S&\longrightarrow
 U=K[[a,b,c]]/(ab,ac,bc),& x&\longmapsto a,\quad
 y\longmapsto b+c.
 \end{aligned}
\]
Indeed, \(x\) satisfies \(x^2-tx=0\), while \(b\) satisfies
\(b^2-yb=0\) and \(c=y-b\).  Thus the first map has two generic branches
of degree one.  Over its \(x\)-branch the second map has one generic branch,
whereas over its \(y\)-branch it has two, again all with ramification and
residue degree one.

Take \(\mathfrak a=(t)\), and give the \(a,b,c\) branches of the completed
link of \(U\) the values \(-1,-2,-3\), respectively.  The direct normalized
average from \(U\) to \(R\) is
\[
                          \frac{-1-2-3}{3}=-2.
\]
Averaging first to \(S\) gives the pair \((-1,-5/2)\), whose subsequent
average to \(R\) is
\[
                          \frac{-1-5/2}{2}=-\frac74.
\]
Thus normalized generic averages need not compose when the upper generic
degree varies between intermediate branches.  The unnormalized traces do
compose: both routes give \(-6\), exactly as asserted in
\Cref{thm:p126-curve-generic-norm-trace-descent}.  The dual obstruction is
equally explicit.  A unit mass on the unique branch of \(R\) lifts directly
to masses
\[
                         (1/3,1/3,1/3)
\]
on the \(a,b,c\) branches, whereas successive lifts through \(S\) give
\[
                         (1/2,1/4,1/4).
\]
Thus the signed-measure lifts also fail to compose, with the same sharp
constant-upper-degree boundary.
\end{example}

\begin{theorem}[Finite-flat norm descent on reduced curve links]
\label{thm:p126-curve-finite-flat-norm-descent}
Let
\(\varphi:(R,\mathfrak m)\to(S,\mathfrak n)\) be a finite flat local
homomorphism of constant rank \(d\geq1\) between one-dimensional reduced
Nagata noetherian local rings, and let \(\mathfrak a\subset R\) be
\(\mathfrak m\)-primary.  Put \(D=R^\nu\) and \(E=S^\nu\), and write
\(v_{\mathfrak q}^{\mathfrak a}
=w_{\mathfrak q}/d_{\mathfrak q}(\mathfrak a)\) and
\(v_{\mathfrak q'}^{\mathfrak aS}
=w_{\mathfrak q'}/d_{\mathfrak q'}(\mathfrak aS)\).
For a branch \(\mathfrak q'\) of \(S\) over a branch
\(\mathfrak q\) of \(R\), let
\[
 e_{\mathfrak q'/\mathfrak q}\geq1,
 \qquad
 f_{\mathfrak q'/\mathfrak q}
   :=[\kappa(\mathfrak q'): \kappa(\mathfrak q)]
\]
be its ramification index and residue degree.  Then:
\begin{enumerate}
\item The branch map \(\rho_{\mathfrak a}\) is surjective, the generic
      degrees of
      \Cref{thm:p126-curve-generic-norm-trace-descent} all equal \(d\), and
      over every \(\mathfrak q\) one has
      \[
       \sum_{\mathfrak q'\mid\mathfrak q}
          e_{\mathfrak q'/\mathfrak q}
          f_{\mathfrak q'/\mathfrak q}=d.
       \tag{\(\triangleleft\)}\label{eq:p126-curve-ef-degree}
      \]
\item For every nonzerodivisor \(g\in S\), its finite-flat determinant
      norm satisfies
      \[
       w_{\mathfrak q}\!\left(\mathrm N_{S/R}(g)\right)
          =\sum_{\mathfrak q'\mid\mathfrak q}
             f_{\mathfrak q'/\mathfrak q}\,
             w_{\mathfrak q'}(g).
       \tag{\(\blacktriangle\)}\label{eq:p126-curve-norm-order}
      \]
      If \(g\in\mathfrak n\), then \(\mathrm N_{S/R}(g)\) is an
      \(R\)-nonzerodivisor in \(\mathfrak m\), so its principal ideal is
      \(\mathfrak m\)-primary.
\item On the continuous or completed ideal-psh classes define the
      finite-flat norm average by
      \[
       (\operatorname{Av}_{\varphi,\mathfrak a}u)
          (v_{\mathfrak q}^{\mathfrak a})
       :=\frac1d\sum_{\mathfrak q'\mid\mathfrak q}
          e_{\mathfrak q'/\mathfrak q}
          f_{\mathfrak q'/\mathfrak q}\,
          u(v_{\mathfrak q'}^{\mathfrak aS}),
      \tag{\(\triangleright\)}\label{eq:p126-curve-norm-average}
      \]
      This is exactly the generic-degree average
      \eqref{eq:p126-curve-generic-average} in the finite-flat case; a
      positive-weighted sum containing \(-\infty\) is interpreted as
      \(-\infty\).  This is order preserving, preserves decreasing limits,
      and satisfies
      \[
       \rho_{\mathfrak a,*}u
          \leq\operatorname{Av}_{\varphi,\mathfrak a}u
          \leq\rho_{\mathfrak a,!}u,
       \qquad
       \operatorname{Av}_{\varphi,\mathfrak a}
          \rho_{\mathfrak a}^*=\operatorname{id}.
       \tag{\(\blacktriangleright\)}\label{eq:p126-curve-average-sandwich}
      \]
      It commutes with the change-of-normalizing-ideal transports
      \eqref{eq:p122-curve-psh-normalization-transport}, and the norm
      averages compose in towers of finite flat local maps.
\item Norm averaging sends principal models to principal models exactly:
      for every \(S\)-nonzerodivisor \(g\in\mathfrak n\) and \(N\geq1\),
      \[
       \operatorname{Av}_{\varphi,\mathfrak a}
          u^{\mathfrak aS}_{g,N}
        =u^{\mathfrak a}_{\mathrm N_{S/R}(g),dN}.
       \tag{\(\blacklozenge\)}\label{eq:p126-curve-principal-norm-descent}
      \]
\end{enumerate}
Consequently, the source conductor realization of every sufficiently deep
integral branch-order vector descends by one determinant norm to the
weighted-average principal profile downstairs.  This supplies a genuine
pushforward and conductor-compatible range of P12.6 for finite flat curve
maps.  It asserts no norm pushforward for a nonflat finite map.
\end{theorem}

\begin{proof}
Because \(R\) is local and \(S\) is finite flat of constant rank \(d\),
the \(R\)-module \(S\) is free of rank \(d\), and
\(\mathrm N_{S/R}(g)=\det(m_g)\) is the determinant norm of
\eqref{eq:appA-element-norm}.

Fix \(\mathfrak q\), let \(A=D_{\mathfrak q}\) be the corresponding
normalized discrete valuation ring of \(R\), and write \(K\) for its
fraction field.  If \(\mathfrak p\) is the minimal prime of \(R\)
corresponding to \(\mathfrak q\), then \(K\simeq R_{\mathfrak p}\).
Consequently the generic algebra
\[
                          T:=S\otimes_RK
\]
is the localization \(S_{R\setminus\mathfrak p}\).  It is therefore a
reduced \(K\)-algebra of dimension \(d\), hence a product of finite field
extensions of \(K\).  The ring \(A\) is Nagata, being a localization of the
finite normalization of the Nagata ring \(R\).  Thus the integral closure
\(B\) of \(A\) in \(T\) is finite; it is torsion-free, hence a finite free
\(A\)-module of rank \(d\).  Its maximal ideals are exactly the branches
\(\mathfrak q'\) above \(\mathfrak q\), and
\(B_{\mathfrak q'}\) is their normalized discrete valuation ring.

Let \(\pi\) be a uniformizer of \(A\).  The quotient \(B/\pi B\) has
\(A\)-length \(d\).  Localizing it at the maximal ideals of \(B\), and
filtering each
\(B_{\mathfrak q'}/\pi B_{\mathfrak q'}
 =B_{\mathfrak q'}/\mathfrak q'^{e_{\mathfrak q'/\mathfrak q}}\),
gives
\[
 d=\operatorname{length}_A(B/\pi B)
   =\sum_{\mathfrak q'\mid\mathfrak q}
      e_{\mathfrak q'/\mathfrak q}
      [\kappa(\mathfrak q'): \kappa(\mathfrak q)].
\]
This proves \eqref{eq:p126-curve-ef-degree}, and in particular proves that
every base branch has a branch above it.

If \(g\) is a nonzerodivisor of \(S\), it is nonzero in every field factor
of \(T\), so multiplication by \(g\) is injective on \(B\).  Determinants
commute with base change by \Cref{prop:appA-norm-identities}.  After
extending scalars to \(K\), multiplication on \(S\) and on the
\(A\)-lattice \(B\) becomes the same endomorphism of \(T\); its determinant
is independent of the chosen \(K\)-basis.  Hence the determinant of
\(m_g\) on \(B\) is the image of \(\mathrm N_{S/R}(g)\) in \(A\).  The DVR length--determinant formula
\cite[\emph{Lengths and determinants}, Tag~02QA]{stacks-project} and the
same residue-field filtration give
\[
 \begin{aligned}
 w_{\mathfrak q}\!\left(\mathrm N_{S/R}(g)\right)
   &=\operatorname{length}_A(B/gB)\\
   &=\sum_{\mathfrak q'\mid\mathfrak q}
       f_{\mathfrak q'/\mathfrak q}\,
       \operatorname{length}_{B_{\mathfrak q'}}
          (B_{\mathfrak q'}/gB_{\mathfrak q'})\\
   &=\sum_{\mathfrak q'\mid\mathfrak q}
       f_{\mathfrak q'/\mathfrak q}\,
       w_{\mathfrak q'}(g),
 \end{aligned}
\]
which is \eqref{eq:p126-curve-norm-order}.  It is finite on every branch,
so the norm avoids every minimal prime of the reduced ring \(R\) and is a
nonzerodivisor.  If \(g\in\mathfrak n\), its image in the artinian local
algebra \(S/\mathfrak mS\) is nilpotent.  Multiplication by that image has
zero determinant, so \(\mathrm N_{S/R}(g)\in\mathfrak m\).  The resulting
principal ideal is \(\mathfrak m\)-primary by the one-dimensional argument
in \Cref{prop:p126-curve-conductor-principal-realization}.

Formula \eqref{eq:p126-curve-ef-degree} says that the coefficients in
\eqref{eq:p126-curve-norm-average} are positive and sum to one on each
fiber.  Thus the average lies between the fiberwise infimum and supremum,
which proves the first part of
\eqref{eq:p126-curve-average-sandwich}; it plainly preserves finite values,
order, and decreasing limits.  The ramification-cancellation formula
\eqref{eq:p126-ramification-cancellation} makes a pullback constant on every
fiber, so averaging and using \eqref{eq:p126-curve-ef-degree} proves the
retraction identity.  A change of normalizing ideal multiplies all source
coordinates over a fixed base branch by the same positive factor, proving
compatibility with \eqref{eq:p122-curve-psh-normalization-transport}.
Ramification indices and residue degrees multiply in a tower; expanding the
two successive finite sums therefore proves composition of the averages.

Finally,
\[
 \begin{aligned}
 (\operatorname{Av}_{\varphi,\mathfrak a}
       u^{\mathfrak aS}_{g,N})(v_{\mathfrak q}^{\mathfrak a})
  &=-\frac1d\sum_{\mathfrak q'\mid\mathfrak q}
       e_{\mathfrak q'/\mathfrak q}f_{\mathfrak q'/\mathfrak q}
       \frac{w_{\mathfrak q'}(g)}
            {N e_{\mathfrak q'/\mathfrak q}
               d_{\mathfrak q}(\mathfrak a)}\\
  &=-\frac{w_{\mathfrak q}(\mathrm N_{S/R}(g))}
            {dN d_{\mathfrak q}(\mathfrak a)}
   =u^{\mathfrak a}_{\mathrm N_{S/R}(g),dN}
       (v_{\mathfrak q}^{\mathfrak a}),
 \end{aligned}
\]
where the middle equality is \eqref{eq:p126-curve-norm-order}.  This proves
\eqref{eq:p126-curve-principal-norm-descent}; combining it with the source
version of \Cref{prop:p126-curve-conductor-principal-realization} gives the
last assertion.
\end{proof}

\begin{example}[The two endpoint weights in finite-flat norm descent]
\label{ex:p126-curve-norm-endpoint-weights}
Let \(K\) be a field.  For the totally ramified map
\[
 K[[t]]\longrightarrow K[[u]],\qquad t\longmapsto u^e,
\]
the rank is \(d=e\), the residue degree is one, and
\(\mathrm N(u)=\pm t\).  With \(\mathfrak a=(t)\), the normalized source
value of \(u^{(u^e)}_{u,N}\) is \(-1/(eN)\); the weight \(ef/d=1\) leaves
it unchanged, and the descended model
\(u^{(t)}_{\pm t,eN}\) has the same value.  Thus the ramification index
belongs in the averaging weight but not in
\eqref{eq:p126-curve-norm-order}.

At the other endpoint, if \(K'/K\) is a finite extension of degree \(f\),
then
\[
 K[[t]]\longrightarrow K'[[t]]
\]
has \(e=1\), \(d=f\), and \(\mathrm N(t)=t^f\).  The source model
\(u^{(t)}_{t,N}\) and the descended model
\(u^{(t)}_{t^f,fN}\) both have value \(-1/N\).  This time the residue
degree appears both in the norm order and in the rank denominator and
cancels there.
\end{example}

\begin{example}[A boundary profile on the two-branch node]
\label{ex:p122-curve-ideal-psh-boundary}
For \(R=\mathbf C[[x,y]]/(xy)\) and
\(\mathfrak a=(x,y)\), label the two points of
\(V_{\mathfrak a}\) by the \(x\)- and \(y\)-axes.  The profile
\[
                             u=(-1,0)
\]
belongs to \(\operatorname{CPSH}^{\mathrm{id}}_{\mathfrak a}(R)\) and
is an increasing uniform limit of principal-ideal profiles.  It is not a
decreasing limit of principal-ideal profiles: every such profile is
strictly negative on both branches.  Thus uniform ideal closure and
decreasing model regularization genuinely differ at a reducible boundary,
even though the completed psh envelope is the coordinatewise formula
\eqref{eq:p123-curve-ideal-psh-envelope}.
\end{example}

\begin{example}[A nonsurjective branch pullback reaches the bottom]
\label{ex:p122-curve-psh-pullback-bottom}
Let
\[
 R=\mathbf C[[x,y]]/(xy),\qquad S=\mathbf C[[x]],
 \qquad
 \varphi:R\longrightarrow S,\qquad x\longmapsto x,\quad y\longmapsto0,
\]
and take \(\mathfrak a=(x,y)\), so that
\(\mathfrak aS=(x)\).  This is a finite local homomorphism of reduced
excellent curve rings.  Its normalized-link map sends the unique branch of
\(S\) to the \(x\)-axis branch of \(R\), and misses the \(y\)-axis branch.
On the two-point link of \(R\), ordered as the \(x\)- and \(y\)-axes, the
function
\[
                              u=(-\infty,0)
\]
belongs to \(\operatorname{PSH}^{\mathrm{id},\bot}_{\mathfrak a}(R)\)
but is not its bottom element.  Nevertheless,
\[
                 \rho_{\mathfrak a}^*u\equiv-\infty
                 \quad\text{on }V_{(x)}(S).
\]
Thus finite pullback preserves the bottom-inclusive completion but cannot in
general preserve the class obtained by deleting its bottom element.  This is
the precise reason for retaining \(-\infty\) in the functorial statement.
Here \(\rho_{\mathfrak a}\) is injective but not surjective, so pullback is
surjective but not injective, exactly as the preceding theorem predicts.
\end{example}

The word ``psh'' in the preceding notation cannot be justified by extending
every homogeneous branch profile radially across an analytic singularity.
On any isolated curve singularity whose normalization is a finite union of
discs and whose branches are separated by the conductor, however, the exact
branch-compatibility condition can be computed.

\begin{theorem}[Conductorial radial classification on a disc-normalized curve]
\label{thm:p122-conductorial-disc-radial-psh-classification}
Let \(K\) be a complete nontrivially valued non-Archimedean field, choose
\(\varpi\in K\) with \(0<|\varpi|<1\), and put
\(r=|\varpi|\) and \(a=-\log r\).  Let \(A\) be a reduced strict
\(K\)-affinoid algebra of pure dimension one whose normalization is finite
and has the form
\[
 A\lhook\joinrel\longrightarrow
 B=\prod_{i=1}^q K\langle t_i\rangle,
 \qquad
 \nu:\widetilde U:=\mathcal M(B)\longrightarrow U:=\mathcal M(A).
\]
Suppose that there is a \(K\)-rational point \(o\in U\) such that
\(\nu^{-1}(o)=\{p_1,\ldots,p_q\}\), where \(p_i\) is the origin of the
\(i\)-th disc \(D_i:=\mathcal M(K\langle t_i\rangle)\), and that \(\nu\)
is an isomorphism away from these points.  Assume that the conductor
contains a separated branch ideal
\[
 \mathfrak d_{\boldsymbol e}
   :=\prod_{i=1}^q t_i^{e_i}K\langle t_i\rangle
   \ \subseteq\
 \mathfrak c_{B/A}:=\{b\in B:bB\subseteq A\}
 \tag{\(\blacktriangledown\)}
 \label{eq:p122-disc-conductor-branch-ideal}
\]
for some integers \(e_i\geq1\).  This conductor hypothesis is automatic if
the finite \(A\)-module \(B/A\) is supported at the maximal ideal of \(o\).

For \(c=(c_1,\ldots,c_q)\in[-\infty,0]^q\), define
\(h_c^\circ:U\setminus\{o\}\to[-\infty,0]\) through the normalization by
\[
 h_c^\circ(\nu(z)):=
 \begin{cases}
  -\infty,&c_i=-\infty,\\[1mm]
  -\dfrac{c_i}{a}\log|t_i(z)|,&c_i> -\infty,
 \end{cases}
 \qquad z\in D_i\setminus\{p_i\}.
 \tag{\(\blacktriangleleft\)}
 \label{eq:p122-disc-normalized-radial-function}
\]
Then the following are equivalent:
\begin{enumerate}
\item \(h_c^\circ\) extends to a classically psh function on \(U\);
\item \(h_c^\circ\) extends to a CLD-v2 psh function on \(U\);
\item either \(c=0\), or \(c_i<0\) for every \(i\), with
      \(-\infty<0\) allowed.
\end{enumerate}
The extension is unique.  Its value at \(o\) is zero for \(c=0\), and is
\(-\infty\) otherwise.  Thus the exact radial-extendable cone is
\[
                         \{0\}\cup[-\infty,0)^q.
 \tag{\(\blacktriangleright\)}
 \label{eq:p122-disc-normalized-radial-cone}
\]
Every finite strictly negative rational profile is realized by one scaled
logarithm \((aN)^{-1}\log|f|\) with \(f\in A\), and all-negative extended
profiles are obtained by decreasing limits after putting zero in the
normalization components corresponding to their \(-\infty\) coordinates.
\end{theorem}

\begin{proof}
First note why the final sentence of the conductor hypothesis holds.  Let
\(\mathfrak m\subset A\) be the maximal ideal of \(o\).  If
\(\operatorname{Supp}_A(B/A)\subseteq\{\mathfrak m\}\), then noetherianity
gives an integer \(M\geq1\) with
\(\mathfrak m^M(B/A)=0\), hence
\[
                            \mathfrak m^M B\subseteq A.
\]
Writing \(I_i\) for the image of \(\mathfrak mB\) in
\(K\langle t_i\rangle\), the fiber hypothesis and the affinoid
Nullstellensatz give \(\sqrt{I_i}=(t_i)\).  Choose \(d_i\geq1\) with
\(t_i^{d_i}\in I_i\), and set \(e_i=Md_i\).  Then
\[
 \prod_i t_i^{e_i}K\langle t_i\rangle
       \subseteq (\mathfrak mB)^M=\mathfrak m^M B\subseteq A.
\]
The left-hand side is a \(B\)-ideal, so it is contained in
\(\mathfrak c_{B/A}\), proving
\eqref{eq:p122-disc-conductor-branch-ideal}.

For the psh classification, the zero profile has the constant zero
extension.  Suppose next that every \(c_i<0\), and put
\[
 I:=\{i:c_i> -\infty\},\qquad J:=\{i:c_i=-\infty\}.
\]
If \(I=\varnothing\), use the constant \(-\infty\) extension.  Otherwise,
put \(\lambda_i=-c_i\in(0,+\infty)\) for \(i\in I\).  Choose positive
rationals
\[
                  \lambda_i^{(n)}=k_{n,i}/N_n\nearrow\lambda_i
                  \qquad(i\in I)
\]
with one denominator \(N_n\) at each stage and with
\(k_{n,i}\geq e_i\).  The latter condition costs nothing: after choosing
common-denominator rational approximants, multiply all their numerators and
their common denominator by one sufficiently large integer.  Define
\(b_n=(b_{n,1},\ldots,b_{n,q})\in B\) by
\[
 b_{n,i}:=
 \begin{cases}
   t_i^{k_{n,i}},&i\in I,\\
   0,&i\in J.
 \end{cases}
\]
By \eqref{eq:p122-disc-conductor-branch-ideal}, \(b_n\in A\); denote the
corresponding analytic function on \(U\) by \(f_n\).  Then
\(\nu^*f_n=b_n\), and therefore
\[
 g_n:=(aN_n)^{-1}\log|f_n|
\]
pulls back to
\(a^{-1}\lambda_i^{(n)}\log|t_i|\) on \(D_i\) for \(i\in I\), and to
\(-\infty\) on \(D_j\) for \(j\in J\).  It also takes the value
\(-\infty\) at \(o\).  Each \(g_n\) is CLD-v2 psh by
\Cref{prop:ch15-log-approachable,prop:ch15-psh-cone}.
\Cref{prop:ch15-log-approachable} applies to every holomorphic function, so
this remains valid when \(J\ne\varnothing\) and \(f_n\) is a zero-divisor;
the proposition's additional non-zero-divisor hypothesis concerns only local
integrability.  Since \(\log|t_i|\leq0\), the sequence decreases pointwise
to the claimed
extension of \(h_c^\circ\), which is CLD-v2 psh by decreasing closedness.
This proves (3)\(\Rightarrow\)(2), while (2)\(\Rightarrow\)(1) follows from
\Cref{thm:cld-to-classical-psh}.

Conversely, let \(h\) be a classically psh extension and suppose that
\(c_i=0\) for some \(i\).  Its pullback \(\widetilde h=h\circ\nu\) is
classically psh by \cite[Proposition~4.2(6)]{GR25Classical}.  The points
\(t_i=\varpi^m\) tend to \(p_i\), so upper
semicontinuity gives \(h(o)=\widetilde h(p_i)\geq0\).  If \(c\ne0\), choose
\(j\) with \(c_j<0\).  On the closed radius-\(r\) subdisc of \(D_j\), the
unique Shilov point has value \(c_j<0\), and every other noncentral value is
at most \(c_j\) (or is \(-\infty\) when \(c_j=-\infty\)).  The affinoid
maximum principle \cite[Theorem~6.4]{GR25Classical} contradicts
\(\widetilde h(p_j)=h(o)\geq0\).  This proves (1)\(\Rightarrow\)(3).

For uniqueness, the zero branch limit gives \(h(o)\geq0\) when \(c=0\),
and the maximum principle on any branch disc gives \(h(o)\leq0\).  In the
all-negative case, a finite coordinate \(c_j\) and the radius-\(r^m\) disc
give \(h(o)\leq m c_j\) for every \(m\), hence \(h(o)=-\infty\); if every
coordinate is \(-\infty\), one branch disc gives the same conclusion.
Finally, a finite rational vector \(c_i=-k_i/N\) can be represented with
all \(k_i\geq e_i\) by scaling numerator and denominator, and the preceding
conductor construction realizes it by one \((aN)^{-1}\log|f|\).  The same
coordinatewise rational approximation proves the asserted decreasing
closure on the radius-\(r\) branch link.
\end{proof}

\begin{example}[The one-branch cusp]
\label{ex:p122-cusp-radial-psh-classification}
Let
\[
 A_{\mathrm{cusp}}:=K\langle u,v\rangle/(v^2-u^3)
      \lhook\joinrel\longrightarrow B:=K\langle t\rangle,
 \qquad u\longmapsto t^2,\quad v\longmapsto t^3.
\]
The decomposition
\(B=K\langle t^2\rangle\oplus tK\langle t^2\rangle\) makes this extension
finite, the element \(t\) is integral because \(t^2=u\), and it equals
\(v/u\) in the fraction field.  Since \(B\) is normal, this is the finite
normalization.  Inside \(B\) one has
\[
 A_{\mathrm{cusp}}=K+t^2K\langle t\rangle,
 \qquad
 \mathfrak c_{B/A_{\mathrm{cusp}}}=t^2K\langle t\rangle.
\]
Indeed, even--odd decomposition gives
\[
 t^2K\langle t\rangle
   =t^2K\langle t^2\rangle\oplus t^3K\langle t^2\rangle
   \subseteq A_{\mathrm{cusp}},
\]
whereas every nonconstant monomial in the generators \(t^2,t^3\) has order
at least two; this proves the first equality.  The displayed ideal is
therefore contained in the conductor.  Conversely, if
\(b=\gamma+t^2g(t)\) lies in the conductor, then
\(tb\in A_{\mathrm{cusp}}\); the coefficient of \(t\) forces
\(\gamma=0\), proving the reverse inclusion.
The normalization is an isomorphism away from the origin, so
\Cref{thm:p122-conductorial-disc-radial-psh-classification} applies with
\(q=1\) and \(e_1=2\).  Hence every
\(c\in[-\infty,0]=\{0\}\cup[-\infty,0)\) has a unique radial extension
which is both CLD-v2 and classically psh.  Its value at the cusp is zero for
\(c=0\) and \(-\infty\) otherwise.  More explicitly, if
\(c=-p/s<0\) with \(p,s\in\mathbf Z_{>0}\), then
\[
       \frac{1}{2as}\log|u^p|
          =\frac{p}{as}\log|t|
\]
is the required extension; increasing rational approximation to \(-c\)
gives the general finite negative case by a decreasing sequence.  Thus the
mixed-zero obstruction is genuinely multibranch and is absent on this
unibranch singularity.
\end{example}

The ordinary multinode makes the conductor construction completely
explicit and will be useful as a permanent boundary example.

\begin{theorem}[Radial psh classification on an ordinary multinode]
\label{thm:p122-multinode-radial-psh-classification}
Let \(K\) be a complete nontrivially valued non-Archimedean field, choose
\(\varpi\in K\) with \(0<|\varpi|<1\), and put
\(r=|\varpi|\) and \(a=-\log r\).  For \(q\geq2\), set
\[
 A_q:=K\langle x_1,\ldots,x_q\rangle/
             (x_ix_j:i\ne j),
 \qquad U_q:=\mathcal M(A_q),
\]
let \(D_i=V(x_j:j\ne i)\) be the \(i\)-th closed-disc branch, and let
\(o=V(x_1,\ldots,x_q)\).  Then
\(U_q=\bigcup_iD_i\) and the branches meet only at \(o\).  Given
\(c=(c_1,\ldots,c_q)\in[-\infty,0]^q\), define
\(h_c^\circ:U_q\setminus\{o\}\to[-\infty,0]\) on the \(i\)-th branch by
\[
 h_c^\circ(z):=
 \begin{cases}
  -\infty,&c_i=-\infty,\\[1mm]
  -\dfrac{c_i}{a}\log|x_i(z)|,&c_i> -\infty.
 \end{cases}
 \tag{\(\blacktriangle\)}\label{eq:p122-multinode-radial-function}
\]
At the radius-\(r\) Gauss point of \(D_i\), this function has value \(c_i\).
The following are equivalent:
\begin{enumerate}
\item \(h_c^\circ\) extends to a classically psh function on \(U_q\);
\item \(h_c^\circ\) extends to a CLD-v2 psh function on \(U_q\);
\item either \(c=0\), or \(c_i<0\) for every \(i\), with
      \(-\infty<0\) allowed.
\end{enumerate}
When it exists, the extension is unique.  Its nodal value is \(0\) when
\(c=0\), and \(-\infty\) otherwise.

Consequently, after identifying the branch link with \(\{1,\ldots,q\}\),
the radial-extendable part of the completed ideal-psh orthant is exactly
\[
                 \{0\}\ \cup\ [-\infty,0)^q.
 \tag{\(\blacklozenge\)}\label{eq:p122-multinode-radial-psh-cone}
\]
It is the zero profile together with the pointwise decreasing closure of
the primary principal-ideal profiles.  In particular, the mixed-zero
profile in \Cref{ex:p122-curve-ideal-psh-boundary} is excluded, whereas
every strictly negative finite profile and every limit obtained by sending
some of its coordinates to \(-\infty\) has a radial CLD-v2 psh extension.
This classifies only the displayed radial realization on the ordinary
multinode; it does not identify arbitrary boundary traces or settle P12.2.
\end{theorem}

\begin{proof}
Assume first that \(c=0\).  The constant zero function is a CLD-v2 psh
extension.  Now suppose that every \(c_i<0\), and put
\[
 I:=\{i:c_i> -\infty\},\qquad J:=\{i:c_i=-\infty\}.
\]
If \(I=\varnothing\), the constant \(-\infty\) function is a CLD-v2 psh
extension.  Otherwise, put \(\lambda_i=-c_i\in(0,+\infty)\) for
\(i\in I\).  Choose positive rational numbers
\[
       \lambda_i^{(n)}=k_{n,i}/N_n\nearrow\lambda_i
       \qquad(i\in I),
\]
using one positive denominator \(N_n\) at each stage, and set
\[
                  f_n:=\sum_{i\in I}x_i^{k_{n,i}}\in A_q,
       \qquad g_n:=(aN_n)^{-1}\log|f_n|.
\]
On \(D_i\) for \(i\in I\), all summands except the \(i\)-th vanish, and
hence
\[
                  g_n|_{D_i}=a^{-1}\lambda_i^{(n)}\log|x_i|.
\]
On \(D_j\) for \(j\in J\), the function \(f_n\) vanishes identically, so
\(g_n|_{D_j}=-\infty\).  In particular, omitting a summand---rather than
sending its exponent to infinity---is what gives the prescribed value at
the outer points with \(|x_j|=1\).
Each \(g_n\) is CLD-v2 psh by
\Cref{prop:ch15-log-approachable,prop:ch15-psh-cone}.  Since
\(\log|x_i|\leq0\), the sequence \((g_n)\) decreases pointwise on every
branch.  Its limit is \(h_c^\circ\) off \(o\), takes the value \(-\infty\)
at \(o\), and is CLD-v2 psh by decreasing closedness.  This proves
(3)\(\Rightarrow\)(2), and (2)\(\Rightarrow\)(1) is
\Cref{thm:cld-to-classical-psh}.

For the converse, suppose that a classically psh extension \(h\) exists and
that \(c_i=0\) for some \(i\).  The rigid points of \(D_i\setminus\{o\}\)
with \(x_i=\varpi^m\) converge to \(o\), so upper semicontinuity gives
\(h(o)\geq0\).  If \(c\ne0\), choose \(j\) with \(c_j<0\).  On the closed
radius-\(r\) disc
\[
                    D_{j,r}:=\{z\in D_j:|x_j(z)|\leq r\},
\]
the unique Shilov point is its radius-\(r\) Gauss point, where \(h=c_j<0\).
When \(c_j\) is finite, every other nonnodal value on this disc is at most
\(c_j\); when \(c_j=-\infty\), all its nonnodal values are \(-\infty\).
Thus \(h(o)\geq0\) contradicts the affinoid maximum principle
\cite[Theorem~6.4]{GR25Classical}.  This proves (1)\(\Rightarrow\)(3).

It remains only to prove uniqueness at the node.  For \(c=0\), upper
semicontinuity gives \(h(o)\geq0\), while the same maximum principle on any
branch disc gives \(h(o)\leq0\).  If all \(c_i<0\) and one \(c_j\) is
finite, applying the maximum principle on the radius-\(r^m\) disc in
\(D_j\) gives \(h(o)\leq m c_j\) for every \(m\), hence
\(h(o)=-\infty\).  If every \(c_i=-\infty\), one branch disc gives the same
conclusion immediately.

Finally, strictly negative rational vectors in \(\mathbf Q_{<0}^q\) are
exactly the radius-\(r\) profiles of the functions
\((aN)^{-1}\log|\sum_i x_i^{k_i}|\).  Coordinatewise
increasing rational approximation of
\((\lambda_i)\in(0,+\infty]^q\) therefore shows that
\([-\infty,0)^q\) is precisely their decreasing pointwise closure.
This last closure statement is on the finite radius-\(r\) branch link; for
the analytic extension of an infinite coordinate, the preceding construction
omits the corresponding summand so that the whole branch, including its
outer boundary, has value \(-\infty\).
The zero profile must be adjoined separately, since every primary
principal-ideal profile is strictly negative on every branch.
\end{proof}


The finite-map proposition first gives link-level functoriality: with the
extended primary ideal on the source, ramification changes numerator and
normalizing denominator by the same factor.  The conductor proposition then
realizes every sufficiently deep integral vector of normalization-branch
orders by one downstairs $\mathfrak m$-primary principal ideal.  The last
two theorems identify the resulting uniform closure and decreasing
completion with a book-local ideal-psh candidate; normalization transport,
finite-map pullback, and the complete obstacle envelope are exact on this
finite link.  For every finite map surjective on branches, generic degrees
give an exact weighted trace, a pullback retraction, and
conductor-controlled principal descent; branch surjectivity is also
necessary for any retraction.  The pushforward satisfies a sharp Jensen
inequality with the ideal envelope, with equality exactly outside mixed-sign
fibers.  It also induces a norm-one projection onto the pullback image and a
dual total-variation-isometric lift of signed measures on the finite link;
the associated source-side Markov projection is positive, norm one, and
adjoint to the function projection.  Relative to every full-support base
probability, the function projection is the finite conditional expectation:
it satisfies the pullback-module identity, contracts every $L^p$ norm, is
orthogonal in $L^2$, and has an exact fiberwise variance decomposition.
Normalized averages and their dual lifts compose exactly when the next-stage
generic degree is fiberwise constant.
For finite flat maps, determinant norms make that descent canonical
without a conductor scaling.  These results do not
compare this class with nef b-divisors or analytic
psh functions.  More sharply,
\Cref{thm:p122-conductorial-disc-radial-psh-classification} computes the
exact radial comparison whenever the finite normalization is a union of
closed discs, is an isomorphism away from one point, and the conductor
separates sufficiently high branch powers.  A branch profile extends
radially as a classically psh function if and only if it also extends as a
CLD-v2 psh function, and this happens exactly for the zero profile or when
every branch coordinate is strictly negative, with $-\infty$ allowed.
The ordinary multinode calculation in
\Cref{thm:p122-multinode-radial-psh-classification} is its explicit
coordinate-ring model.  Thus mixed-zero faces are excluded by the Shilov
maximum principle.  This leaves open a different boundary-trace comparison,
does not by itself close the corresponding P12 atoms, and does not construct
a Monge--Amp\`ere product or energy.

\section{Finite transfer and conductor descent on singular links}
\label{sec:supplied-p126}

Let $(R,\mathfrak m)$ be an excellent reduced Noetherian local ring,
let $\mathfrak a$ be a proper $\mathfrak m$-primary ideal, and let
$R\to B$ be finite with $B\ne0$ reduced. This section follows the supplied
manuscript \cite{SuppliedP12620260924}. Its operators act on continuous
valuative functions and explicit ideal-generated cones. They do not
identify those cones with an independently geometric singular psh class.

Use the primary sections $V_{\mathfrak a}(R)$ of
\cref{sec:local-excellent-boundary-links}; on $B$ use the extended ideal
$\mathfrak aB$. Semivaluations with nonzero prime support are retained.
For a finitely generated ideal $I$, write $g_I(v)=v(I)$. Ideal arithmetic
 gives $g_{IJ}=g_I+g_J$ and $g_{I+J}=\min(g_I,g_J)$, including infinite
values. A gauge of an ideal containing a power of the boundary ideal is
finite and continuous on its section.

\subsection{Restriction and fiber extrema}

\begin{lemma}\label{lem:supplied-p126-finite-restriction}
Restriction is a continuous closed surjection
\[
 f:V_{\mathfrak aB}(B)\longrightarrow
 Z:=\{v\in V_{\mathfrak a}(R):v(\ker(R\to B))=\infty\}.
\]
In particular it is a quotient map onto its image.
\end{lemma}

\begin{proof}
Compactness is the closed-product argument of
\Cref{thm:p121-excellent-primary-links}, which applies equally to the
finite $R$-algebra $B$ and the finitely generated ideal $\mathfrak aB$.
Restriction is continuous and its image lies in $Z$. To lift $v\in Z$,
put $\mathfrak p=\operatorname{supp}(v)$. Lying over for the integral
injection $R/\ker(R\to B)\hookrightarrow B$ supplies a prime
$\mathfrak P$ over $\mathfrak p$. The extension
$\operatorname{Frac}(B/\mathfrak P)/\operatorname{Frac}(R/\mathfrak p)$
is finite. Extend the valuation ring of $v$ to a dominating valuation
ring in this finite field extension, as in
\Cref{lem:p121-normalization-semivaluation-lift}.
An algebraic equation shows that a positive integer multiple of each
new value lies in the original real value group: among its terms at
least two attain the minimal value. Dividing by that integer gives a
well-defined order-preserving real-valued extension. Integral elements
have nonnegative value, since a negative value would make the leading
term in a monic equation uniquely smallest. Restore $\mathfrak P$ with
value $\infty$ to obtain a lift on $B$. Its boundary gauge is one.
Compactness and the Hausdorff property now give the closed quotient map.
\end{proof}

For $H\in C(V_{\mathfrak aB}(B),\mathbf R)$ define, on the exact image $Z$,
\[
 f_{\min}H(v)=\min_{w\in f^{-1}(v)}H(w),\qquad
 f_{\max}H(v)=\max_{w\in f^{-1}(v)}H(w).
\]
The first function is lower semicontinuous and the second upper
semicontinuous. Indeed,
$\{f_{\min}H\leq t\}=f(\{H\leq t\})$ is closed, and the analogous
superlevel identity applies to the maximum. Both operations preserve
order, contract the uniform norm, commute with nonnegative scaling and
addition of a pulled-back base function, and recover that base function
from its pullback. They compose in finite towers: minimize or maximize
first on each intermediate fiber and then on the intermediate points
in the image of the upper map. No finite value is assigned outside the
image of the composite.

\subsection{Field-norm coefficient ideals}

For the next construction let $A$ be an excellent normal domain and
$A\subset B$ a finite domain extension. Put $K=\operatorname{Frac}(A)$,
$L=\operatorname{Frac}(B)$ and $n=[L:K]$. If $J$ is a $B$-ideal, choose
finite $A$-module generators $b_1,\ldots,b_r$ and put
\[
 P_J(T)=N_{L(T)/K(T)}\Bigl(\sum_iT_ib_i\Bigr),\qquad
 N_{B/A}(J)=(\text{coefficients of }P_J)\subset A.
\]
The norm in this definition is a field norm; it does not require the
$A$-module $B$ to be free.

\begin{proposition}\label{prop:supplied-p126-norm-coefficients}
The coefficients belong to $A$, and $N_{B/A}(J)$ is independent of the
chosen generators. Let $\Sigma$ be the set of $K$-embeddings of $L$ in
an algebraic closure, put $q=n/|\Sigma|$, and choose a finite field
extension $M/K$ containing all their images. Let $C$ be the integral
closure of $A$ in $M$. For every nonnegative real semivaluation $v$ of
$A$, choose an extension $w$ to $C$ and set $w_\sigma=w\circ\sigma$ on
$B$. Then
\[
 g_{N_{B/A}(J)}(v)=q\sum_{\sigma\in\Sigma}g_J(w_\sigma).
 \tag{70.7f}\label{eq:supplied-p126-norm-gauge}
\]
This includes semivaluations with nonzero support and infinite ideal
values. In particular, as identities of gauges,
\[
 g_{N_{B/A}(IJ)}=g_{N_{B/A}(I)}+g_{N_{B/A}(J)},\qquad
 g_{N_{B/A}(IB)}=n g_I
\]
for $B$-ideals in the first formula and an $A$-ideal in the second.
\end{proposition}

\begin{proof}
Over an algebraic closure $\Omega$ of $K$, the algebra
$L\otimes_K\Omega$ is a product of Artinian local algebras, indexed by
$\Sigma$, each of length $q$ and residue field $\Omega$. Multiplication
by an element on a factor is its residue scalar plus a nilpotent map.
Consequently
\[
 P_J(T)=\prod_{\sigma\in\Sigma}
                  \Bigl(\sum_iT_i\sigma(b_i)\Bigr)^q.
\]
Each coefficient is integral over $A$ and belongs to $K$; normality puts
it in $A$. Replacing generators by $A$-linear combinations substitutes
linear forms in $P_J$, so the new coefficient ideal is contained in the
old one. Substitution in both directions proves independence.

Excellence makes $C$ finite, and
\Cref{lem:supplied-p126-finite-restriction} supplies $w$; its proof does
not require a local base before imposing a primary section. The minimum
valuation of the coefficients of a polynomial is multiplicative. To
verify this with prime support, first quotient by the support and pass
to the valued fraction field. Divide each nonzero polynomial by a
coefficient of least value. Its reduction over the residue field is
nonzero, and the product of two such reductions is nonzero. This proves
multiplicativity. If a polynomial vanishes modulo the support, the
corresponding values are infinite and the same identity holds.
Apply this identity to the product above. The coefficient minimum of
its $\sigma$-factor is $\min_iw(\sigma(b_i))=g_J(w_\sigma)$, proving
\eqref{eq:supplied-p126-norm-gauge}. The last two formulas follow from
ideal-gauge arithmetic and restriction to $A$.
\end{proof}

The multiplicities $q$ are generic inseparable lengths. Replacing them
by ramification index times residue degree at an arbitrary real
valuation can lose valuation defect and is not part of this formula.

\subsection{Continuous positive trace over a normal target}

Assume now that $(A,\mathfrak m)$ is normal local and excellent, and
normalize by a proper primary ideal $\mathfrak a$. Write
$X=V_{\mathfrak a}(A)$ and $Y=V_{\mathfrak aB}(B)$.

\begin{lemma}\label{lem:supplied-p126-gauge-density}
The rational linear span of $g_J$, for $B$-ideals $J\supseteq\mathfrak a^N B$
with some $N$, is uniformly dense in $C(Y,\mathbf R)$.
\end{lemma}

\begin{proof}
These finite continuous gauges include the constant
$g_{\mathfrak aB}=1$ and separate points: for $b\in B$ the gauge
$g_{(b)+\mathfrak a^N B}(w)=\min(w(b),N)$ distinguishes different values
when $N$ is sufficiently large. After a common denominator, every
rational linear combination has the form $(g_I-g_J)/d$ by products,
powers and the zero gauge $g_B$. Moreover,
\[
 \min(g_I-g_J,g_K-g_H)=g_{IH+KJ}-g_{JH}.
\]
All ideals on the right contain a boundary power. The real linear span
is a separating vector lattice containing constants; the lattice
Stone--Weierstrass argument makes it dense. Explicitly, match values
at pairs of points with affine combinations of separating gauges,
then take finite minima and maxima using compactness. Approximate the
finitely many real coefficients by rationals to obtain the same closure
for the rational span.
\end{proof}

\begin{theorem}\label{thm:supplied-p126-normal-trace}
With the notation of \Cref{prop:supplied-p126-norm-coefficients},
\[
 \mu_v=q\sum_{\sigma\in\Sigma}\delta_{w_\sigma},\qquad
 T_{B/A}H(v)=\int_YH\,\rd\mu_v
 \tag{70.7g}\label{eq:supplied-p126-trace}
\]
is independent of all choices. The measure has mass $n$ and support
exactly the restriction fiber over $v$. The operator
$T_{B/A}:C(Y,\mathbf R)\to C(X,\mathbf R)$ is positive and bounded,
with norm $n$ when $X\ne\varnothing$. It satisfies
\[
 T_{B/A}1=n,\qquad T_{B/A}(f^*h)=nh,\qquad
 T_{B/A}g_J=g_{N_{B/A}(J)}
\]
for every ideal $J$ containing a boundary power. No flatness or
separability hypothesis is imposed.
\end{theorem}

\begin{proof}
By \eqref{eq:supplied-p126-norm-gauge}, the integral of every generating
gauge against $\mu_v$ is independent of the choices. All candidate
measures have mass $n$, so \Cref{lem:supplied-p126-gauge-density} gives
equality on every continuous test function and hence equality of the
finite-support measures.
Each summand lies over $v$. Conversely, given a point $u$ of that fiber,
choose the common field $M$ with the identity embedding of $B$, and
extend $u$ to its integral closure $C$. This is possible because $C$ is
integral and finite over $B$. The identity summand is $u$ and has
positive weight. Choice independence proves that every point of the
fiber has positive weight; in particular there are at most $n$ points.

For $J\supseteq\mathfrak a^N B$, adjoining any $a_0\in\mathfrak a^N$
as a redundant generator shows that $N_{B/A}(J)$ contains $a_0^n$.
Finite generation then implies that it contains a power of
$\mathfrak a$. Its gauge is finite and continuous. Formula
\eqref{eq:supplied-p126-norm-gauge}, density and the common mass bound
therefore prove continuity of $T_{B/A}H$ for every continuous $H$.
Positivity and the displayed identities follow from the fiber measures.
The norm bound is attained on $1$ whenever $X$ is nonempty.
\end{proof}

Let $\mathcal P_{\mathfrak a}(A)$ be the cone of finite nonnegative real
combinations of ideal gauges containing a boundary power; define
$\mathcal P_{\mathfrak aB}(B)$ in the same way. Trace preserves these
cones, and pullback preserves them by extension of ideals. This is an
explicit ideal-generated class.

For a finite tower $A\subset B\subset D$ of excellent normal domains,
with $A$ local, the traces compose:
\[
 T_{D/A}=T_{B/A}\circ T_{D/B}.
 \tag{70.7h}\label{eq:supplied-p126-tower}
\]
To prove this, choose one finite extension of $\operatorname{Frac}(A)$
containing every embedding of $\operatorname{Frac}(D)$, and extend the
base semivaluation once. Group the top embeddings by their restriction
to the intermediate field. Every intermediate embedding extends, and
both field degrees and separable embedding counts multiply. Thus the
inseparable multiplicities multiply as well, and the resulting Dirac
sum is the iterated fiber measure. Choice independence proves the
formula. The proof also applies to a normal semilocal intermediate
ring on the common pulled-back section, using its finitely many local
center pieces. In this domain setting the degree-normalized averages
also compose, since the degrees multiply.

\begin{example}\label{ex:supplied-p126-nonflat-normal-trace}
For $\operatorname{char}k\ne2$, put
\[
 A=k[x^2,xy,y^2]_{(x^2,xy,y^2)},\qquad B=k[x,y]_{(x,y)}.
\]
The sign-invariant ring $A$ is normal: an element of its fraction field
integral over $A$ is in the normal ring $B$, and invariance puts it in
$A$. It is excellent, and $B/A$ is finite of generic degree two.
The residue space $B/\mathfrak m_AB$ has basis $1,x,y$, so the finite
module is not flat: a finite flat module over a local ring is free,
whereas its minimal generator number is three and its generic rank two.
For $J=(x,y)B$, use the $A$-module generators $x,y,x^2,xy,y^2$.
Writing $E=T_3x^2+T_4xy+T_5y^2$, the norm polynomial is
$E^2-(T_1x+T_2y)^2$, whose coefficient ideal is $\mathfrak m_A$.
Thus $T_{B/A}g_J=g_{\mathfrak m_A}$ on the whole section, including
semivaluations with nonzero support.
\end{example}

\subsection{The exact conductor-compatible domain}

Return to the finite map $R\to B$ at the start of the section. Write
the finite normalizations as $S_R=\prod_iA_i$ and $S_B=\prod_jD_j$,
where the factors are normal domains, and let $q_R:W_R\to X$ and
$q_B:W_B\to Y$ be their primary-link quotients. A source component $D_j$
is active over $A_i$ if its minimal prime contracts to the minimal prime
defining $A_i$. The induced inclusion $A_i\subset D_j$ is finite.
Components with positive-codimension image are called vertical here;
no generic degree is assigned to them. On normalization data set
\[
 (T_f^\nu H)|_{A_i}=\sum_{j\text{ active over }i}T_{D_j/A_i}(H|_{D_j}),
 \tag{70.7i}\label{eq:supplied-p126-branch-trace}
\]
with an empty sum equal to zero. The common-section proofs above apply
to these normal semilocal factors. This defines a positive bounded
operator. Fiber extrema still detect vertical components.

Define $\mathcal D_f\subset C(Y,\mathbf R)$ to consist of those $h$ for
which $T_f^\nu(q_B^*h)$ is constant on every fiber of $q_R$.
For such an $h$ define $\operatorname{Tr}_fh$ by
$q_R^*(\operatorname{Tr}_fh)=T_f^\nu(q_B^*h)$.

\begin{theorem}\label{thm:supplied-p126-conductor-trace}
The space $\mathcal D_f$ is a closed vector subspace, and
$\operatorname{Tr}_fh$ is the unique continuous descent of the branch
trace. Its compatibility equalities occur exactly on the normalization
fibers described by the conductor quotient in
\Cref{thm:p121-excellent-conductor-descent}.
For a finite tower $R\to B\to D$, if $h\in\mathcal D_{D/B}$ and
$\operatorname{Tr}_{D/B}h\in\mathcal D_{B/R}$, then
$h\in\mathcal D_{D/R}$ and
\[
 \operatorname{Tr}_{D/R}h
   =\operatorname{Tr}_{B/R}(\operatorname{Tr}_{D/B}h).
\]
Direct admissibility alone is not an intermediate descent hypothesis.
\end{theorem}

\begin{proof}
Fiber constancy is a closed linear condition for the uniform norm.
The compact Hausdorff quotient $q_R$ makes it necessary and sufficient
for a unique continuous descent. A finite-valued conductor point has
one normalization lift; on the infinite-conductor stratum the exact
fibers are those of restriction through $S_R/\mathfrak c_R$ to
$R/\mathfrak c_R$, by \Cref{thm:p121-excellent-conductor-descent}.

A minimal component of $D$ dominating a minimal component of $R$ passes
through a unique minimal component of $B$ dominating that component:
incomparability in an integral extension makes a prime over a minimal
target prime minimal in the source. Conversely such an active chain
is active for the composite. Summing \eqref{eq:supplied-p126-tower}
over these chains gives the identity on normalizations. The two assumed
descents and uniqueness yield the displayed identity downstairs and
prove direct admissibility.
\end{proof}

Define $\mathcal P_{\mathfrak a}^{\nu}(R)$ to be the continuous functions
whose pullbacks to all normal components lie in their ideal-generated
cones. This includes both the branchwise ideal condition and the fiber
equalities. Equation \eqref{eq:supplied-p126-branch-trace} gives
\[
 \operatorname{Tr}_f\bigl(\mathcal P_{\mathfrak aB}^{\nu}(B)
                         \cap\mathcal D_f\bigr)
                  \subseteq\mathcal P_{\mathfrak a}^{\nu}(R).
\]
This normalization-ideal cone is not identified with the singular nef
or analytic psh cone.

There is a useful sufficient criterion for realization by an actual
ideal. If $S$ is the normalization, $\mathfrak c$ its conductor,
$J\supseteq\mathfrak c$ an $S$-ideal and $K\subset R/\mathfrak c$ an
ideal with $J/\mathfrak c=K(S/\mathfrak c)$, let $I$ be the inverse
image of $K$ in $R$. Then $IS=J$: both contain $\mathfrak c$ and have
the same image modulo it. Conversely every $R$-ideal containing the
conductor has this description after extension. Fiber-compatible
continuous gauges need not be gauges of individual downstairs ideals.
For the cusp $k[t^2,t^3]_{(t^2,t^3)}\subset k[t]_{(t)}$, the gauge of
$(t)$ on the maximal-ideal section is $1/2$, whereas every proper
nonzero downstairs ideal has value at least one. The former is still
the real multiple $g_{\mathfrak m}/2$; this example distinguishes a
single ideal from the real cone.

\subsection{The obstruction and the finite-flat endpoint}

\begin{example}\label{ex:supplied-p126-conductor-obstruction}
Let
\[
 \begin{gathered}
 R=k[x,y,z]_{(x,y,z)}/(xy),\qquad
 S=k[x,z]_{(x,z)}\times k[y,z]_{(y,z)},\\
 J=(x,z)\times(y,z^2).
 \end{gathered}
\]
Normalize by $\mathfrak m=(x,y,z)$. On the first branch take monomial
semivaluations $u_t(x)=t$, $u_t(z)=1$; on the second take
$w_t(y)=t$, $w_t(z)=1$, for integers $t\geq2$. Their restrictions both
converge pointwise to the semivaluation supported on $(x,y)$ and inducing
$\operatorname{ord}_z$ on $R/(x,y)$. Indeed, transverse terms acquire
unbounded value and local denominators have value zero. But
$g_J(u_t)=1$ and $g_J(w_t)=2$. On either generic branch the degree-one
trace must equal this source value. No continuous function on the target
can have both limits at the common conductor point. Thus a finite
birational surface map need not admit a continuous degree trace on all
source ideal gauges.
\end{example}

\begin{proposition}\label{prop:supplied-p126-finite-flat-trace}
If $R\to B$ is finite locally free of constant rank $n$, then
$\mathcal D_f=C(Y,\mathbf R)$. Its trace is the continuous rank-$n$
determinant trace. For a $B$-ideal $J$ containing a boundary power, its
image is the gauge of the coefficient ideal of
\[
 \det\Bigl(m_{\sum_iT_ib_i}:B[T]\longrightarrow B[T]\Bigr),
\]
where the $b_i$ generate $J$ as an $R$-module.
\end{proposition}

\begin{proof}
The determinant exists and is basis independent. Linear substitution
in both directions proves generator independence of its coefficient
ideal. Flat going down shows that every minimal source component
dominates a minimal target component. After localization at a minimal
prime of the reduced base, the reduced finite algebra splits into the
fraction fields of those components. Its determinant polynomial is
the product of their field norms. The coefficient-minimum argument in
\Cref{prop:supplied-p126-norm-coefficients} shows that the branch trace
of $g_J$ is the pullback of the downstairs coefficient-ideal gauge.
Thus every such gauge is admissible. The density proof of
\Cref{lem:supplied-p126-gauge-density} uses only compactness and ideal
arithmetic, so it applies on $Y$ here. Closedness of $\mathcal D_f$
therefore makes every continuous function admissible. This is the
finite locally free norm of \cite[Tag~0BD2]{StacksProject}; flat going
down is \cite[Tag~00HS]{StacksProject}.
\end{proof}

For a reduced excellent curve, the conductor quotient is zero-dimensional
and its primary link is empty. Thus the normalization quotient is the
finite set of DVR branches. If $A$ is an excellent DVR with uniformizer
$\pi$ and $D$ its normalization in a finite field extension, write
$e_i=\operatorname{ord}_i(\pi)$ and
$f_i=[D/\mathfrak n_i:A/(\pi)]$. Finite freeness and the length of
$D/\pi D$ give $[\operatorname{Frac}(D):\operatorname{Frac}(A)]=\sum_ie_if_i$.
Smith normal form for multiplication gives
$\operatorname{ord}_A N(b)=\sum_i f_i\operatorname{ord}_i(b)$.
On the $\pi^r$-section the branch point is
$\operatorname{ord}_i/(re_i)$, so its trace weight is $e_if_i$.
These formulas use finite freeness over a DVR, not a defectless
hypothesis for arbitrary valuation rings.

Degree-normalized averages need not compose when intermediate ranks
vary by component. Let $A=k[[t]]$, $B=A\times A$ with diagonal map,
and $D=A\times k[[u]]$, with the second map $t\mapsto u^2$.
The direct average has weights $(1/3,2/3)$, whereas iterating the
componentwise averages gives $(1/2,1/2)$. Their values at $(0,1)$ are
$2/3$ and $1/2$. The unnormalized traces both have weights $(1,2)$
and compose, including in characteristic two.

\subsection{Changing the primary section}

For a homogeneous function $H$ of degree one on the centered cone,
let $h_{\mathfrak a}$ be its representative on $V_{\mathfrak a}$.
The exact change to the $\mathfrak b$-section is
\[
 h_{\mathfrak b}\bigl(v/v(\mathfrak b)\bigr)
                  =h_{\mathfrak a}(v)/v(\mathfrak b).
 \tag{70.7j}\label{eq:supplied-p126-radial-transfer}
\]
Every extension of $v$ has the same $\mathfrak b$-gauge. Hence the
formula commutes with trace, fiber extrema and conductor descent.
For a separately chosen source boundary $\mathfrak d$, first express
its representative on the $\mathfrak aB$-section as
$w(\mathfrak d)h_{\mathfrak d}(w/w(\mathfrak d))$.
For nonprimary gauges this conversion is confined to the domain where
the gauge is positive and finite, as in
\eqref{eq:p121-nonprimary-gauge-domain}.

The section establishes finite-map calculus on continuous valuation-link
functions, its exact conductor domain and a counterexample to a universal
continuous degree trace. Constructing independently geometric singular
psh cones and their finite-map, envelope and Monge--Amp\`ere compatibility
remains a separate question.


\section{Exceptional b-divisors and ideal functions}
\label{sec:local-exceptional-ideals}

We return to the smooth complex germ and the simple modifications fixed
at the beginning of the chapter.
Let $\mathfrak X_0=\varprojlim_\pi Y_\pi$ be the local
Riemann--Zariski space.  An exceptional Weil divisor is a compatible family
\[
                         Z=(Z_\pi)_\pi,
 \qquad
                         \mu_*Z_{\pi'}=Z_\pi
\tag{70.8}\label{eq:local-weil-bdivisor}
\]
whenever $\pi'=\pi\circ\mu$.  A Cartier divisor is determined on one
model and pulled back to every higher model.  Coefficients along prime
exceptional divisors identify Weil divisors with functions on the
divisorial locus:
\[
 g_Z(\nu_E)
 :=\frac{\ord_E(Z)}{\ord_E(\mathfrak m)}.
\tag{70.9}\label{eq:local-divisor-function}
\]

For an $\mathfrak m$-primary ideal $\mathfrak a\subset R$ (identified
with its extension to $\widehat R$), define
\[
 \log|\mathfrak a|(\nu):=-\nu(\mathfrak a),
 \qquad
 \ord_E Z(\mathfrak a):=-\ord_E(\mathfrak a).
\tag{70.10}\label{eq:local-ideal-function}
\]
For a single formal function, $Z(f)$ means the exceptional Weil divisor
whose coefficient at $E$ is $-\ord_E(f)$; equivalently, it is the
divisor attached to the principal ideal $(f)$.
Thus $Z(\mathfrak a)\leq0$, and
$g_{Z(\mathfrak a)}=\log|\mathfrak a|$ on divisorial valuations.
The sign agrees with the global convention of this book: an effective ideal
has a nonpositive logarithmic norm.

The tropical identities are
\[
\begin{aligned}
 \log|\mathfrak a\mathfrak b|
   &=\log|\mathfrak a|+\log|\mathfrak b|,\\
 \log|\mathfrak a+\mathfrak b|
   &=\max\{\log|\mathfrak a|,\log|\mathfrak b|\}.
\end{aligned}
\tag{70.11}\label{eq:local-tropical-ideal-laws}
\]
Cartier b-divisors with real coefficients correspond exactly to continuous
functions that become real affine on every face of one sufficiently high
dual complex.  The
integral subgroup is generated by the functions
$\log|\mathfrak a|$ for primary ideals $\mathfrak a$.  This is the
local counterpart of the model-function/vertical-Cartier dictionary in
\Cref{chap:models,chap:mp-dual-complexes}; see
\cite[Propositions~1.14--1.15]{BFJ08} for the local dictionary and its
uniform-density statement.

\begin{example}[The distinguished local constant]
The normalization $\nu(\mathfrak m)=1$ gives
\[
                         \log|\mathfrak m|\equiv-1
                         \quad\text{on }\mathcal V_0.
\tag{70.12}\label{eq:local-maximal-ideal-constant}
\]
Hence the maximal ideal plays simultaneously the role of radial coordinate,
normalizing divisor, and constant model potential.  This is why adding an
arbitrary constant before choosing a normalization is not an intrinsic local
operation.
\end{example}

\begin{example}[A monomial ideal]
On $(\mathbf C^2,0)$, take
$\mathfrak a=(x^p,y^q)$, where $p,q\geq1$.  For the monomial valuation of weight
$(a,b)$, normalized by $\min\{a,b\}=1$,
\[
                   \log|\mathfrak a|(\nu_{a,b})
                    =-\min\{pa,qb\}
                    =\max\{-pa,-qb\}.
\tag{70.13}\label{eq:local-monomial-ideal}
\]
The corner occurs where $pa=qb$.  It is a local model-psh function, but
it is not smooth across that wall.  This is the same regularity distinction
that appears for Fubini--Study maxima in \Cref{chap:bj-fs-pl-psh}.
\end{example}

\section{Formal psh functions}
\label{sec:local-formal-psh}

The cone of nef Weil divisors is the closed convex cone generated by
$Z(\mathfrak a)$ as $\mathfrak a$ ranges over ideals of the completed
regular local ring.  It is enough to use the unit ideal, which gives the zero
divisor, together with $\mathfrak m$-primary ideals.  If a nef divisor is
Cartier and determined on $Y_\pi$, nefness is equivalent to
$\pi$-nefness: its degree on every complete curve in $\pi^{-1}(0)$ is
nonnegative \cite[Proposition~2.2]{BFJ08}.

\begin{definition}
\label{def:local-formal-psh}
A formal psh function is a function
$g:\mathcal V_{0,\mathrm{qm}}\to\mathbf R$ that is continuous in the
strong topology and whose divisorial values define a nef Weil divisor
$Z(g)$.  It is extended to all of $\mathcal V_0$ by
\[
                           g(\nu):=\inf_\pi g(r_\pi(\nu)).
\tag{70.14}\label{eq:local-formal-extension}
\]
\end{definition}

The extension takes values in $[ -\infty,0]$ and is upper
semicontinuous for the weak topology.  On every simplex it is finite,
continuous, and convex.  The sign may initially look reversed: the
functions are nonpositive, while their restrictions to affine weight
coordinates are convex.

We need one discrepancy term for the approximation theorem.  If $E$ is an
exceptional prime on $Y_\pi$, define the \emph{thin divisor}
$\mathbf A$ and its normalized thinness function by
\[
 \ord_E(\mathbf A)=1+\ord_E(K_{Y_\pi/Y}),
 \qquad
 A(\nu_E)=\frac{1+\ord_E(K_{Y_\pi/Y})}{b_E}.
\]
The function $A$ is finite, affine, and bounded on each fixed simplex,
though it can take the value $+\infty$ on
$\mathcal V_0\setminus\mathcal V_{0,\mathrm{qm}}$.
For a nef Weil divisor $W$, its \emph{valuative multiplier ideal} is the
stationary value, for $0<\epsilon\ll1$, of
\[
 \mathcal L^2(W)
   =\{f\in\widehat R:
          Z(f)\leq(1+\epsilon)W+\mathbf A\}.
\]
Noetherianity makes the increasing family on the right stationary as
$\epsilon\downarrow0$.

\begin{theorem}\label{thm:local-nef-psh-equivalence}
The assignments $g\mapsto Z(g)$ and $Z\mapsto g_Z$ give a bijection
between formal psh functions and nef Weil divisors.  Every formal psh
function admits a decreasing ideal approximation
\[
                  g_k=c_k\log|\mathfrak a_k|\searrow g,
 \qquad c_k>0,
\tag{70.15}\label{eq:local-decreasing-ideal-approximation}
\]
where every $\mathfrak a_k\subset\widehat R$ is either the unit ideal or
$\mathfrak m$-primary.  If $g\not\equiv0$, the sequence may be reindexed
so that every $\mathfrak a_k$ is $\mathfrak m$-primary; if $g\equiv0$,
one may take $\mathfrak a_k=\widehat R$ and $g_k=0$ for all $k$.
\end{theorem}

\begin{proof}
For a nef Weil divisor, divisorial coefficients give a convex function on
the rational points of each dual face.  Continuity up to the boundary of a
closed face is the subtle point.  BFJ attach valuative multiplier ideals
$\mathcal L^2(kZ)$ and prove the uniform estimate
\[
 \left|g_Z-\frac1k\log|\mathcal L^2(kZ)|\right|
                   \leq\frac{A}{k}
\tag{70.16}\label{eq:local-multiplier-uniform-estimate}
\]
on every fixed simplex, where the log-discrepancy function $A$ is bounded
there \cite[Theorem~3.9]{BFJ08}.  The ideal functions therefore converge
uniformly on each face and
produce the unique strongly continuous extension.  The formula
\eqref{eq:local-formal-extension} gives its weakly usc extension to all
valuations.

For monotone approximation, BFJ take, after dyadic indexing,
\[
 \mathfrak a_k=\mathcal L^2(2^kZ(g))+\mathfrak m^{4^k},
 \qquad c_k=2^{-k}.
\]
The valuative subadditivity theorem makes the resulting functions decrease
\cite[Theorem~3.10]{BFJ08},
and the sum is either the unit ideal or $\mathfrak m$-primary.  For
$g\equiv0$, one has $\mathcal L^2(0)=\widehat R$, so the displayed
sequence is identically zero.  If $g\not\equiv0$, then
$\max_{\mathcal V_0}g<0$ by \cite[Proposition~3.14]{BFJ08}.  The upper
estimate in \cite[Theorem~3.9]{BFJ08}, evaluated at any fixed
quasimonomial valuation, therefore forces $\mathcal L^2(2^kZ(g))$ to be
proper for all sufficiently large $k$.  Adding $\mathfrak m^{4^k}$
then makes it $\mathfrak m$-primary.  Discarding the finite initial segment
gives the asserted nonzero case.  The limit on non-quasimonomial valuations follows from
\eqref{eq:local-formal-extension}.  See
\cite[Propositions~3.11 and 3.13]{BFJ08}.
\end{proof}

The formal psh cone is closed under positive linear combinations and
maxima.  Nefness is also weakly closed under coefficientwise limits whose
coefficients remain finite,
but such a limit represents a formal psh function only after strong
continuity on every simplex has been checked; no unconditional
monotone-closure assertion is intended.  Every formal psh function satisfies
\[
                              g\leq g\circ r_\pi.
\tag{70.17}\label{eq:local-retraction-inequality}
\]
The right side decreases to $g$ as the model increases
\cite[Proposition~2.7]{BFJ08}.

\begin{remark}[Retraction trap]
\label{rem:local-retraction-trap}
In dimension two, $g\circ r_\pi$ is again formal psh.  In higher
dimension this may fail: the Cartier incarnation of a nef Weil divisor on a
fixed model need not be nef.  Consequently local regularization cannot be
reduced to a single fixed skeleton by formal analogy with a metric graph.
\end{remark}

\section{Nef envelopes on a fixed local model}
\label{sec:local-nef-envelopes}

Suppose a bounded function $h$ is determined on one dual complex, meaning
$h=h\circ r_\pi$.  Its associated Weil divisor is determined by the
divisorial values of $h$.  BFJ define its nef envelope by
\[
 \widehat Z(h)
 :=\sup\{W:W\text{ is a nef Weil divisor and }W\leq Z(h)\}.
\tag{70.18}\label{eq:local-nef-envelope}
\]
The admissible family is nonempty: since $h$ is bounded, a sufficiently
large multiple of $Z(\mathfrak m)$ lies below $Z(h)$.
Finite coefficientwise maxima are nef.  Direct the finite subfamilies by
inclusion; their maxima increase coefficientwise to the displayed
supremum, which remains finite because it is bounded above by $Z(h)$.
Closedness of the nef cone therefore makes the supremum nef.  Equivalently,
one obtains the largest formal psh function dominated by $h$ on
divisorial valuations.

For a Weil divisor $W$, put
\[
 \mathfrak a(W)
 :=\{f\in\widehat R:Z(f)\leq W\}.
\tag{70.19}\label{eq:local-envelope-ideal}
\]
If $W$ is determined on a bounded dual complex, then
$\mathfrak a(W)\neq0$, since a sufficiently high power of
$\mathfrak m$ satisfies the required inequalities.  For the nef envelope
$W=\widehat Z(h)$, the ideals $\mathfrak a(kW)$ recover $W$ with
quantitative multiplier-ideal estimates
\cite[Section~3.5, especially Proposition~3.18]{BFJ08}.

This is a genuine local envelope construction, but its signature matters.
The obstacle is bounded and carried by a fixed dual complex; the output is a
nef b-divisor on the tower over one smooth point.  It is not a proof that the
usc regularization of an arbitrary local FS supremum is psh on an arbitrary
Berkovich analytic germ.  In fact,
\Cref{prop:ch15-nodal-local-envelope} shows that the unrestricted assertion
fails for a continuous obstacle on a reduced nodal affinoid curve.

\section{Intersection products and local Monge--Amp\`ere measures}
\label{sec:local-relative-ma}

On one modification, exceptional Cartier divisors have their usual
intersection number
\[
                         \langle Z_1,\ldots,Z_n\rangle\in\mathbf R.
\tag{70.20}\label{eq:local-cartier-intersection}
\]
For nef Cartier divisors it is nonpositive and monotone in every variable.
BFJ extend it to nef Weil divisors by
\[
 \langle Z_1,\ldots,Z_n\rangle
 :=\inf_{W_i\geq Z_i}
       \langle W_1,\ldots,W_n\rangle\in[-\infty,0],
\tag{70.21}\label{eq:local-nef-intersection}
\]
where the $W_i$ range over nef Cartier divisors.  The product is symmetric,
multilinear, upper semicontinuous, and continuous along decreasing families
\cite[Proposition~4.4]{BFJ08}.  Multilinearity at the Weil level is a deep use of
the ideal approximation in \Cref{thm:local-nef-psh-equivalence}; it is not a
formal property of arbitrary inverse-limit divisors.

\begin{theorem}\label{thm:local-formal-ma}
For $1\leq i\leq n-1$, let $g_i$ be a formal psh function with nef
divisor $Z_i$.  Assume
\[
        \langle Z_1,\ldots,Z_{n-1},Z(\mathfrak m)\rangle> -\infty.
\tag{70.22}\label{eq:local-ma-finiteness}
\]
There is a unique positive Radon measure
$\MA_0(g_1,\ldots,g_{n-1})$ on $\mathcal V_0$ such that
\[
 \int_{\mathcal V_0}g\,
       \MA_0(g_1,\ldots,g_{n-1})
   =\langle Z(g),Z_1,\ldots,Z_{n-1}\rangle
\tag{70.23}\label{eq:local-ma-characterization}
\]
for every formal psh test function $g$.  The hypothesis holds when all
$g_i$ are bounded.
\end{theorem}

\begin{proof}
First take $g$ piecewise affine, so that $Z(g)$ is Cartier.  Any Cartier
divisor is a difference of two nef Cartier divisors, and
\eqref{eq:local-ma-finiteness} makes the resulting linear functional finite.
It is positive on nonnegative piecewise-affine functions.  Such functions
form a uniformly dense lattice in $C^0(\mathcal V_0)$, by the ideal
functions and the lattice form of Stone--Weierstrass.  The functional is
therefore bounded: if $M=L(1)=-\langle
Z(\mathfrak m),Z_1,\ldots,Z_{n-1}\rangle$, positivity applied to
$\|g\|_\infty\pm g$ gives $|L(g)|\leq M\|g\|_\infty$.
It extends uniquely to $C^0(\mathcal V_0)$.  The Riesz representation theorem gives the
measure.  Decreasing ideal approximation then proves
\eqref{eq:local-ma-characterization} for every formal psh test function.  See
\cite[Proposition~4.6]{BFJ08}.
If the $g_i$ are bounded, choose $C_i\geq0$ with
$g_i\geq-C_i$.  Then $Z_i\geq C_iZ(\mathfrak m)$, so monotonicity and
$-\langle Z(\mathfrak m),\ldots,Z(\mathfrak m)\rangle
=e(\mathfrak m)=1$ give the required finite lower bound.
\end{proof}

Taking the test function $\log|\mathfrak m|=-1$ yields the total mass
\[
 \bigl\|\MA_0(g_1,\ldots,g_{n-1})\bigr\|
 =-\langle Z(\mathfrak m),Z_1,\ldots,Z_{n-1}\rangle.
\tag{70.24}\label{eq:local-ma-total-mass}
\]
For primary ideals $\mathfrak a_1,\ldots,\mathfrak a_n$, write
$e\langle\mathfrak a_1,\ldots,\mathfrak a_n\rangle$ for their mixed
Hilbert--Samuel multiplicity.  The identity is
\[
 e\langle\mathfrak a_1,\ldots,\mathfrak a_n\rangle
 =\int_{\mathcal V_0}-\log|\mathfrak a_1|\,
   \MA_0(\log|\mathfrak a_2|,\ldots,
                 \log|\mathfrak a_n|).
\tag{70.25}\label{eq:local-mixed-multiplicity}
\]
This is \cite[Proposition~4.7]{BFJ08}.
If the $g_i$ are piecewise affine on one model, the measure is atomic at
its divisorial vertices.  More precisely, the mass at
$\nu_E=b_E^{-1}\ord_E$ is
\[
 b_E\bigl\langle E,(Z_1)_\pi,\ldots,(Z_{n-1})_\pi
             \bigr\rangle_{Y_\pi}.
\tag{70.26}\label{eq:local-atomic-ma}
\]
This model formula is the calculation in
\cite[Proposition~4.9]{BFJ08}.
Finally, decreasing sequences in all $g_i$ give weak convergence of
their Monge--Amp\`ere measures whenever the limiting measure is defined
\cite[Proposition~4.8]{BFJ08}.

\begin{example}[The maximal ideal gives a Dirac mass]
Set $g_1=\cdots=g_{n-1}=\log|\mathfrak m|=-1$.  This model function is
determined, when $n\geq2$, by the blowup of $0$, whose exceptional
divisor gives the single vertex $\nu_{\mathfrak m}$.  Since
$e(\mathfrak m)=1$, \eqref{eq:local-mixed-multiplicity} and
\eqref{eq:local-atomic-ma} give
\[
              \MA_0(-1,\ldots,-1)=\delta_{\nu_{\mathfrak m}}.
\tag{70.27}\label{eq:local-maximal-ideal-dirac}
\]
For $n=1$, one instead uses the divisor $\{0\}$ on the identity model;
the left side has no potential arguments, and the formula recovers the
one-point calculation at the start of the chapter.
\end{example}

The operator takes $n-1$ potentials because $\mathcal V_0$ is a
normalized link of dimension $n-1$.  The extra divisor in
\eqref{eq:local-ma-characterization} is supplied by the test function.  This
differs from the global Monge--Amp\`ere operator on an $n$-dimensional
projective variety, which takes $n$ curvature factors.

\section{Complex psh germs and their valuative transforms}
\label{sec:local-valuative-transform}

There is a convention change between \cite{BFJ08} and the complex part of
this book.  BFJ write singularities using $u\sim\log|f|$ and define
multiplier ideals by $|f|^2\rme^{-2u}$.  In
\Cref{chap:complex-prerequisites}, we write
$\varphi=2u\sim\log|f|^2$ and use
$|f|^2\rme^{-\varphi}$.  Thus
\[
             u=\frac12\varphi,
 \qquad
             \mathcal J_{\mathrm{BFJ}}(u)
                =\mathcal I_{\mathrm{book}}(\varphi).
\tag{70.28}\label{eq:local-convention-firewall}
\]
All statements in this section are first written in the BFJ variable $u$.

Let $u$ be a psh germ at $0$.  For every exceptional prime $E$ on a
modification, let
\[
 \nu_E(u):=b_E^{-1}\nu^{\mathrm{gen}}(u\circ\pi,E)
\]
be the normalized generic Lelong number of $u\circ\pi$ along $E$.
The coefficients
\[
                         \ord_E Z(u):=-b_E\nu_E(u)
\tag{70.29}\label{eq:local-psh-divisor}
\]
define a Weil divisor.

\begin{theorem}\label{thm:local-valuative-transform}
The divisor $Z(u)$ is nef.  Its corresponding formal psh function
$\widehat u$ is characterized on divisorial valuations by
\[
                            \widehat u(\nu_E)=-\nu_E(u).
\tag{70.30}\label{eq:local-valuative-transform-divisorial}
\]
If $u=c\max_j\log|f_j|+O(1)$ and
$\mathfrak a=(f_j)$, then
\[
                            \widehat u=c\log|\mathfrak a|.
\tag{70.31}\label{eq:local-analytic-ideal-transform}
\]
On a quasimonomial simplex, $-\widehat u$ is the corresponding Kiselman
number of $u\circ\pi$.
\end{theorem}

\begin{proof}
Demailly's Hilbert-space approximation gives psh germs $u_k$ with ideal
singularities determined by $\mathcal J(ku)^{1/k}$.  The divisorial
estimate
\[
 Z(\mathcal J(u))\geq Z(u)
            \geq Z(\mathcal J(u))-\mathbf A
\tag{70.32}\label{eq:local-demailly-divisor-estimate}
\]
implies
$Z(u)=\lim_k k^{-1}Z(\mathcal J(ku))$.  Each ideal divisor is nef, and
the nef cone is closed, proving nefness.  The first characterization is the
definition of the divisor coefficients; the ideal case follows by taking
orders of the generators.  Rational monomial weights are divisorial, and
continuity on each simplex extends the equality with Kiselman numbers to all
real weights.  These are \cite[Theorem~5.1 and
Proposition~5.6]{BFJ08}.
\end{proof}

The valuative transform remembers the full collection of generic Lelong
numbers over every model.  It also controls multiplier ideals.  In BFJ's
notation,
\[
                 \mathcal L^2(Z(u))\cap\cO_{Y,0}
                    =\mathcal J_+(u),
\tag{70.33}\label{eq:local-valuative-multiplier}
\]
where $\mathcal J_+(u)$ is the stationary value of
$\mathcal J((1+\epsilon)u)$; this is
\cite[Theorem~5.5]{BFJ08}.  Strong openness, proved later by
Guan--Zhou \cite[Theorem~1.1]{GZ15}, in the form recorded in
\Cref{thm:complex-strong-openness}, identifies this ideal with $\mathcal J(u)$.
Thus one historical question recorded in \cite{BFJ08} is now resolved; this
does not resolve the separate generalized-Lelong-number question below.

A \emph{psh weight} $\psi$ has an isolated pole at $0$ and
$\rme^\psi$ continuous there.  It is \emph{maximal} when
$(\ddc\psi)^n=0$ on a punctured neighborhood of $0$, and it is
\emph{tame} when there is $C>0$ such that, for all $t>0$ and
$f\in\mathcal J(t\psi)$,
\[
                    \log|f|\leq(t-C)\psi+O(1).
\]
For a psh germ $u$, its relative type is
\[
 \sigma(u,\psi):=\sup\{c\geq0:u\leq c\psi+O(1)\}.
\]
For a tame maximal psh weight $\psi$, BFJ prove
\[
\sigma(u,\psi)
 =\sup\{c\geq0:\widehat u\leq c\widehat\psi\},
\tag{70.34}\label{eq:local-relative-type}
\]
and every quasimonomial valuation $\nu$ is represented by a tame maximal
weight $\psi_\nu$ satisfying
\[
                      -\widehat u(\nu)=\sigma(u,\psi_\nu)
\]
for every psh germ $u$.  These are
\cite[Theorems~5.12--5.13]{BFJ08}.  If $\psi$ is a tame psh weight, then
the generalized Lelong number is
\[
 \nu_\psi(u)
 =\bigl(\ddc u\wedge(\ddc\psi)^{n-1}\bigr)\{0\}
 =\int_{\mathcal V_0}-\widehat u\,
             \MA_0(\widehat\psi,\ldots,\widehat\psi).
\tag{70.35}\label{eq:local-generalized-lelong}
\]
This is \cite[Theorem~B]{BFJ08}.  The displayed $\ddc$ and its
normalization are those of the BFJ source.  If
$\Phi=2u$ and $\Psi=2\psi$ are the book variables while the differential
normalization is held fixed, then
\[
 \ddc\Phi\wedge(\ddc\Psi)^{n-1}
    =2^n\ddc u\wedge(\ddc\psi)^{n-1}.
\]
On the valuative side $\widehat\Phi=2\widehat u$, while the
$(n-1)$-linear measure scales by $2^{n-1}$; hence the same factor
$2^n$ appears.  Converting only one factor would be incorrect.

\begin{theorem}\label{thm:local-generalized-lelong-separation}
Let \(u,v:(\mathbf C^n,0)\to\mathbf R\cup\{-\infty\}\) be psh germs,
where \(n\geq1\), and normalize the centered valuation space by
\(\nu(\mathfrak m)=1\).  If
\[
                         \nu_\psi(u)=\nu_\psi(v)
\]
for every tame psh weight \(\psi\), then their valuative transforms agree
pointwise:
\[
                         \widehat u=\widehat v
                  \quad\text{on }\mathcal V_0.
\]
It is enough to test the analytic weights
\[
 \psi_{\mathfrak a}=\log\max_j|f_j|,
 \qquad \mathfrak a=(f_1,\ldots,f_N),
\]
as \(\mathfrak a\) ranges over the \(\mathfrak m\)-primary ideals.
\end{theorem}

\begin{proof}
We work first in BFJ's \(\log|f|\)-normalization, as in
\eqref{eq:local-generalized-lelong}.  The preceding factor audit then shows
that the conclusion is unchanged in the book's squared-norm normalization.
If \(n=1\), the normalized centered valuation space consists only of
\(\ord_0\), and the generalized Lelong number of a germ is its ordinary
Lelong number.  We may therefore assume \(n\geq2\).

For an \(\mathfrak m\)-primary ideal
\(\mathfrak a=(f_1,\ldots,f_N)\), set
\[
                         \psi_{\mathfrak a}=\log\max_j|f_j|.
\]
This weight has an isolated pole.  Moreover,
\(\rme^{\psi_{\mathfrak a}}=\max_j|f_j|\) is locally Lipschitz, so
\(\psi_{\mathfrak a}\) is tame by
\cite[Lemma~5.10]{BFJ08}.  Its formal psh function and Cartier
\(b\)-divisor are
\[
 \widehat\psi_{\mathfrak a}=\log|\mathfrak a|,
 \qquad Z(\psi_{\mathfrak a})=Z(\mathfrak a).
\]
Products of primary ideals are primary, and for primary
\(\mathfrak a_i\) and positive integers \(k_i\),
\[
 Z\!\left(\prod_i\mathfrak a_i^{k_i}\right)
                       =\sum_i k_i Z(\mathfrak a_i).
\]

We next isolate the multilinear operation that will be polarized.  If
\(W\) is a nef Weil \(b\)-divisor and
\(C_1,\ldots,C_{n-1}\) are arbitrary exceptional Cartier
\(b\)-divisors, choose a common determination
\(\pi:X_\pi\to(\mathbf C^n,0)\) of the \(C_i\) and put
\[
 I_W(C_1,\ldots,C_{n-1})
 :=W_\pi\cdot(C_1)_\pi\cdots(C_{n-1})_\pi.
\]
This is independent of the common determination.  Indeed, on a higher model
\(\mu:X_{\pi'}\to X_\pi\), one has
\(\mu_*W_{\pi'}=W_\pi\), whereas the Cartier inputs pull back; the
projection formula gives the same number.  Thus \(I_W\) is symmetric and
multilinear in its Cartier arguments.  When those arguments are nef,
\cite[Propositions~4.6 and~4.9]{BFJ08} identify it with the integral of
\(g_W\) against their mixed formal Monge--Amp\`ere measure.  There is only
one possibly non-Cartier Weil input here.

Define
\[
 \Lambda(C_1,\ldots,C_{n-1})
 :=I_{Z(u)}(C_1,\ldots,C_{n-1})
   -I_{Z(v)}(C_1,\ldots,C_{n-1}).
\]
The hypothesis, \eqref{eq:local-generalized-lelong}, and the preceding
intersection formula yield
\[
 \Lambda\bigl(Z(\mathfrak a),\ldots,Z(\mathfrak a)\bigr)=0
\]
for every \(\mathfrak m\)-primary ideal \(\mathfrak a\).
Fix primary ideals \(\mathfrak a_1,\ldots,\mathfrak a_s\), write
\(W_i=Z(\mathfrak a_i)\), and consider the homogeneous polynomial
\[
 P(t_1,\ldots,t_s)
 :=\Lambda\!\left(\sum_i t_iW_i,\ldots,\sum_i t_iW_i\right).
\]
The product-ideal formula shows that
\(P(k_1,\ldots,k_s)=0\) for all positive integers \(k_i\).  Induction on
\(s\), using the one-variable identity principle for polynomials, gives
\(P=0\).  Each polarized coefficient therefore vanishes:
\[
             \Lambda(W_{i_1},\ldots,W_{i_{n-1}})=0.
\]
By \cite[Proposition~1.14]{BFJ08}, the integral exceptional Cartier
\(b\)-divisor group is generated by \(Z(\mathfrak a)\) for primary ideals
\(\mathfrak a\).  Tensoring with \(\mathbf R\) and using multilinearity
gives the exact identity
\[
             \Lambda(C_1,\ldots,C_{n-1})=0
\]
for all real exceptional Cartier \(b\)-divisors \(C_i\).  No density
argument and no polarization of singular currents has been used.

It remains to show that these mixed intersections separate the trace of a
Weil \(b\)-divisor on each model.  Fix a BFJ modification
\(\pi:Y\to(\mathbf C^n,0)\), with exceptional prime divisors
\(E_1,\ldots,E_r\).  Choose an exceptional \(\pi\)-ample real Cartier
divisor \(A\).  Such a divisor exists: a relatively ample Cartier divisor
has, under the identification
\[
             \operatorname{Div}(\pi)\simeq
             N^1(Y/(\mathbf C^n,0))_{\mathbf R}
\]
of \cite[\S1.2]{BFJ08}, a unique exceptional representative with the same
relative numerical class.  On \(\bigoplus_i\mathbf R E_i\), put
\[
                         Q_A(D,C)=D\cdot C\cdot A^{n-2}.
\]
We claim that \(Q_A\) is negative definite.

To verify the claim without hiding the local-to-projective step, compactify
\(\pi\) and resolve the boundary without changing a neighborhood of
\(\pi^{-1}(0)\).  This gives a projective birational morphism
\(\bar\pi:\bar Y\to\bar X\) of smooth projective \(n\)-folds.  Choose a
\(\bar\pi\)-ample Cartier divisor \(\bar A\), and take for \(A\) the
exceptional representative of its relative numerical class near the fibre
over \(0\).  If \(B\) is a sufficiently positive ample divisor on
\(\bar X\), then
\[
                         H=\bar A+\bar\pi^*B
\]
is ample on \(\bar Y\).  For exceptional divisors \(D,C\), every
intersection cycle under consideration is supported over \(0\).  The
classes \(\bar A\) and \(A\) have the same degree on its fibre curves, and
\(\bar\pi^*B\) has degree zero there.  Expanding the powers term by term
therefore gives
\[
       D\cdot C\cdot H^{n-2}
        =D\cdot C\cdot\bar A^{n-2}
        =D\cdot C\cdot A^{n-2}.
\]

The Hodge index theorem on \(\bar Y\) says that
\[
 q_H(L_1,L_2)=L_1\cdot L_2\cdot H^{n-2}
\]
has signature \((1,\rho(\bar Y)-1)\) on
\(N^1(\bar Y)_{\mathbf R}\).  The class \(\bar\pi^*B\) is nef and big,
so \((\bar\pi^*B)^2\cdot H^{n-2}>0\).  Every exceptional divisor is
\(q_H\)-orthogonal to \(\bar\pi^*B\), since the latter restricts
numerically trivially to the fibre over \(0\).  Hence \(q_H\) is negative
definite on the numerical span of the \(E_i\).  This span has no
coefficient-space kernel: a globally numerically trivial exceptional divisor
is relatively numerically trivial, and the displayed BFJ identification then
makes it zero.  The preceding equality of intersection forms proves the
claim.  For the blowup of a smooth point the sign is visible from
\(A=-E\) and \(E^2(-E)^{n-2}=-1\).

Now write
\[
              D=Z(u)_\pi-Z(v)_\pi=\sum_{i=1}^r d_iE_i.
\]
For every exceptional real divisor \(C\) on \(Y\), the mixed vanishing
already proved, applied to the Cartier \(b\)-divisors determined by
\(C,A\), gives
\[
 0=\Lambda(C,A,\ldots,A)
   =D\cdot C\cdot A^{n-2}=Q_A(D,C).
\]
Nondegeneracy implies \(D=0\).  Since \(\pi\) was arbitrary, \(Z(u)\)
and \(Z(v)\) have equal traces on every modification.  Their formal psh
functions consequently agree on each quasimonomial simplex by continuity,
and \cite[Definition~2.8 and Proposition~3.11]{BFJ08} extends them to all
normalized valuations by the same formula
\[
                          g(\nu)=\inf_\pi g(r_\pi(\nu)).
\]
Thus \(\widehat u=\widehat v\) pointwise on \(\mathcal V_0\).  The proof
used only the primary-ideal weights displayed in the statement.
\end{proof}

The theorem concerns valuative transforms, not the original germs: bounded
perturbations are invisible on both sides.  It supplies the converse left
open after Theorem~A of \cite{BFJ08} by a finite-model intersection
argument, without requiring formal Monge--Amp\`ere fundamental solutions.
The frozen pre-integration statement and proof passed production Rethlas run
\nolinkurl{20260824T140127Z_ef8316d19c12} with final verdict
\texttt{correct}, zero critical errors, and zero gaps; an independent hostile
audit found no defect.  This verifies the mathematical proof object, while a
separate post-integration byte-identity gate for the typographical live copy
is not claimed.  This book-level derivation is not accompanied by a priority
claim.

\section{A verified Gauss-shell realization of the local measure}
\label{sec:local-cld-gauss-shell}

The BFJ construction above is birational and trivially valued. To compare
it with a CLD Bedford--Taylor product, one must first place both measures on
one analytic object. The comparison in this section does so on an exact
model locus. It is the result covered by the Rethlas receipt
\texttt{local\_singular\_ma\_comparison.json}; no statement beyond the
displayed ideal, algebraizable rational-model, and joint bounded domains is
being inferred.

Assume now \(n\geq2\) and put \(d=n-1\). Fix
\[
 K=\mathbf C((t)),
 \qquad |t|=\rme^{-1},
\]
and choose regular parameters \(z_1,\ldots,z_n\) at the germ. Let \(B\) be
a closed good, possibly non-strict, \(K\)-analytic polydisc of polyradius
\(\rho\), where \(|t|<\rho<1\), chosen so that the finitely many analytic
germs under consideration converge on \(B\). For an
\(\mathfrak m\)-primary ideal \(\mathfrak a=(f_1,\ldots,f_N)\), set
\[
 u_{\mathfrak a}:=\log\max_j|f_j|,
 \qquad
 u_{\mathfrak m}:=\log\max_i|z_i|,
 \qquad
 \chi_{\mathfrak m}:=\max\{u_{\mathfrak m},-1\},
\tag{70.36}\label{eq:local-cld-model-functions}
\]
and define the shell
\[
                S_{\mathfrak m}:=\{u_{\mathfrak m}=-1\}.
\tag{70.37}\label{eq:local-cld-shell}
\]
An \emph{algebraizable log model} means the analytification, after
shrinking the germ, of a projective modification of a Zariski neighborhood
of the origin in \(\mathbf A^n_{\mathbf C}\), with the relevant divisors and
line bundles descending to that model.

For continuous semipositive model potentials \(u_1,\ldots,u_d\), put
\(T=\dprimeddprime u_1\wedge\cdots\wedge\dprimeddprime u_d\) and define
the swept boundary measure by
\[
 \begin{split}
 \MA^{\partial}_{\mathrm{CLD},\mathfrak m}(u_1,\ldots,u_d)
 :={}&T\wedge\dprimeddprime\chi_{\mathfrak m}\\
 &-\one_{\{u_{\mathfrak m}>-1\}}
       \bigl(T\wedge\dprimeddprime u_{\mathfrak m}\bigr),
 \end{split}
\tag{70.38}\label{eq:local-cld-swept-product}
\]
where the second measure is extended by zero. Both terms are continuous
Bedford--Taylor products in the established CLD domain. On
\(\{u_{\mathfrak m}>-1\}\), locality makes the two terms equal; on
\(\{u_{\mathfrak m}<-1\}\), the cutoff is constant and both vanish. Thus
the difference is shell-supported. This argument deliberately does not say
that \(\dprimeddprime\chi_{\mathfrak m}\) itself is shell-supported; that
claim is false in dimension greater than one.

\begin{lemma}\label{lem:local-cld-gauss-core}
There are maps
\[
 \iota:\mathcal V_0\hookrightarrow S_{\mathfrak m},
 \qquad
 p:S_{\mathfrak m}\longrightarrow\mathcal V_0,
 \qquad p\circ\iota=\operatorname{id}_{\mathcal V_0},
\]
such that
\[
                 u_{\mathfrak a}(\iota(v))=-v(\mathfrak a)
\tag{70.39}\label{eq:local-cld-gauss-ideal}
\]
for every \(\mathfrak m\)-primary ideal \(\mathfrak a\). If
\(\pi:Y\to X\) is an algebraizable simultaneous log resolution and
\(\mathfrak m\mathcal O_Y=\mathcal O_Y(-D)\), every normalized face of its
dual complex is carried by \(\iota\) onto a polyhedral skeleton in a formal
model of the shell. Moreover, \(u_{\mathfrak a}\) is a continuous
semipositive model function, belongs to the revised CLD psh sheaf on a
collar of the shell, pulls back to a subharmonic function on every analytic
curve, and restricts through \(\iota\) to the BFJ function
\(g_{\mathfrak a}=-v(\mathfrak a)\). The same assertions hold for rational
semipositive model functions carried by an algebraizable common log model.
\end{lemma}

\begin{proof}
For \(v\in\mathcal V_0\), define its Gauss extension on convergent Laurent
series by
\[
 \widetilde v\!\left(\sum_{j\geq j_0}t^jf_j\right)
       :=\min_j\{j+v(f_j)\}.
\tag{70.40}\label{eq:local-cld-gauss-extension}
\]
The minimum is attained because \(v(f_j)\geq0\) and \(j\to+\infty\).
Multiplicativity follows from the associated graded domain: the initial
form is a nonzero Laurent polynomial in a degree-one indeterminate over
\(\operatorname{gr}_vR\), so initial forms of nonzero factors cannot
multiply to zero. Since \(v(z_i)\geq1\), the seminorm satisfies
\(|z_i|\leq|t|<\rho\), is bounded on the affinoid algebra of \(B\), and
extends by continuity. Its restriction to \(R\) is \(v\), and
\(\min_i\widetilde v(z_i)=\widetilde v(t)=1\), so it defines
\(\iota(v)\in S_{\mathfrak m}\). Conversely, restriction of a shell point
to \(R\) gives a normalized valuation; call this map \(p\). Formula
\eqref{eq:local-cld-gauss-extension} gives \(p\circ\iota=\operatorname{id}\),
and taking the minimum over ideal generators proves
\eqref{eq:local-cld-gauss-ideal}.

If \(D=\sum_{j=1}^rb_jE_j\) near a stratum and \(x_j=0\) are adapted
equations, a monomial valuation with weights \(w_j\) lies on the normalized
face precisely when
\[
                         \sum_{j=1}^rb_jw_j=1.
\]
Complete \(Y\times\Spec K^\circ\) along its special fibre and take the
normalized admissible blowup of
\((t,\mathcal O_Y(-D))\). Locally the center is \((t,s_D)\), and near the
stratum
\(s_D=\varepsilon x_1^{b_1}\cdots x_r^{b_r}\). The exceptional skeleton at
\(v(t)=1\) is therefore the displayed normalized face. Toroidal blowups of
strata give common subdivisions, and
\eqref{eq:local-cld-gauss-extension} identifies their points with the Gauss
extensions.

It remains to justify algebraizability and positivity. Choose \(N\) with
\(\mathfrak m^N\subset\mathfrak m\mathfrak a\), and replace each analytic
generator by its Taylor polynomial modulo \(\mathfrak m^N\). If
\(\mathfrak b\) is generated by those polynomials, then
\[
 \mathfrak b\subset\mathfrak a
 \subset\mathfrak b+\mathfrak m^N
 \subset\mathfrak b+\mathfrak m\mathfrak a.
\]
Nakayama's lemma applied to \(\mathfrak a/\mathfrak b\) gives
\(\mathfrak a=\mathfrak b\). A finite collection is treated with one
choice of \(N\). Transition coefficients between finite generating sets
are power bounded on \(B\); the ultrametric inequality in both directions
shows that their maximum norms, and hence their analytic weights, agree.

On a simultaneous principalization, write
\(\mathfrak a\mathcal O_Y=\mathcal O_Y(-A)\) and locally
\(f_j=s_Ah_j\). The \(h_j\) generate the unit ideal. On every formal affine
chart, a bounded relation \(1=\sum c_jh_j\) gives
\(\max_j|h_j|=1\), hence
\[
                       u_{\mathfrak a}=\log|s_A|.
\]
Thus the weight is exactly the model weight of the relatively generated,
therefore relatively nef, bundle \(\mathcal O_Y(-A)\). It is one of the
semipositive model generators of the revised CLD psh sheaf. On a curve,
the outgoing-slope sum is the nonnegative degree of the pulled-back model
bundle, so the pullback is subharmonic. Finally,
\eqref{eq:local-cld-gauss-ideal} identifies its Gauss restriction with the
BFJ model function. Tensor powers, multilinearity, and a common
algebraizable model prove the rational assertion.
\end{proof}

\begin{theorem}\label{thm:local-cld-model-comparison}
Let \(\mathfrak a_1,\ldots,\mathfrak a_d\) be
\(\mathfrak m\)-primary ideals. On one simultaneous log resolution write
\[
 \mathfrak m\mathcal O_Y=\mathcal O_Y\!\left(-\sum_Eb_EE\right),
 \qquad
 \mathfrak a_i\mathcal O_Y=\mathcal O_Y(-A_i).
\]
Then
\[
 \MA^{\partial}_{\mathrm{CLD},\mathfrak m}
       (u_{\mathfrak a_1},\ldots,u_{\mathfrak a_d})
 =\sum_E b_E\bigl((-A_1)\cdots(-A_d)\cdot E\bigr)
       \delta_{\iota(v_E)}.
\tag{70.41}\label{eq:local-cld-atomic}
\]
After pushforward by \(p\), this is exactly the BFJ local measure
\[
 \MA_{0}(g_{\mathfrak a_1},\ldots,g_{\mathfrak a_d})
 =\sum_E b_E\bigl((-A_1)\cdots(-A_d)\cdot E\bigr)
       \delta_{v_E}.
\tag{70.42}\label{eq:local-cld-bfj-atomic}
\]
The same comparison holds for rational semipositive model functions carried
by an algebraizable common log model. The common total mass is
\[
 (-A_1)\cdots(-A_d)\cdot D,
 \qquad D=\sum_Eb_EE,
\tag{70.43}\label{eq:local-cld-total-mass}
\]
the mixed Hilbert--Samuel multiplicity with reference ideal
\(\mathfrak m\).
\end{theorem}

\begin{proof}
By \Cref{lem:local-cld-gauss-core}, all ideal data live on an algebraic
simultaneous log resolution \(\pi:Y\to U\), isomorphic away from the origin.
Compactify \(Y\) and the finitely many divisors without changing a
neighborhood of the proper exceptional fibre. Complete
\(Y\times_{\mathbf C}\Spec K^\circ\) \(t\)-adically and form
\[
 q:\mathscr Y:=
 \operatorname{Nor}\!\left(
   \operatorname{Bl}_{(t,\mathcal O_Y(-D))}(Y\times\Spec K^\circ)
 \right)\longrightarrow Y.
\tag{70.44}\label{eq:local-cld-formal-blowup}
\]
After a toroidal resolution, this is a flat formal model of pure relative
dimension \(n\). Every modification is admissible because its center
contains \(t\), so the generic analytic fibre is unchanged. Components
relevant to the swept difference lie over the original proper exceptional
fibre; hence the later compactification contributes no coefficient.

Let \(\mathscr F\) be defined by
\[
 (t,\mathcal O_Y(-D))\mathcal O_{\mathscr Y}
      =\mathcal O_{\mathscr Y}(-\mathscr F).
\]
The model line bundles for the \(n=d+1\) weights in the first term of
\eqref{eq:local-cld-swept-product} are
\[
 \mathcal L_i=q^*\mathcal O_Y(-A_i)\quad(1\leq i\leq d),
 \qquad
 \mathcal L_\chi=\mathcal O_{\mathscr Y}(-\mathscr F).
\tag{70.45}\label{eq:local-cld-model-bundles}
\]
Indeed, if \(\tau\) generates the pulled-back blowup ideal and
\(t=\tau a\), \(s_D=\tau b\), then \((a,b)\) generates the unit ideal.
Thus \(\max\{|a|,|b|\}=1\) on the formal domain and
\[
                \log\max\{|t|,|s_D|\}=\log|\tau|.
\]
The analogous factorization \(f_{ij}=s_{A_i}h_{ij}\), with
\(\max_j|h_{ij}|=1\), identifies each of the first \(d\) bundles. Hence
\eqref{eq:local-cld-model-bundles} describes the metrics themselves, not
only their Cartier divisors.

For model metrics, the CLD formal intersection formula
\cite[Theorem~18.5.5 and Proposition~18.5.8]{CLD25} is
\[
 c_1(\mathcal M_1)\wedge\cdots\wedge c_1(\mathcal M_n)
 =\sum_Gm_G\deg_G(\mathcal M_1\cdots\mathcal M_n)\,\delta_{\xi_G},
\tag{70.46}\label{eq:local-cld-formal-measure}
\]
where \(G\) ranges over special-fibre components and
\(m_G=\ord_G(t)\). The analytic point is normalized by
\(\ord_G/\ord_G(t)\), which explains this factor.

Fix a prime \(E\subset D\). At its generic point, choose a parameter \(x\)
with \(s_D=x^{b_E}\) up to a unit. The normalized blowup of
\((t,x^{b_E})\) has one equality-locus component \(\mathscr E_E\) with
\[
             \ord_{\mathscr E_E}(t)=b_E,
 \qquad
             \ord_{\mathscr E_E}(x)=1.
\tag{70.47}\label{eq:local-cld-boundary-orders}
\]
To see this directly, over \(\mathbf C(E)\) the Rees algebra is
\[
 \mathbf C(E)[x,t][T_0,T_1]/(x^{b_E}T_0-tT_1).
\]
Its exceptional fibre is \(\mathbf P^1_{\mathbf C(E)}\). The chart
\(t=x^{b_E}u\) exhibits the primitive exceptional valuation in
\eqref{eq:local-cld-boundary-orders}; the other chart is normal in
codimension one, so normalization does not change the generic exceptional
curve.

Although \(\mathscr F\) has coefficient \(b_E\) on \(\mathscr E_E\), the
tautological line bundle \(\mathcal L_\chi\) has degree one on the generic
\(\mathbf P^1\)-fibre. The remaining bundles pull back from \(Y\). If
\(p_E:\mathscr E_E\to E\) is the projection, then
\[
 \mathcal L_i|_{\mathscr E_E}
    =p_E^*(\mathcal O_Y(-A_i)|_E),
\]
and the projective-bundle formula gives
\[
 \deg_{\mathscr E_E}
       (\mathcal L_1\cdots\mathcal L_d\mathcal L_\chi)
       =(-A_1)\cdots(-A_d)\cdot E.
\tag{70.48}\label{eq:local-cld-boundary-degree}
\]
Combining \eqref{eq:local-cld-formal-measure},
\eqref{eq:local-cld-boundary-orders}, and
\eqref{eq:local-cld-boundary-degree} yields the \(E\)-atom of
\eqref{eq:local-cld-atomic} at \(\iota(v_E)\).

We now account for every other component. On
\(\{u_{\mathfrak m}>-1\}\), the section \(s_D\) dominates \(t\), so the
model bundle of \(\chi_{\mathfrak m}\) agrees with that of
\(u_{\mathfrak m}\); the second term of
\eqref{eq:local-cld-swept-product} cancels the coefficient. On
\(\{u_{\mathfrak m}<-1\}\), the section \(t\) dominates and
\(\mathcal L_\chi\) is trivial, so the coefficient vanishes. If an
equality-locus component \(G\) maps to a center \(Z\subset Y\) of
codimension at least two, put
\(\alpha_G=c_1(\mathcal L_\chi|_G)\cap[G]\). Since
\(\dim Z\leq n-2\), proper pushforward gives \(q_*\alpha_G=0\) in
\(A_{n-1}(Y)\). The projection formula with the remaining \(d=n-1\)
pulled-back bundles then gives zero degree on \(G\). Thus precisely the
components \(\mathscr E_E\) survive, proving
\eqref{eq:local-cld-atomic}. Relative nefness makes every coefficient
nonnegative.

On the BFJ side, \(g_{\mathfrak a_i}\) is the function of the nef Cartier
b-divisor \(-A_i\). The local atomic formula
\eqref{eq:local-atomic-ma} assigns to \(v_E\) exactly the coefficient in
\eqref{eq:local-cld-bfj-atomic}. Since
\(p(\iota(v_E))=v_E\), the measures agree term by term. Multilinearity
proves the rational common-model case. Finally, summing the atoms gives
\((-A_1)\cdots(-A_d)\cdot D\). Taking general sections of the ideals and
using the projection formula identifies this intersection number with the
mixed Hilbert--Samuel multiplicity. This proves
\eqref{eq:local-cld-total-mass} and completes the proof.
\end{proof}

\begin{proposition}\label{prop:local-cld-independence}
The common model measure of
\Cref{thm:local-cld-model-comparison} is independent of the chosen log
resolution, formal model, and tropical chart. With
\((\mathfrak m,-1)\) fixed, it is local on \(\mathcal V_0\), and literal
restriction to an open subset is restriction of the measure without
renormalizing its mass.
\end{proposition}

\begin{proof}
Two resolutions and formal models admit a common domination
\(q:Y'\to Y\). If \(E'\) is the strict transform of \(E\), the projection
formula gives
\[
 ((-q^*A_1)\cdots(-q^*A_d)\cdot E')
 =((-A_1)\cdots(-A_d)\cdot E).
\]
For a new \(q\)-exceptional prime \(F\), one has \(q_*F=0\) as a divisor
cycle, so its intersection with all the pulled-back factors is zero.
Refinement therefore preserves each old atom and adds no new mass. The
normalization \(\ord_E/\ord_E(\mathfrak m)\) is unchanged as well.

Two tropical charts have a common polyhedral refinement. On it, the weights
are the same intrinsic model functions and
\eqref{eq:local-cld-model-bundles} is unchanged. Finally, both terms of
\eqref{eq:local-cld-swept-product} are local Bedford--Taylor measures and
their difference is shell-supported. Model potentials from the germ factor
through \(p\) on the shell. Hence agreement near an open subset of
\(\mathcal V_0\) gives agreement near its inverse image, and pushforward by
\(p\) proves locality. Ordinary restriction of a measure has this property;
dividing by its new local mass would not.
\end{proof}

Define \(\mathcal C_b^{\mathrm{com}}\) to be the set of bounded continuous
BFJ formal psh functions \(g\) for which there are decreasing semipositive
model functions \(g^{(k)}\downarrow g\), all in the preceding
ideal/algebraizable model locus, whose Gauss extensions \(u^{(k)}\) live on
one collar of the shell and satisfy:
\begin{enumerate}
\item \(u^{(k)}\) belongs to the revised CLD psh sheaf;
\item \(u^{(k)}\) decreases locally boundedly to a continuous revised-CLD
      psh function; and
\item every pullback of \(u^{(k)}\) to an analytic curve is subharmonic.
\end{enumerate}
For finitely many inputs, the collar and algebraizable refinements are
common and synchronized. For a further decreasing sequence, ``joint'' means
that one decreasing two-index array controls both limiting operations.

\begin{proposition}\label{prop:local-cld-bounded-closure}
For \(g_1,\ldots,g_d\in\mathcal C_b^{\mathrm{com}}\), the weak limit
\[
 \MA^{\mathrm{com}}_{\mathfrak m}(g_1,\ldots,g_d)
 :=\lim_{k\to\infty}
 p_*\MA^{\partial}_{\mathrm{CLD},\mathfrak m}
       (u_1^{(k)},\ldots,u_d^{(k)})
\tag{70.49}\label{eq:local-cld-bounded-limit}
\]
exists, is independent of the joint approximants, and equals the BFJ mixed
local intersection measure. It is continuous along decreasing sequences
admitting one joint bounded approximation of the stated kind.
\end{proposition}

\begin{proof}
At every model index, \Cref{thm:local-cld-model-comparison} gives
\[
 p_*\MA^{\partial}_{\mathrm{CLD},\mathfrak m}
       (u_1^{(k)},\ldots,u_d^{(k)})
 =\MA_0(g_1^{(k)},\ldots,g_d^{(k)}).
\tag{70.50}\label{eq:local-cld-model-index}
\]
Let \(B\) denote the symmetric \(d\)-linear BFJ measure operator. Its
polarization identity is
\[
 d!B(h_1,\ldots,h_d)
 =\sum_{J\subset\{1,\ldots,d\}}(-1)^{d-|J|}
 B\!\left(\sum_{j\in J}h_j,\ldots,\sum_{j\in J}h_j\right).
\tag{70.51}\label{eq:local-cld-polarization}
\]
For each \(J\), the diagonal input on the right decreases to the
corresponding sum of limiting functions. BFJ bounded monotone continuity
\cite[Proposition~4.8]{BFJ08}, followed by the finite sum in
\eqref{eq:local-cld-polarization}, proves weak convergence of the mixed
measures. Together with \eqref{eq:local-cld-model-index}, this proves
existence of \eqref{eq:local-cld-bounded-limit}. Notice that no
characteristic function is commuted with a weak limit: the argument passes
through the already swept model measures.

The limiting BFJ product depends only on \(g_1,\ldots,g_d\), proving
independence. For a further decreasing sequence, take the diagonal of the
joint two-index array. It is cofinal and decreases to the final limit, so
the same polarization argument gives the asserted continuity. The curve
submean inequalities pass to the continuous decreasing limit. This proves
the proposition only on \(\mathcal C_b^{\mathrm{com}}\), not for an
arbitrary bounded curve-test psh function.
\end{proof}

If \(I\) is any \(\mathfrak m\)-primary ideal, the same construction uses
the link \(\mathcal V_I=\{v:v(I)=1\}\), the shell \(u_I=-1\), and the
cutoff \(\max\{u_I,-1\}\). Principalizing \(I\) with the input ideals
replaces the coefficient \(b_E\) in
\eqref{eq:local-cld-atomic} by \(\ord_E(I)\). Denote the resulting common
measure on the verified domain by \(\MA_I^{\mathrm{com}}\).

\begin{proposition}\label{prop:local-cld-finite-map}
Let \(F:(X',0')\to(X,0)\) be a finite map of smooth \(n\)-dimensional germs
of local degree \(q\), and put \(I=\mathfrak m\mathcal O_{X',0'}\). Normalize
the source first by \(I\), and define
\(F_\#:\mathcal V_I\to\mathcal V_0\) by \(w\mapsto w\circ F^\#\). For ideal
model functions, for rational functions whose target models and pulled-back
source models lie in the algebraizable locus, and for members of
\(\mathcal C_b^{\mathrm{com}}\) whose pullbacks remain in that class,
\[
 (F_\#)_*\MA_I^{\mathrm{com}}(F^*g_1,\ldots,F^*g_d)
       =q\,\MA_{\mathfrak m}^{\mathrm{com}}(g_1,\ldots,g_d).
\tag{70.52}\label{eq:local-cld-finite-map}
\]
On the standard source link \(\mathcal V_{\mathfrak m'}\), set
\[
 c(v):=v(I),
 \qquad R_I(v):=v/c(v),
 \qquad \widehat F(v):=(v\circ F^\#)/v(I).
\]
For ideal model functions, the exact standard-link formula is
\[
 \widehat F_*\!\left(
 c\,\MA_{\mathfrak m'}^{\mathrm{com}}
   (g_{F^*\mathfrak a_1},\ldots,g_{F^*\mathfrak a_d})\right)
 =q\,\MA_{\mathfrak m}^{\mathrm{com}}
   (g_{\mathfrak a_1},\ldots,g_{\mathfrak a_d}).
\tag{70.53}\label{eq:local-cld-standard-link}
\]
For an \emph{\'etale} map of pointed complex germs, \(q=1\),
\(I=\mathfrak m'\), and \(c\equiv1\).
\end{proposition}

\begin{proof}
Choose compatible analytic log resolutions, and let \(h:Y'\to Y\) be the
normalization of the analytic fibre product before resolving the source.
It is finite of degree \(q\). Divisors introduced only by the final
resolution map to codimension at least two and have zero coefficient by the
projection formula. Let \(E'\) dominate a target prime \(E\), with
ramification and residue degrees \(e(E'/E)\) and \(f(E'/E)\). Then
\[
 \ord_{E'}(I)=e(E'/E)b_E,
 \qquad
 h_*E'=f(E'/E)E.
\tag{70.54}\label{eq:local-cld-ef}
\]
If \(\alpha_E=(-A_1)\cdots(-A_d)\cdot E\), the projection formula gives
\[
 ((-h^*A_1)\cdots(-h^*A_d)\cdot E')=f(E'/E)\alpha_E.
\]
The source atom over \(E'\) therefore has coefficient
\(e(E'/E)f(E'/E)b_E\alpha_E\) and maps to \(v_E\). Summing over
\(E'\mid E\) and using
\(\sum_{E'\mid E}e(E'/E)f(E'/E)=q\) yields
\(qb_E\alpha_E\), which is exactly the target coefficient in
\eqref{eq:local-cld-finite-map}. Multilinearity gives the algebraizable
rational-model case. Joint bounded approximation, bounded monotone
continuity, and continuity of pushforward of finite measures give the
stated bounded extension.

To pass to the standard source link, let
\(v_{E'}=\ord_{E'}/b'_{E'}\) and put
\(c_{E'}=\ord_{E'}(I)\). Then
\[
 c(v_{E'})=c_{E'}/b'_{E'},
 \qquad
 R_I(v_{E'})=\ord_{E'}/c_{E'}.
\]
If \(\beta_{E'}=(-A'_1)\cdots(-A'_d)\cdot E'\), the atom has coefficient
\(b'_{E'}\beta_{E'}\) with reference \(\mathfrak m'\) and
\(c_{E'}\beta_{E'}\) with reference \(I\). Multiplication by
\(c(v_{E'})\) changes the first coefficient into the second, while \(R_I\)
changes only its location. Hence
\[
 \MA_I^{\mathrm{com}}
 =R_{I*}\!\left(c\,\MA_{\mathfrak m'}^{\mathrm{com}}\right)
\]
for the pulled-back ideal inputs. Combining this identity with
\eqref{eq:local-cld-finite-map} proves
\eqref{eq:local-cld-standard-link}. If \(F\) is \'etale
at the marked point, the completed regular local rings are isomorphic, so
all three simplifying assertions follow.
\end{proof}

\paragraph{Worked reference-change audit.}
For
\(F(z_1,\ldots,z_n)=(z_1^q,z_2,\ldots,z_n)\), one has
\[
 v(\mathfrak m\mathcal O_{X',0'})
 =\min\{qv(z_1),v(z_2),\ldots,v(z_n)\}.
\]
This function is not constant on \(\mathcal V_{\mathfrak m'}\). Therefore
the factor \(c\) in \eqref{eq:local-cld-standard-link} cannot be replaced by
the global degree \(q\); the degree appears only after pushforward and the
sum of ramification-residue degrees.

\paragraph{Why no probability normalization is local.}
On \((\mathbf C^2,0)\), blow up the origin, and then blow up the point of
the first exceptional curve lying in the tangent direction \(y=0\), that is,
where the strict transform of \(\{y=0\}\) meets it.  The choice matters: in
adapted coordinates \(x=cd\), \(y=cd^2\) one has \(\ord_{E_2}(x)=1\) and
\(\ord_{E_2}(y)=2\), whereas blowing up a general point of the first
exceptional curve gives \(\ord_{E_2}(x)=\ord_{E_2}(y)=1\) and hence a
different coefficient on \(E_2\) below.  Write \(E_1\) for the strict
transform and \(E_2\) for the new exceptional curve. Then
\[
 E_1^2=-2,
 \qquad E_2^2=-1,
 \qquad E_1\cdot E_2=1,
 \qquad b_{E_1}=b_{E_2}=1.
\]
For \(\mathfrak a=(x^3,xy,y^2)\), the ideal divisor is
\(A=2E_1+3E_2\), and
\[
                      (-A)\cdot E_1=(-A)\cdot E_2=1.
\]
Thus the common measure is
\(\delta_{v_{E_1}}+\delta_{v_{E_2}}\). On an open set containing both
atoms, probability normalization assigns mass \(1/2\) to each; restricting
to a smaller open containing only \(v_{E_1}\) keeps mass \(1/2\), whereas
renormalizing there gives mass one. The unnormalized measure is local, and
the probability-normalized assignment is not a sheaf under literal
restriction.

\paragraph{The unbounded polar endpoint.}
On \((\mathbf C^2,0)\), let \(v_s(x)=1\), \(v_s(y)=s\) for
\(1\leq s<+\infty\), and let \(v_\infty(y)=+\infty\). For
\(\mathfrak a_k=(y,x^k)\),
\[
 g_k(v_s)=-\min\{s,k\}\searrow-s=:g(v_s),
 \qquad
 \MA_0(g_k)=\delta_{v_k}.
\tag{70.55}\label{eq:local-cld-polar-endpoint}
\]
The full decreasing limit of the measures is \(\delta_{v_\infty}\). But
\(g_k=\max\{g,-k\}\), and the nonpluripolar truncation
\[
 \one_{\{g>-k\}}\MA_0(\max\{g,-k\})
\]
vanishes because its only atom lies on \(\{g=-k\}\). Hence full monotone
continuation charges the polar endpoint while nonpluripolar continuation
discards it. The verified model comparison does not choose between these
two extensions.

\section{Polarized affine cones and finite energy}
\label{sec:local-polarized-affine-cones}

There is a second local theory whose carrier, positivity datum, and measure
all differ from the BFJ construction above.  It is the theory of polarized
affine cones initiated in \cite{Wu22AffineEnergy} and extended to finite
energy in \cite{Wu24AffineCones}.  It should be read as a local
Boucksom--Jonsson branch with a Reeb polarization, not as a completion of the
BFJ theory for arbitrary germs.

Throughout this section, $k$ is an algebraically closed trivially valued
field of characteristic zero.  Let
\[
 X=\Spec R,
 \qquad
 T=(\mathbf G_{\mathrm m,k})^r,
 \qquad
 R=\bigoplus_{\alpha\in M}R_\alpha
\]
be a normal affine $T$-variety of dimension $n$.  The action is assumed
effective and \emph{good}: it has a unique fixed point $o$, and $o$ lies
in the closure of every orbit.  Write $\mathfrak m$ for the maximal ideal
of $o$, let $M=\operatorname{Hom}(T,\mathbf G_{\mathrm m})$, and put
$N=\operatorname{Hom}(\mathbf G_{\mathrm m},T)$.  If
$\Lambda=\{\alpha:R_\alpha\ne0\}$, the Reeb cone is
\[
 \mathfrak t^+_R
 :=\left\{\xi\in N_{\mathbf R}\ \middle|\
       \langle\xi,\alpha\rangle>0
       \text{ for every }\alpha\in\Lambda\setminus\{0\}\right\}.
\tag{70.56}\label{eq:local-cone-reeb-cone}
\]
A \emph{polarized affine cone} is the datum $(X,T;\xi)$ with
$\xi\in\mathfrak t^+_R$.  The adjective ``polarized'' is literal: changing
$\xi$ changes the psh and energy classes even though it does not change
the underlying analytic link.

Write a point of $X^{\An}$ additively as a semivaluation
$v:R\to\mathbf R\cup\{+\infty\}$, so that
$\log|f|(v)=-v(f)$.  The torus Gauss map associates to
$\eta\in N_{\mathbf R}$ the semivaluation
\[
 (\eta*v)\!\left(\sum_\alpha f_\alpha\right)
  =\min_\alpha\bigl\{v(f_\alpha)
                  +\langle\eta,\alpha\rangle\bigr\}.
\tag{70.57}\label{eq:local-cone-reeb-action}
\]
At $\eta=0$, this map is the torification
$\rho_T(v)(\sum f_\alpha)=\min_\alpha v(f_\alpha)$, rather than the
identity at an arbitrary analytic point.  It is a continuous retraction
onto the $T$-invariant locus: invariance means exactly
$v(\sum f_\alpha)=\min_\alpha v(f_\alpha)$
\cite[Section~2.2]{Wu24AffineCones}.  Continuity follows because each
evaluation is a finite minimum of continuous evaluations.  Formula
\eqref{eq:local-cone-reeb-action} is an $N_{\mathbf R}$-action on this
invariant locus.  The radial direction relevant to the polarization there
is $t\mapsto t\xi*v$.
A function $\varphi$ is $\xi$-equivariant when
$\varphi(t\xi*v)=\varphi(v)-t$ wherever both sides are finite.

\begin{definition}
\label{def:local-cone-link}
The non-Archimedean link of the affine cone is
\[
                         X_0:=\{v\in X^{\An}:v(\mathfrak m)=0\}.
\tag{70.58}\label{eq:local-cone-link}
\]
\end{definition}

This is a compact analytic domain, independent of $\xi$
\cite[Introduction and Definition~2.6]{Wu24AffineCones}.  Its normalization must not be
confused with $\mathcal V_0=\{v:v(\mathfrak m)=1\}$ in
\eqref{eq:local-normalized-space}: the latter cuts the positive-scaling cone
of centered valuations, while, after $\xi$ has been fixed,
$X_0^T:=X_0\cap(X^{\An})^T$ is a cross-section for Reeb-translation
orbits in the invariant punctured cone.  Indeed, the reference function
below changes by $-t$, so each such orbit meets its zero set exactly once.
The full link $X_0$ also contains non-invariant points.  Since homogeneous
generators compute $v(\mathfrak m)$, torification preserves $X_0$;
its image $X_0^T$ is a compact closed retract of that link.

For a nonzero homogeneous element $f\in R_\alpha$, set
$w_\xi(f)=\langle\xi,\alpha\rangle$.  The reference function is
\[
 \varphi_\xi(v)
 :=\sup_{\substack{0\ne f\in R_\alpha\\ \alpha\ne0}}
             \frac{\log|f|(v)}{w_\xi(f)}.
\tag{70.59}\label{eq:local-cone-reference}
\]
The supremum is a finite maximum after choosing finitely many homogeneous
algebra generators of $R$ in positive weight
\cite[Lemma~5.3]{Wu22AffineEnergy}.

\begin{lemma}
\label{lem:local-cone-reference-link}
One has $X_0=\{\varphi_\xi=0\}$.  In particular, the zero set is
independent of the Reeb field although the function $\varphi_\xi$ is not.
\end{lemma}

\begin{proof}
Choose homogeneous elements $g_1,\ldots,g_N$ of positive weights which
generate $R$ as a $k$-algebra.  They generate $\mathfrak m$ as an
ideal.  Every homogeneous $f$ is a linear combination of monomials in the
$g_j$ of the same $\xi$-weight.  The ultrametric inequality and the
weighted-average inequality therefore give
\[
 \frac{-v(f)}{w_\xi(f)}
 \leq \max_j\frac{-v(g_j)}{w_\xi(g_j)}.
\]
Consequently the right-hand side is exactly the supremum in
\eqref{eq:local-cone-reference}.  Now
$v(\mathfrak m)=\min_jv(g_j)$.  This minimum is zero exactly when every
$v(g_j)$ is nonnegative and at least one is zero, which is equivalent to
$\max_j\{-v(g_j)/w_\xi(g_j)\}=0$.
\end{proof}

\subsection{Fubini--Study, continuous psh, and general psh functions}
\label{subsec:local-cone-psh-classes}

Let $f_j\in R_{\alpha_j}$ be homogeneous elements whose common zero set is
$\{o\}$, and let $\lambda_j\in\mathbf R$.  A Fubini--Study function is
of the form
\[
 \varphi
 =\max_{1\leq j\leq N}
       \frac{\log|f_j|+\lambda_j}
            {\langle\xi,\alpha_j\rangle}.
\tag{70.60}\label{eq:local-cone-fs}
\]
Write $\mathcal H(\xi)$ for this class.  Each member is continuous on
$X^{\An}\setminus\{o\}$, is $\xi$-equivariant, and is determined by its
restriction to $X_0^T$.  In fact $\varphi=\varphi\circ\rho_T$, since
the presentation uses only homogeneous elements.  Constants and finite maxima preserve
$\mathcal H(\xi)$.  Addition necessarily changes the equivariance weight,
so it does not preserve a fixed $\mathcal H(\xi)$.

Let $C^0(X^{\An}\setminus\{o\},\xi)$ denote the space of continuous
$\xi$-equivariant functions with the topology of uniform convergence on
compact subsets.  Following \cite[Definition~3.1]{Wu24AffineCones}, put
\[
 \mathrm{CPSH}(\xi)
 :=\overline{\mathcal H(\xi)}^{\,C^0_{\mathrm{loc}}}.
\tag{70.61}\label{eq:local-cone-cpsh}
\]
The rational approximation of the Reeb field is part of this closure
mechanism.  It yields stability under constants, maxima, positive rescaling,
and convex combinations.  The exact proportional-polarization calculus is
\[
 a\varphi\in\mathrm{CPSH}(a^{-1}\xi),
 \qquad
 \left.
 \begin{array}{c}
  \varphi\in\mathrm{CPSH}(a\xi),\\
  \psi\in\mathrm{CPSH}(b\xi)
 \end{array}
 \right\}
 \Longrightarrow
 \varphi+\psi\in
 \mathrm{CPSH}\!\left(\frac{ab}{a+b}\xi\right)
\]
for $a,b>0$ \cite[Proposition~3.3]{Wu24AffineCones}.  This statement
concerns proportional rescalings of one Reeb ray; it gives no comparison
between nonproportional Reeb fields.  The density argument requires the
invariant carrier, as follows.

\begin{proposition}\label{prop:local-cone-invariant-density}
Differences of continuous psh functions with a common proportional
polarization form a vector lattice whose restrictions are uniformly dense
in $C^0(X_0^T)$.  Every such difference factors through $\rho_T$.
\end{proposition}

\begin{proof}
The factorization holds for Fubini--Study functions and passes to their
locally uniform limits.  Proportional rescaling and addition show that the
differences form a real vector space containing constants.  If
$u=\varphi_1-\varphi_2$ and $v=\psi_1-\psi_2$, with the two pairs
polarized by $a\xi$ and $b\xi$, respectively, then
\[
 \max\{u,v\}
 =\max\{\varphi_1+\psi_2,\psi_1+\varphi_2\}
       -(\varphi_2+\psi_2).
\]
Both terms on the right have polarization $ab\xi/(a+b)$.  Negation gives
closure under minima.

Distinct invariant semivaluations differ on a nonzero homogeneous element
$f$ of positive weight $w=\langle\xi,\alpha\rangle$.  For $C>0$,
the Fubini--Study function
$\max\{\log|f|/w,\varphi_\xi-C\}$ restricts on $X_0$ to
$\max\{-v(f)/w,-C\}$.  Choose $C$ large enough to distinguish the
two values, including when one value is infinite.  Thus the lattice
separates points of the compact Hausdorff space $X_0^T$, and lattice
Stone--Weierstrass proves density.  This gives the invariant-carrier version
of the argument in \cite[Theorem~3.5]{Wu24AffineCones}.
\end{proof}

\begin{example}[Why the full link needs torification]
\label{ex:local-cone-full-link-density-failure}
Take $X=\mathbf A_k^1$, $T=\mathbf G_{\mathrm m}$, and $\xi=1$.
The trivial valuation and the rigid point $z=1$ both lie in $X_0$,
since they give $z$ value zero.  They agree on every homogeneous element
$cz^d$, so every Fubini--Study function, continuous psh function, and
their differences has equal values at the two points.  They are distinct:
$v_{\mathrm{triv}}(z-1)=0$, whereas $v_{z=1}(z-1)=+\infty$.
The continuous function $|z-1|$ on $X_0$ separates them.  Consequently
these differences are not dense in $C^0(X_0)$.  The full-link wording of
\cite[Theorem~3.5]{Wu24AffineCones} needs this carrier restriction.
\end{example}

\begin{definition}
\label{def:local-cone-psh}
A function $\varphi:X^{\An}\setminus\{o\}\to
\mathbf R\cup\{-\infty\}$ belongs to $\PSH(\xi)$ if it is not identically
$-\infty$ and is the pointwise limit of a decreasing net in
$\mathrm{CPSH}(\xi)$.
\end{definition}

Nets, rather than only sequences, are natural because $X_0$ need not be
metrizable.  The definition implies upper semicontinuity.  It also gives
decreasing approximation directly in the model class
$\mathcal H(\xi)$, and a psh function is determined by its values on the
quasimonomial locus \cite[Proposition~5.2 and Theorem~5.8]{Wu24AffineCones}.
Thus divisorial and quasimonomial tests enter the topology, but irrational
Reeb fields force one to retain genuinely quasimonomial points.

\begin{example}[The standard affine cone]
\label{ex:local-cone-affine-space}
Take $R=k[z_1,\ldots,z_n]$, let $T=\mathbf G_{\mathrm m}$ act with
weight one on every coordinate, and take $\xi=1$.  Then
\[
 X_0=\left\{v:\min_i v(z_i)=0\right\},
 \qquad
 \varphi_\xi(v)=\max_i\{-v(z_i)\}.
\]
The homogeneous function
\[
 \max_j\frac{-v(f_j)+\lambda_j}{\deg f_j}
\]
is a member of $\mathcal H(1)$ whenever the homogeneous $f_j$ generate
an $\mathfrak m$-primary ideal.  For a rational Reeb field this cone is
the affine cone over $\mathbf P^{n-1}$; the appearance of $n-1$ curvature
inputs below is therefore the expected projective dimension.
\end{example}

\begin{example}[An irrational Reeb field]
\label{ex:local-cone-irrational-reeb}
Let $X=\mathbf A_k^2$, let $T=\mathbf G_{\mathrm m}^2$ act diagonally,
and take $\xi=(1,\sqrt2)$.  The link remains
$\{\min(v(x),v(y))=0\}$, whereas
\[
 \varphi_\xi(v)=\max\{-v(x),-v(y)/\sqrt2\}.
\]
The incommensurable denominators explain why a sum of finite-max
presentations need not remain a finite-max presentation for one integral
grading.  Passing first to $\mathrm{CPSH}(\xi)$ restores convexity, but it
does not turn the irrational polarization into a quasi-regular one.  In
particular, the associated Monge--Amp\`ere measure may be supported at a
genuinely quasimonomial rather than divisorial point
\cite[Introduction]{Wu24AffineCones}.
\end{example}

\subsection{The shell factor, energy, and finite-energy measures}
\label{subsec:local-cone-shell-energy}

Put $\varphi_\xi^+=\max\{\varphi_\xi,0\}$.  For
$\varphi_1,\ldots,\varphi_{n-1}\in\mathcal H(\xi)$, Wu's model measure is
the CLD product
\[
 \MA_\xi(\varphi_1,\ldots,\varphi_{n-1})
 :=\frac{1}{\vol(\xi)}
   \dprimeddprime\varphi_1\wedge\cdots\wedge
   \dprimeddprime\varphi_{n-1}\wedge
   \dprimeddprime\varphi_\xi^+.
\tag{70.62}\label{eq:local-cone-model-ma}
\]
Here $\vol(\xi)>0$ is the Reeb volume.  The last factor is not decorative:
it sweeps the radial current onto the boundary
$\{\varphi_\xi=0\}=X_0$.  With the normalization in
\eqref{eq:local-cone-model-ma}, the result is a probability measure supported
on $X_0^T$ \cite[Corollary~5.7]{Wu22AffineEnergy}
\cite[Theorem~2.12 and Proposition~4.2(1)]{Wu24AffineCones}.  Formula
\eqref{eq:local-cone-model-ma} is therefore a second, intrinsic occurrence of
the same ``one reference factor plus $n-1$ tangential factors'' pattern as
the Gauss-shell construction in \eqref{eq:local-cld-swept-product}.  The two
formulas have the same differential architecture, but they live on different
links and no equality between them is implied.

For $\varphi\in\mathrm{CPSH}(\xi)$, the energy relative to
$\varphi_\xi$ is
\[
 E_\xi(\varphi)
 :=\frac{1}{n\vol(\xi)}
   \sum_{j=0}^{n-1}
   \int_{X^{\An}}
    (\varphi-\varphi_\xi)
    (\dprimeddprime\varphi)^j
    \wedge(\dprimeddprime\varphi_\xi)^{n-1-j}
    \wedge\dprimeddprime\varphi_\xi^+.
\tag{70.63}\label{eq:local-cone-continuous-energy}
\]
Integration by parts and the local Hodge-index inequality give monotonicity,
concavity, and the first-variation formula on this continuous domain
\cite[Propositions~3.8--3.9 and~4.2]{Wu24AffineCones}.  For a general psh
function the energy is the monotone extension
\[
\begin{aligned}
 E_\xi(\varphi)
 &:=\inf\left\{E_\xi(\psi)\ \middle|\
       \psi\in\mathrm{CPSH}(\xi),\ \varphi\leq\psi\right\},\\
 \mathcal E^1(\xi)
 &:=\left\{\varphi\in\PSH(\xi):E_\xi(\varphi)>-\infty\right\}.
\end{aligned}
\tag{70.64}\label{eq:local-cone-finite-energy}
\]
It is increasing, concave, upper semicontinuous, continuous along decreasing
nets, and satisfies $E_\xi(\varphi+c)=E_\xi(\varphi)+c$
\cite[Proposition~5.10]{Wu24AffineCones}.

Let $\mathcal P(X_0)$ be the Radon probability measures on the link.  The
dual energy and its finite-energy domain are
\[
\begin{aligned}
 E_\xi^\vee(\mu)
 &:=\sup_{\varphi\in\mathcal E^1(\xi)}
       \left\{E_\xi(\varphi)
       -\int_{X_0}\varphi\,\rd\mu\right\},\\
 \mathcal M^1(\xi)
 &:=\left\{\mu\in\mathcal P(X_0):E_\xi^\vee(\mu)<+\infty\right\}.
\end{aligned}
\tag{70.65}\label{eq:local-cone-dual-energy}
\]
For $\varphi,\psi\in\mathcal E^1(\xi)$, set
\[
 I(\varphi,\psi)
 :=\int_{X_0}(\varphi-\psi)
             \bigl(\MA_\xi(\psi)-\MA_\xi(\varphi)\bigr).
\tag{70.66}\label{eq:local-cone-I}
\]
On the measure side, put
\[
 J_\mu(\varphi)
 :=E_\xi^\vee(\mu)-E_\xi(\varphi)
        +\int_{X_0}\varphi\,\rd\mu,
 \qquad
 I^\vee(\mu,\nu)
 :=\inf_{\varphi\in\mathcal E^1(\xi)}
       \bigl(J_\mu(\varphi)+J_\nu(\varphi)\bigr).
\tag{70.67}\label{eq:local-cone-dual-I}
\]
Both $I$ and $I^\vee$ are quasi-pseudometrics; constants are invisible to
$I$, as they must be because they do not change the Monge--Amp\`ere
measure.

\refstepcounter{theorem}\label{thm:local-cone-finite-energy-ma}
\begin{quote}
\textbf{Foundational input \thetheorem\ (source boundary).}
For a polarized affine cone $(X,T;\xi)$ as above, the mixed operator in
\eqref{eq:local-cone-model-ma} has a unique extension to tuples in
$\mathcal E^1(\xi)$.  Its values are Radon probability measures supported
on $X_0^T$, and the mixed pairing is continuous when the potentials and the
test function decrease in $\mathcal E^1(\xi)$.  In particular,
\[
 \MA_\xi:\bigl(\mathcal E^1(\xi),I\bigr)
       \longrightarrow\bigl(\mathcal M^1(\xi),I^\vee\bigr)
\]
extends the operator on $\mathcal H(\xi)$, is continuous along decreasing
nets, and is bi-Lipschitz onto its image, with dimensional constants for the
two quasi-pseudometrics.  The statement is imported from
\cite[Theorem~1.1, Theorem~6.3, and Corollary~7.8]{Wu24AffineCones}; no part
of the finite-energy estimates is proved here.  We use the invariant-carrier
version of the construction: \Cref{prop:local-cone-invariant-density}
supplies the separating test lattice, and the model measures have invariant
support.  That closed support persists under weak limits, so the extension
is viewed as a measure on $X_0$ by inclusion.
\end{quote}

The words ``onto its image'' are essential.  The cited results do not
characterize the image of $\MA_\xi$ inside $\mathcal M^1(\xi)$, and
they do not solve the local non-Archimedean Monge--Amp\`ere equation for an
arbitrary finite-energy measure with invariant support.  The missing input identified in the
source is a uniform differentiability theorem strong enough for the
variational argument.

There is also a carrier constraint.  All psh functions in this section
factor through $\rho_T$, by decreasing approximation.  Hence
$E_\xi^\vee(\mu)=E_\xi^\vee((\rho_T)_*\mu)$, and
$J_\mu$ and $I^\vee$ depend only on the torified measures.  The
integral identity follows pointwise from
$\varphi=\varphi\circ\rho_T$, including extended values.  The density
proposition recovers measures supported on $X_0^T$ from these tests; it
does not distinguish arbitrary measures within a torification fiber.
In particular, a measure outside that invariant carrier cannot be a value
of the canonical operator, even if its torification has finite energy.

The proof in the source has three layers worth keeping distinct.  On
$\mathrm{CPSH}(\xi)$, CLD approximation, integration by parts, and the
Hodge-index inequality produce the energy and the $I$--$J$ estimates.
Mixed products are then extended one finite-energy input at a time, using
uniform energy bounds to obtain decreasing-net continuity.  Finally, the
Legendre transform \eqref{eq:local-cone-dual-energy} and the dual
$I^\vee$-estimates yield the last assertion.  This architecture explains
why a bounded BFJ measure alone does not automatically supply the theorem:
the shell factor, the Reeb polarization, and the proportional-rescaling
estimates all enter the construction.

Two normalization checks are immediate from the formulas.  Replacing
$\varphi$ by $\varphi+c$ leaves
\eqref{eq:local-cone-model-ma} unchanged and adds $c$ to the energy; hence
the expression in \eqref{eq:local-cone-dual-energy} is constant-invariant.
Also, the top degree in \eqref{eq:local-cone-model-ma} is $n$: there are
$n-1$ psh inputs and exactly one radial reference factor.  Omitting the
last factor produces neither a probability measure on $X_0$ nor the
energy in \eqref{eq:local-cone-continuous-energy}.

\subsection{The exact BJ and CLD interfaces}
\label{subsec:local-cone-interfaces}

The construction has two proved interfaces.

\paragraph{The CLD interface.}
The Fubini--Study functions in \eqref{eq:local-cone-fs} are
psh-approachable in the CLD sense, so
\eqref{eq:local-cone-model-ma} is an actual CLD Bedford--Taylor product, not
only an intersection-theoretic analogy.  Local uniform closure gives the
continuous products used in \eqref{eq:local-cone-continuous-energy}, and the
source proves integration by parts on differences of continuous psh
functions.  This is a concrete common domain between local differential
theory and a finite-energy variational theory.  It does not identify every
function in the transfinite CLD psh hierarchy with $\PSH(\xi)$, nor does it
remove the Reeb equivariance or the shell factor.

\paragraph{The quasi-regular BJ interface.}
If $\xi$ is rational, then after scaling it to integral weights the cone is
the affine cone over a polarized orbifold pair $(V,B_V;L)$, with
$\dim V=n-1$ \cite[Proposition~2.4]{Wu24AffineCones}.  There is an embedding
$\iota:V^{\An}\hookrightarrow X_0^{\langle\xi\rangle}$
\cite[Lemma~5.11]{Wu22AffineEnergy}.  Under this
embedding, the $T$-homogeneous Fubini--Study functions correspond to
$T$-invariant Fubini--Study metrics on $L$
\cite[Lemma~5.12]{Wu22AffineEnergy}.  The action on the projective quotient
factors through $T/\langle\xi\rangle$.  If $\phi_j$ corresponds to
$\varphi_j$, then, with the integral normalization for which
$\vol(L)=\vol(\xi)$ \cite[Lemma~3.16]{Wu22AffineEnergy},
\[
 \iota_*\!\left(
  \frac{1}{\vol(L)}
  \dprimeddprime\phi_1\wedge\cdots\wedge
  \dprimeddprime\phi_{n-1}
 \right)
 =\MA_\xi(\varphi_1,\ldots,\varphi_{n-1}).
\]
This is the normalized form of
\cite[Proposition~5.13]{Wu22AffineEnergy}.  The Fubini--Study energies then
agree, and decreasing completion recovers the BJ finite-energy class and
energy on this invariant branch
\cite[Corollary~5.14]{Wu22AffineEnergy}
\cite[Corollary~6.6(4)]{Wu24AffineCones}.  These are the precise recovered
objects.  When $T=\langle\xi\rangle$, the residual action is trivial,
and this restriction permits all projective Fubini--Study metrics.  For a
larger fixed torus, rationality of $\xi$ does not remove invariance.

For an irrational Reeb field there is no projective quotient with one
integral grading to which this argument can be applied verbatim.  Separately,
for any two nonproportional Reeb fields---whether rational or
irrational---there is no established comparison or addition calculus between
their psh classes.  Rational approximation transports many
fixed-polarization statements but does not fill that gap.  Together with the
unresolved image of $\MA_\xi$, these are two distinct boundaries of the
current affine-cone theory.

\begin{remark}[Two local links]
\label{rem:local-two-links}
The BFJ and affine-cone theories should not be joined by identifying their
links.  The BFJ link $v(\mathfrak m)=1$ consists of centered valuations and
normalizes positive scalar multiplication.  The affine-cone link
$v(\mathfrak m)=0$ is an analytic domain containing the trivial valuation
and its invariant retract normalizes Reeb translation.  The former is defined for a smooth local
germ without a torus; the latter requires a normal affine cone, a good torus
action, and a Reeb field.  Their model functions and top-degree measures can
be compared only after additional data produce a common carrier.  No such
general comparison theorem is asserted here.
\end{remark}

\section{Comparison with the global and local analytic approaches}
\label{sec:local-comparison-ledger}

The same words---model, psh, envelope, and Monge--Amp\`ere---refer to
different signatures in the theories surveyed in this book.  The following
ledger records the exact interfaces.

\begin{description}[style=nextline,leftmargin=1.5em,labelindent=0pt,
                    itemsep=.75\baselineskip]
\item[Underlying space]
\emph{Local BFJ theory.} Normalized link $\nu(\mathfrak m)=1$.

\emph{Polarized affine-cone theory.} Analytic link
$v(\mathfrak m)=0$, together with Reeb-equivariant functions on the
punctured cone.

\emph{Global BFJ/BJ theory.} Full analytification or homogeneous cone.

\emph{CLD/GR local theory.} Analytic open subsets and tropical charts.

\item[Positivity datum]
\emph{Local BFJ theory.} Relative nefness over one closed point.

\emph{Polarized affine-cone theory.} A fixed Reeb field and local-uniform
closure of homogeneous Fubini--Study functions.

\emph{Global BFJ/BJ theory.} Semipositive metric in a fixed class.

\emph{CLD/GR local theory.} Hessian, current, or curve-test positivity.

\item[Model core]
\emph{Local BFJ theory.} The unit ideal, $\mathfrak m$-primary ideals,
and exceptional Cartier b-divisors.

\emph{Polarized affine-cone theory.} Homogeneous primary ideals, finite-type
norms, and Reeb-weighted Fubini--Study functions.

\emph{Global BFJ/BJ theory.} Models, vertical divisors, FS metrics.

\emph{CLD/GR local theory.} Smooth psh, semipositive PL, local FS.

\item[Approximation]
\emph{Local BFJ theory.} Decreasing unit-or-primary ideal functions; away
from the zero formal psh function, all sufficiently late ideals are
$\mathfrak m$-primary.

\emph{Polarized affine-cone theory.} Local-uniform Fubini--Study
approximation for continuous psh functions and decreasing nets for general
psh functions.

\emph{Global BFJ/BJ theory.} Decreasing model/FS regularization in stated
global domains.

\emph{CLD/GR local theory.} CLD decreasing closure or GR local PL/LFS
approximation.

\item[Measure]
\emph{Local BFJ theory.} $n-1$ inputs on the normalized link.

\emph{Polarized affine-cone theory.} $n-1$ tangential inputs and one Reeb
shell factor on an $n$-dimensional cone.

\emph{Global BFJ/BJ theory.} $n$ curvature inputs on an $n$-fold.

\emph{CLD/GR local theory.} Continuous Bedford--Taylor products in the
established CLD domain.

\item[Boundary data]
\emph{Local BFJ theory.} The ideal $\mathfrak m$ and relative exceptional
tower.

For an excellent local germ,
\cref{sec:local-excellent-boundary-links} supplies the intrinsic ray space,
compact primary sections, exact changes of scale, and conductor quotient.
Nonprimary ideals give partial sections; reduction preserves the datum.
For a one-dimensional reduced excellent germ,
\Cref{thm:p121-reduced-nagata-curve-link} identifies the ideal-normalized
link with the finite normalization branches and gives the exact rescaling
under a change of primary ideal; \Cref{prop:p126-finite-reduced-curve-link}
gives its finite-map branch formula.  Moreover,
\Cref{prop:p126-curve-conductor-principal-realization} realizes every
sufficiently deep integral branch-order vector by a downstairs primary
principal ideal, while \Cref{thm:p122-curve-principal-ideal-envelope}
computes the closure and raw envelope of the resulting principal-ideal
profiles.  The book-local completion
\Cref{thm:p122-curve-ideal-psh-functoriality} realizes that closure as a
normalization-independent homogeneous branch-cone object, completes it under
decreasing limits, and proves exact envelope and finite-map pullback formulas.
For branch-surjective finite curve maps,
\Cref{thm:p126-curve-generic-norm-trace-descent} adds the generic-degree
weighted trace, normalized pullback retraction, and exact
conductor-controlled descent of principal profiles, together with the sharp
pushforward--envelope inequality, norm-one barycentric projection, and dual
total-variation-isometric signed-measure lift with its adjoint norm-one
projection.  The same barycentric operator is the finite conditional
expectation for each lifted full-support base probability, with the exact
module law, all-$L^p$ contraction, $L^2$ orthogonality, and variance
identity.  For finite flat maps,
\Cref{thm:p126-curve-finite-flat-norm-descent} adds the determinant-norm
specialization, including a canonical unscaled downstairs norm and its
tower law.

\emph{Polarized affine-cone theory.} The Reeb field, the reference function
$\varphi_\xi$, and the shell current
$\dprimeddprime\varphi_\xi^+$.

\emph{Global BFJ/BJ theory.} Polarization and total mass.

\emph{CLD/GR local theory.} Analytic boundary, compact support, or Shilov
boundary.
\end{description}

There are six proved overlap statements.

\begin{enumerate}
\item Ideal functions have the same valuation formula in every approach:
      $\log|\mathfrak a|=-\nu(\mathfrak a)$.
\item On the model locus, both the local measure and the global
      Chambert--Loir measure are determined by intersection numbers, but on
      different spaces and with different numbers of inputs.
\item The local normalized link embeds into the centered cone of the
      trivially valued analytification; positive rescaling recovers exactly
      the locus $0<\nu(\mathfrak m)<+\infty$, not the semivaluation
      supported at the closed point.
\item The Gauss section of \Cref{lem:local-cld-gauss-core} realizes the
      normalized link as a Gauss core inside the shell $u_{\mathfrak m}=-1$ over
      $\mathbf C((t))$. On the ideal and algebraizable rational-model
      locus, the swept CLD boundary measure agrees term by term with the BFJ
      measure; \Cref{prop:local-cld-bounded-closure} extends this equality
      only to the explicitly defined joint bounded closure.
\item On a polarized affine cone, homogeneous Fubini--Study functions are
      CLD-approachable, and their affine-cone Monge--Amp\`ere measure is the
      literal CLD product \eqref{eq:local-cone-model-ma}.  The extension to
      $\mathcal E^1(\xi)$ is the imported result
      \Cref{thm:local-cone-finite-energy-ma}, not a formal consequence of
      the continuous product.
\item For a rational Reeb field, the projective quotient identifies the
      affine-cone Fubini--Study functions, mixed measure, energy, and
      finite-energy class with their $T$-invariant global BJ counterparts on the
      quasi-regular branch, as recorded after
      \eqref{eq:local-cone-dual-I}.
\end{enumerate}

The fourth statement is a genuine common analytic realization, but it does
not identify arbitrary formal psh functions with arbitrary CLD-v2 or
classically psh functions. The BFJ local theory begins trivially valued and
birational; the realization adds the specific nontrivially valued Gauss
shell and the swept reference factor. The strongest GR comparison theorems
in \Cref{chap:verified-extensions} still require their own good, strict
analytic domains. Dropping algebraizability, the joint approximation, or
the boundary datum would enlarge the theorem and is not licensed by the
comparison above.  The fifth and sixth statements concern a different
carrier with a Reeb polarization.  They supply continuous and finite-energy
theory for polarized affine cones, but they do not extend the fourth
statement to arbitrary BFJ formal psh functions or to arbitrary analytic
germs.

\section{What a full local Boucksom--Jonsson theory would require}
\label{sec:local-general-theory-requirements}

The BFJ construction supplies the first three rows below, and the exact
common model core supplies the fourth.  Wu's polarized affine-cone theory
supplies the fifth row under additional symmetry and polarization
hypotheses.  The remaining rows describe, rather than assume, what is still
missing for a general local theory.

\begin{enumerate}
\item A compact normalized valuative link and a dense divisorial/model locus.
      This is \Cref{thm:local-valuation-compactness}.
\item A psh cone defined by relative nef b-divisors, stable under maxima,
      with a decreasing unit-or-primary ideal regularization for each
      element, and an eventually $\mathfrak m$-primary regularization for
      each nonzero element.  This is
      \Cref{thm:local-nef-psh-equivalence}.
\item Mixed positive measures for bounded formal psh data, with ideal
      approximation and local intersection formulas.  This is
      \Cref{thm:local-formal-ma}.
\item A common analytic realization on ideal and algebraizable rational
      model data, together with the joint bounded monotone closure and the
      finite-map/reference-change formulas. These are
      \Cref{thm:local-cld-model-comparison,prop:local-cld-bounded-closure,prop:local-cld-finite-map}.
\item On a normal affine cone with a good torus action and a fixed Reeb
      field, a continuous psh cone, decreasing-net completion, finite-energy
      class, mixed Monge--Amp\`ere operator, dual finite-energy measures, and
      strong quasi-pseudometrics.  These are the source-bounded results of
      \Cref{sec:local-polarized-affine-cones}.  They recover the invariant BJ
      branch for a rational Reeb field, but the comparison of nonproportional Reeb
      polarizations remains unavailable.
\item For excellent local germs, the normalizing boundary and its change
      with the ideal are now supplied by
      \cref{sec:local-excellent-boundary-links}, including conductor
      descent in arbitrary dimension and the loss of compactness for
      nonprimary sections. For reduced excellent curve germs, \Cref{thm:p121-reduced-nagata-curve-link}
      settles the underlying branch set and primary-ideal scaling, while
      \Cref{prop:p126-finite-reduced-curve-link} settles the normalized
      branch map under finite local morphisms.
      \Cref{prop:p126-curve-conductor-principal-realization} settles
      conductor descent for sufficiently deep integral branch orders, and
      \Cref{thm:p122-curve-principal-ideal-envelope} computes the closure and
      raw envelope of this explicit ideal-function family.
      \Cref{thm:p122-curve-ideal-psh-functoriality} gives its book-local
      ideal-psh completion, homogeneous normalization transport, and
      finite-map-compatible envelope.  The generic trace theorem
      \Cref{thm:p126-curve-generic-norm-trace-descent} gives a pushforward
      retraction, its exact envelope inequality, total-variation-isometric
      signed-measure lift and dual norm-one projection, the associated
      finite conditional-expectation and variance identities, and
      conductor-scaled principal descent exactly when the branch map is
      surjective, while
      \Cref{thm:p126-curve-finite-flat-norm-descent} replaces the conductor
      choice by one canonical determinant norm for finite flat curve maps.
      Identifying that candidate with an
      independently defined singular psh or nef cone, and every
      intersection-theoretic layer, remain beyond these calculations.  In
      fact,
      \Cref{thm:p122-conductorial-disc-radial-psh-classification} identifies
      the radial-extendable cone on every isolated conductor-controlled
      disc-normalized curve singularity as
      $\{0\}\cup[-\infty,0)^q$, excluding precisely the mixed-zero faces
      of the complete ideal-psh orthant; the ordinary multinode is the
      explicit special case of
      \Cref{thm:p122-multinode-radial-psh-classification}.
\item For a general analytic germ without the affine-cone datum, a singular
      finite-energy class would still need a relative energy, capacity,
      monotone nonpluripolar products, and a comparison principle.  These do
      not follow from the bounded BFJ intersection measure, and the
      affine-cone construction cannot be transported after forgetting its
      Reeb field.
\item The unrestricted local-envelope assertion is false by
      \Cref{prop:ch15-nodal-local-envelope}.  Any positive theorem must impose
      a branch-compatibility or unibranch hypothesis and must still prove
      that usc regularization preserves the chosen psh class.
\item Gluing local measures requires boundary or reference data.  Literal
      restriction cannot turn probability measures into a sheaf, because a
      connected open set can contain two disjoint nonempty opens.
\end{enumerate}

This list explains the correct relationship with the global BJ theory.  The
local object is a normalized link with relative intersection theory, not a
global probability theory with the polarization deleted.

\section{Research boundary and source questions}
\label{sec:local-research-boundary}

The generalized-Lelong converse left open in \cite{BFJ08} is now proved in
\Cref{thm:local-generalized-lelong-separation}.  The two questions below
about varying Reeb fields and the image of the affine-cone Monge--Amp\`ere
operator are explicitly isolated in the introduction of
\cite{Wu24AffineCones}.  The other items are comparison problems exposed by
the frameworks in this book.  The list is neither asserted to be exhaustive
nor used as a source of theorems; no answer is claimed here without a
separate dated literature audit and, for a new argument, a Rethlas receipt.

\begin{enumerate}
\item Does every classical psh germ on a suitable non-Archimedean analytic
      space admit decreasing local FS or semipositive-PL regularization in
      higher dimension?  The known curve and boundaryless PL results do not
      imply this for arbitrary singular functions.  The affirmative
      decreasing-net statement on polarized affine cones requires Reeb
      equivariance and therefore does not answer this arbitrary-local
      question.
\item Can the BFJ measure in \Cref{thm:local-formal-ma} be identified with a
      CLD Bedford--Taylor product beyond the ideal/algebraizable model locus
      and the joint bounded closure of
      \Cref{prop:local-cld-bounded-closure}?
\item Can one define a finite-energy completion on the normalized BFJ link of
      a general germ, extending the verified model-domain functoriality of
      \Cref{prop:local-cld-finite-map}, and whose energy contains the correct
      boundary and polar-mass terms?  The construction in
      \Cref{sec:local-polarized-affine-cones} answers the analogous question
      only when a good torus action and a Reeb field are part of the datum.
\item With the normalized boundary and conductor quotient of
      \cref{sec:local-excellent-boundary-links} fixed, how do those data
      enter an independently geometric singular psh class and the full
      Monge--Amp\`ere theory beyond the curve ideal-profile constructions of
      \cref{thm:p122-curve-principal-ideal-envelope,thm:p122-curve-ideal-psh-functoriality}?
\item Can the local nef envelope be extended from bounded fixed-skeleton
      obstacles to arbitrary usc obstacles while preserving continuity or an
      appropriate capacity regularity?
\item For two nonproportional Reeb fields $\xi,\xi'$, what is the precise
      relation between $\PSH(\xi)$ and $\PSH(\xi')$?  In particular, is
      there a cross-polarization operation strong enough to prove the uniform
      differentiability estimate that replaces
      \cite[Theorem~8.5]{BJ22} in the affine-cone setting?
\item What is the image of
      $\MA_\xi:\mathcal E^1(\xi)\to\mathcal M^1(\xi)$, and under which
      hypotheses does a prescribed finite-energy measure on $X_0^T$ solve the local
      non-Archimedean Monge--Amp\`ere equation?  For irrational $\xi$, the
      image of a divisorial norm may already be a genuinely quasimonomial
      measure, so a divisorial reduction alone cannot settle the question.
\end{enumerate}

The Rethlas verification protocol in
\Cref{chap:verified-extensions} is applied to any proposed answer before it
is promoted into this chapter.  Failed attempts and conditional reductions
remain in the editorial research ledger.

\section*{Notes and further reading}

The entire local construction, including the normalized valuation space,
nef Weil divisors, multiplier-ideal approximation, local Monge--Amp\`ere
measure, valuative transforms, and generalized Lelong formula, is due to
Boucksom--Favre--Jonsson \cite{BFJ08}.  The dimension-two predecessor is the
valuative tree theory cited there.  Strong openness is due to Guan--Zhou
\cite{GZ15}.  The converse separation statement
\Cref{thm:local-generalized-lelong-separation} is the book-level
primary-ideal polarization and relative Hodge-index argument proved above;
it is not attributed to those sources and carries no priority claim.  The
global analogues and their different hypotheses are
developed in \Cref{chap:bfj-psh,chap:bj-linear-growth,chap:mp-psh}; the CLD
and classical local calculi are treated in
\Cref{chap:cld-forms,chap:cld-currents,chap:verified-extensions}.  The
polarized affine-cone model energy, its equality with filtration volume, and
its quasi-regular projective comparison are due to Wu
\cite{Wu22AffineEnergy}.  The continuous and decreasing-net psh classes,
finite-energy functions and measures, mixed Monge--Amp\`ere extension, and
the $I$--$I^\vee$ comparison are due to
\cite{Wu24AffineCones}.  The latter source is a version-1 preprint dated
19 June 2024; its explicit cross-polarization and image questions are not
presented here as solved.  The
Gauss-shell comparison is the new manuscript argument verified by
\texttt{gpt-5.6-sol} at \texttt{xhigh} scrutiny on 14 August 2026. Its
publication receipt has statement digest
\texttt{923e70de7a01\ldots}, proof digest
\texttt{64776c912e47\ldots}, and verdict \texttt{correct} with zero gaps.
Its primary-source inputs are BFJ Propositions~4.8--4.9 and the CLD formal
measure formula \cite[Theorem~18.5.5 and Proposition~18.5.8]{CLD25}; no
priority claim is made for the comparison itself.
